% L'emphase ; redacción : carta de motivación (120 palabras) — Espagnol Tle A — Cours signé OnBuch+
% Programme officiel MINESEC. Compiler avec : tectonic E18-enfasis-carta-de-motivacion.tex
\newcommand{\DOCMATIERE}{Espagnol}
\newcommand{\DOCNIVEAU}{Tle A}
\newcommand{\DOCMODULE}{Módulo 4 : Activités économiques et monde du travail}
\newcommand{\DOCLECON}{E18}
\newcommand{\DOCTITRE}{L'emphase ; redacción : carta de motivación (120 palabras)}
% OnBuch+ — préambule « cours signés » ENRICHI (compilé avec tectonic / XeLaTeX)
% Système visuel « Pop » d'OnBuch+ : fond crème (jamais blanc), bordures encre
% épaisses, ombres « dures » décalées, titres Archivo Black en capitales.
% Version MATHÉMATIQUES : graphes (pgfplots/tikz), boîtes pédagogiques
% (définition, propriété, méthode, attention, le savais-tu, exercice résolu,
% exercice, TP, bilan), page de garde et signature OnBuch+ ancrée en bas.
%
% Macros attendues dans main.tex AVANT \input{preamble} :
%   \DOCMATIERE  \DOCNIVEAU  \DOCMODULE  \DOCLECON (numéro)  \DOCTITRE
\documentclass[11pt]{article}

\usepackage{fontspec}
\usepackage[spanish,french]{babel} % français = langue principale ; l'espagnol via \esp{..} / espagnol
\usepackage[a4paper,margin=2cm,top=3.4cm,bottom=2.8cm,headheight=1.4cm,footskip=1.3cm]{geometry}
\usepackage{amsmath,amssymb,amsfonts}
\usepackage{mathtools} % \xmapsto, \coloneqq... souvent utilisés par les modèles
\usepackage{cancel} % simplifications barrées (bilans de forces)
% \abs{z} (module, valeur absolue) et \norm{u} (norme), utilisés par les modèles
\providecommand{\abs}{}\let\abs\relax\DeclarePairedDelimiter{\abs}{\lvert}{\rvert}
\providecommand{\norm}{}\let\norm\relax\DeclarePairedDelimiter{\norm}{\lVert}{\rVert}
\usepackage[table]{xcolor} % [table] : fournit aussi \rowcolors (alternance de lignes)
\usepackage{pagecolor}
\usepackage{tikz}
\usetikzlibrary{arrows.meta,positioning,calc,shapes.geometric,shapes.misc,shapes.arrows,decorations.pathmorphing,decorations.pathreplacing,decorations.markings,patterns,shadows,mindmap,trees,fit,backgrounds,matrix,intersections,angles,quotes}
\usepackage{pgfplots}
\usepackage[version=4]{mhchem} % équations nucléaires \ce{^{235}_{92}U}
\usepackage[european,siunitx]{circuitikz} % schémas électriques
\pgfplotsset{compat=1.18}
\usepgfplotslibrary{fillbetween,groupplots} % aires entre courbes (calcul intégral)
\usepackage{enumitem}
\usepackage{fancyhdr}
% [most] charge toutes les bibliothèques tcolorbox (listings, minted, raster,
% documentation...) et gonfle inutilement la mémoire TeX sur un document long
% (~300 boîtes) : on ne charge que ce qui est réellement utilisé.
\usepackage[skins,breakable]{tcolorbox}
\usepackage{graphicx}
\usepackage{colortbl}
\usepackage{booktabs}
\usepackage{tabularx}
\usepackage{diagbox} % case d'en-tête diagonale des tableaux à double entrée
% Colonnes extensibles centrée / gauche / droite (C, L, R), souvent utilisées par les modèles
\newcolumntype{C}{>{\centering\arraybackslash}X}
\newcolumntype{L}{>{\raggedright\arraybackslash}X}
\newcolumntype{R}{>{\raggedleft\arraybackslash}X}
\usepackage{array}
\usepackage{multicol}
\usepackage{siunitx}
% parse-numbers=false : accepte aussi \SI{2,5 \times 10^{-3}}{...}
\sisetup{locale=FR,output-decimal-marker={,},group-minimum-digits=4,parse-numbers=false}
\DeclareSIUnit\ohm{\ensuremath{\symup{\Omega}}} % U+2126 absent de la police
\DeclareSIUnit\division{div} % réglages d'oscilloscope (ms/div, V/div)
\DeclareSIUnit\tr{tr} % tours (vitesse de rotation, ex. \SI{1500}{\tr\per\minute})
\DeclareSIUnit\molar{M} % concentration molaire (ex. \SI{3}{\molar}, \SI{1}{\micro\molar})
\DeclareSIUnit\an{an} % année (le modèle confond parfois avec \year, un compteur TeX) — ex. \SI{5}{\kilo\meter\per\an}
\DeclareSIUnit\FCFA{FCFA} % franc CFA (ex. \SI{500}{\FCFA\per\kilo\gram})
\DeclareSIUnit\degreeBrix{\textdegree Brix} % teneur en sucre d'un sirop/jus (réfractométrie)
\DeclareSIUnit\month{mois} % le modèle utilise parfois \month, un compteur TeX, comme unité de durée
\usepackage{qrcode}
% \degree : commande fréquemment utilisée par les modèles pour un angle ou une
% température en dehors de \SI{}{} (non fournie par les paquets chargés).
\providecommand{\degree}{^{\circ}}
% \qed (fin de démonstration), utilisé par les modèles sans amsthm
\providecommand{\qed}{\hfill$\square$}
% \barre : double barre des valeurs interdites dans un tableau de variations (tkz-tab)
\providecommand{\barre}{\ensuremath{\|}}
% environnement proof (amsthm non chargé) utilisé par les modèles
\ifdefined\proof\else\newenvironment{proof}[1][Démonstration]{\par\noindent\textit{#1.}\ }{\qed\par}\fi
\usepackage{lastpage}
\usepackage{hyperref}
\hypersetup{hidelinks}

% ============================================================
% Palette « Pop » OnBuch+ — mêmes hex que la classe P (plus_ui.dart)
% ============================================================
\definecolor{poporange}{HTML}{F59321}
\definecolor{popdark}{HTML}{E07A0C}
\definecolor{popcream}{HTML}{FFF6E8}
\definecolor{popink}{HTML}{1C1714}
\definecolor{poppurple}{HTML}{7A5AE0}
\definecolor{popgreen}{HTML}{2E9E6A}
\definecolor{poppink}{HTML}{E0457B}
\definecolor{popblue}{HTML}{2F7DE1}
\definecolor{popgold}{HTML}{C98A12}
\definecolor{popmuted}{HTML}{8A7F76}
\definecolor{popline}{HTML}{EADFD2}
\definecolor{popgray}{HTML}{8A7F76}
\colorlet{popgrey}{popgray}
% Alias tolérants (le modèle nomme parfois les couleurs différemment)
\colorlet{popwhite}{white}
\colorlet{popblack}{popink}
\colorlet{popred}{poppink}
\colorlet{popyellow}{popgold}
% Teintes claires (fonds de tableaux/encadrés) — le modèle nomme parfois
% les couleurs avec un suffixe L (ex. popblueL) sans les avoir définies.
\colorlet{poporangeL}{poporange!20}
\colorlet{popdarkL}{popdark!20}
\colorlet{poppurpleL}{poppurple!20}
\colorlet{popgreenL}{popgreen!20}
\colorlet{poppinkL}{poppink!20}
\colorlet{popblueL}{popblue!20}
\colorlet{popgoldL}{popgold!20}
\colorlet{popredL}{popred!20}
\colorlet{popgrayL}{popgray!20}
% Teintes claires pour fonds de tableaux / zones
\colorlet{popblueL}{popblue!10!white}
\colorlet{poporangeL}{poporange!14!white}
\colorlet{popgreenL}{popgreen!12!white}
\colorlet{poppurpleL}{poppurple!12!white}
\colorlet{poppinkL}{poppink!12!white}
\colorlet{popgoldL}{popgold!14!white}

\pagecolor{popcream}

% ============================================================
% Typographie
% ============================================================
\defaultfontfeatures{Ligatures=TeX}
\setmainfont{PlusJakartaSans-Medium.ttf}[Path=../../../fonts/,BoldFont=PlusJakartaSans-Bold.ttf]
% Maths en sans-serif (Fira Math) avec chiffres et lettres droites de Plus
% Jakarta, pour que les formules se fondent dans le texte.
\usepackage[math-style=ISO,bold-style=ISO]{unicode-math}
\setmathfont{FiraMath-Regular.otf}[Path=../../../fonts/]
\setmathfont{PlusJakartaSans-Medium.ttf}[Path=../../../fonts/,range={up/num,up/{Latin,latin},\mathup}]
\usepackage{icomma} % virgule décimale sans espace en mode math
% Caractères mathématiques Unicode absents des polices de texte (Plus Jakarta,
% Archivo Black) : rendus par leur équivalent mathématique, en texte comme en
% formule — sinon ils s'affichent en carré vide (ex. « Arithmétique dans ℤ »).
\usepackage{newunicodechar}
\newunicodechar{ℝ}{\ensuremath{\mathbb{R}}}
\newunicodechar{ℤ}{\ensuremath{\mathbb{Z}}}
\newunicodechar{ℂ}{\ensuremath{\mathbb{C}}}
\newunicodechar{ℕ}{\ensuremath{\mathbb{N}}}
\newunicodechar{ℚ}{\ensuremath{\mathbb{Q}}}
\newunicodechar{𝔻}{\ensuremath{\mathbb{D}}}
\newunicodechar{α}{\ensuremath{\alpha}}
\newunicodechar{β}{\ensuremath{\beta}}
\newunicodechar{γ}{\ensuremath{\gamma}}
\newunicodechar{δ}{\ensuremath{\delta}}
\newunicodechar{ε}{\ensuremath{\varepsilon}}
\newunicodechar{θ}{\ensuremath{\theta}}
\newunicodechar{λ}{\ensuremath{\lambda}}
\newunicodechar{μ}{\ensuremath{\mu}}
\newunicodechar{σ}{\ensuremath{\sigma}}
\newunicodechar{τ}{\ensuremath{\tau}}
\newunicodechar{φ}{\ensuremath{\varphi}}
\newunicodechar{ψ}{\ensuremath{\psi}}
\newunicodechar{ω}{\ensuremath{\omega}}
\newunicodechar{Σ}{\ensuremath{\Sigma}}
% majuscules produites par \MakeUppercase dans les titres (α-aminés -> Α)
\newunicodechar{Α}{\ensuremath{\alpha}}
\newunicodechar{Β}{\ensuremath{\beta}}
\newunicodechar{Γ}{\ensuremath{\Gamma}}
\newunicodechar{Δ}{\ensuremath{\Delta}}
\newunicodechar{Ε}{\ensuremath{\varepsilon}}
\newunicodechar{Θ}{\ensuremath{\Theta}}
\newunicodechar{Λ}{\ensuremath{\Lambda}}
\newunicodechar{Μ}{\ensuremath{\mu}}
\newunicodechar{Τ}{\ensuremath{\tau}}
\newunicodechar{Φ}{\ensuremath{\Phi}}
\newunicodechar{Ψ}{\ensuremath{\Psi}}
\newunicodechar{Ω}{\ensuremath{\Omega}}
\newunicodechar{↦}{\ensuremath{\mapsto}}
\newunicodechar{∈}{\ensuremath{\in}}
\newunicodechar{∉}{\ensuremath{\notin}}
\newunicodechar{∧}{\ensuremath{\wedge}}
\newunicodechar{⇔}{\ensuremath{\Leftrightarrow}}
\newunicodechar{⟺}{\ensuremath{\Longleftrightarrow}}
\newunicodechar{⇒}{\ensuremath{\Rightarrow}}
\newunicodechar{⟹}{\ensuremath{\Longrightarrow}}
\newunicodechar{∘}{\ensuremath{\circ}}
\newunicodechar{∪}{\ensuremath{\cup}}
\newunicodechar{∩}{\ensuremath{\cap}}
\newunicodechar{≡}{\ensuremath{\equiv}}
\newunicodechar{⊂}{\ensuremath{\subset}}
\newunicodechar{⊥}{\ensuremath{\perp}}
\newunicodechar{∃}{\ensuremath{\exists}}
\newunicodechar{∀}{\ensuremath{\forall}}
\newunicodechar{∛}{\ensuremath{\sqrt[3]{\,}}}
\newunicodechar{∜}{\ensuremath{\sqrt[4]{\,}}}
\newunicodechar{⃗}{\ensuremath{{}^{\to}}}
\newunicodechar{ⁿ}{\ifmmode^{n}\else\textsuperscript{\ensuremath{n}}\fi}
\newunicodechar{ˣ}{\ifmmode^{x}\else\textsuperscript{\ensuremath{x}}\fi}
\newunicodechar{ᵇ}{\ifmmode^{b}\else\textsuperscript{\ensuremath{b}}\fi}
\newunicodechar{ʳ}{\ifmmode^{r}\else\textsuperscript{\ensuremath{r}}\fi}
\newunicodechar{ᵘ}{\ifmmode^{u}\else\textsuperscript{\ensuremath{u}}\fi}
\newunicodechar{ʸ}{\ifmmode^{y}\else\textsuperscript{\ensuremath{y}}\fi}
\newunicodechar{ᵃ}{\ifmmode^{a}\else\textsuperscript{\ensuremath{a}}\fi}
\newunicodechar{ᵉ}{\ifmmode^{e}\else\textsuperscript{\ensuremath{e}}\fi}
\newunicodechar{ᵏ}{\ifmmode^{k}\else\textsuperscript{\ensuremath{k}}\fi}
\newunicodechar{ᵖ}{\ifmmode^{p}\else\textsuperscript{\ensuremath{p}}\fi}
\newunicodechar{ᴺ}{\ifmmode^{N}\else\textsuperscript{\ensuremath{N}}\fi}
\newunicodechar{ᶜ}{\ifmmode^{c}\else\textsuperscript{\ensuremath{c}}\fi}
\newunicodechar{ʰ}{\ifmmode^{h}\else\textsuperscript{\ensuremath{h}}\fi}
\newunicodechar{ᵠ}{\ifmmode^{q}\else\textsuperscript{\ensuremath{q}}\fi}
\newunicodechar{ₙ}{\ifmmode_{n}\else\textsubscript{\ensuremath{n}}\fi}
\newunicodechar{ₐ}{\ifmmode_{a}\else\textsubscript{\ensuremath{a}}\fi}
\newunicodechar{ᵢ}{\ifmmode_{i}\else\textsubscript{\ensuremath{i}}\fi}
\newunicodechar{ᵣ}{\ifmmode_{r}\else\textsubscript{\ensuremath{r}}\fi}
\newunicodechar{ₖ}{\ifmmode_{k}\else\textsubscript{\ensuremath{k}}\fi}
\newunicodechar{ᵦ}{\ifmmode_{\beta}\else\textsubscript{\ensuremath{\beta}}\fi}
\newunicodechar{ₓ}{\ifmmode_{x}\else\textsubscript{\ensuremath{x}}\fi}
\newfontfamily\popdisplay{ArchivoBlack-Regular.ttf}[Path=../../../fonts/]
\newfontfamily\poplogo{SpaceGrotesk-Bold.ttf}[Path=../../../fonts/]
\newfontfamily\popsemi{PlusJakartaSans-SemiBold.ttf}[Path=../../../fonts/]
\newfontfamily\popxbold{PlusJakartaSans-ExtraBold.ttf}[Path=../../../fonts/]
\newfontfamily\popxboldls{PlusJakartaSans-ExtraBold.ttf}[Path=../../../fonts/,LetterSpace=3.0]

\setlist[enumerate]{leftmargin=1.7em,itemsep=5pt,topsep=4pt}
\setlist[itemize]{leftmargin=1.7em,itemsep=4pt,topsep=4pt,label={\color{poporange}\textbullet}}
\setlength{\parindent}{0pt}
\setlength{\parskip}{5pt}
\renewcommand{\arraystretch}{1.25}

% pgfplots : style Pop par défaut
\pgfplotsset{
  every axis/.append style={axis line style={popink,line width=1pt},tick style={popink},
    grid style={popline,line width=0.5pt},label style={font=\small},tick label style={font=\footnotesize}},
  popcurve/.style={poporange,line width=1.8pt,no markers,smooth},
}
% Même style disponible dans un \draw TikZ simple (hors pgfplots)
\tikzset{popcurve/.style={poporange,line width=1.8pt}}
\tikzset{popgrid/.style={popline,very thin}}
\colorlet{popcurve}{poporange} % le nom du style est parfois utilisé comme couleur

% ============================================================
% Pill / squiggle
% ============================================================
\newcommand{\poppill}[3]{%
  \tikz[baseline=-0.55ex]\node[draw=popink,line width=1.2pt,rounded corners=2.6pt,%
    fill=#2,inner xsep=6.5pt,inner ysep=3pt]{%
    \popxboldls\fontsize{7}{7}\selectfont\color{#3}\MakeUppercase{#1}};%
}
\newcommand{\popsquiggle}[1]{%
  \tikz[baseline=-0.4ex]{
    \draw[#1,line width=2.6pt,line cap=round]
      plot [smooth] coordinates {(0,0.09) (0.34,0.01) (0.68,0.21) (1.02,0.01) (1.36,0.10)};
  }%
}
% Mot-clé surligné (marqueur orange sous le texte)
\newcommand{\cle}[1]{{\popxbold\color{popdark}#1}}

% ============================================================
% En-tête / pied
% ============================================================
\pagestyle{fancy}
\fancyhf{}
\renewcommand{\headrulewidth}{0pt}
\renewcommand{\footrulewidth}{0pt}
\lhead{\raisebox{-0.15ex}{\poplogo\fontsize{13}{13}\selectfont\color{popink}OnBuch\color{poporange}+}}
\rhead{\poppill{\DOCMATIERE\ \textbullet\ \DOCNIVEAU\ \textbullet\ Leçon \DOCLECON}{white}{popink}}
\lfoot{\popxbold\fontsize{7.6}{7.6}\selectfont\color{popmuted}\MakeUppercase{Cours signé OnBuch+}\ \textbullet\ {\popsemi\DOCTITRE}}
\rfoot{\tikz[baseline=-0.6ex]\node[rounded corners=8.5pt,fill=popink,minimum height=17pt,minimum width=17pt,inner xsep=5pt,inner sep=0]{\popxbold\fontsize{8}{8}\selectfont\color{popcream}\thepage\,/\,\pageref*{LastPage}};}

% ============================================================
% Titres
% ============================================================
\newcommand{\chaptitre}[2]{%
  \begin{tcolorbox}[enhanced,colback=poporange,colframe=popink,boxrule=2.6pt,arc=11pt,
      drop shadow=popink,left=15pt,right=15pt,top=12pt,bottom=11pt,before skip=3pt,after skip=18pt]
    \poppill{#2}{popink}{popcream}\par\vspace{9pt}
    {\raggedright\hyphenpenalty=10000\exhyphenpenalty=10000\popdisplay\fontsize{22}{24}\selectfont\color{popcream}\MakeUppercase{#1}\par}
  \end{tcolorbox}
}
% Section numérotée : pastille orange avec le numéro + titre Archivo
\newcounter{popsec}
\newcommand{\coursec}[1]{%
  \stepcounter{popsec}\setcounter{popsub}{0}%
  \par\vspace{16pt}\noindent
  \begin{minipage}{\linewidth}
  \tikz[baseline=-0.7ex]\node[draw=popink,line width=1.6pt,fill=poporange,rounded corners=4pt,minimum size=22pt,inner sep=0]{\popdisplay\fontsize{12}{12}\selectfont\color{popcream}\thepopsec};%
  \hspace{8pt}{\popdisplay\fontsize{15}{17}\selectfont\color{popink}\MakeUppercase{#1}}\par
  \vspace{4pt}\noindent{\color{poporange}\rule{36pt}{3.6pt}}
  \end{minipage}\par\nopagebreak\vspace{8pt}
}
\newcounter{popsub}
\newcommand{\courssub}[1]{%
  \stepcounter{popsub}%
  \par\vspace{9pt}\noindent{\popxbold\fontsize{11.5}{13}\selectfont\color{popink}{\color{poporange}\thepopsec.\thepopsub}\hspace{6pt}#1}\par\nopagebreak\vspace{4pt}
}

% ============================================================
% Boîtes pédagogiques — même recette : fond blanc, bordure épaisse colorée,
% ombre dure encre, badge plein. Titre optionnel [#1].
% ============================================================
\tcbset{popbase/.style={breakable,enhanced,colback=white,boxrule=2.3pt,arc=9pt,
  shadow={3pt}{-3pt}{0pt}{popink},left=14pt,right=14pt,top=11pt,bottom=11pt,before skip=11pt,after skip=15pt,
  fontupper=\popsemi\fontsize{10.6}{14.5}\selectfont\color{popink},
  fontlower=\popsemi\fontsize{10.6}{14.5}\selectfont\color{popink}}}
% Coche dessinée (le glyphe ✓ n'existe pas dans les polices du document)
\newcommand{\popcheck}[1]{\tikz[baseline=-0.5ex]\draw[#1,line width=1.6pt,line cap=round,line join=round](-0.13,0.0)--(-0.04,-0.09)--(0.13,0.1);}
\providecommand{\coche}{\popcheck{popgreen}}
\providecommand{\resultat}[1]{\par\smallskip\noindent\fcolorbox{popgreen}{popgreenL}{\parbox{\dimexpr\linewidth-2\fboxsep-2\fboxrule\relax}{\textbf{Résultat.} #1}}\par\smallskip}
\newcommand{\popboxhead}[3]{\poppill{#1}{#2}{white}\ifx\relax#3\relax\else\hspace{6pt}{\popxbold\fontsize{10.6}{12}\selectfont\color{popink}#3}\fi\par\vspace{7pt}}

\newtcolorbox{aretenir}[1][]{popbase,colframe=popgreen,before upper={\popboxhead{À retenir}{popgreen}{#1}}}
% Texte en espagnol (document de lecture, dialogue, poème, extrait) : cadre
% d'étude. Un titre optionnel (source, auteur) s'affiche dans l'en-tête.
\newtcolorbox{textebox}[1][]{popbase,colframe=popblue,before upper={\popboxhead{Texto}{popblue}{#1}\begin{otherlanguage}{spanish}},after upper={\end{otherlanguage}}}
% Marques ✓ / ✗ (forme correcte / fautive) souvent utilisées par les modèles
\providecommand{\cross}{\textcolor{poppink}{\ensuremath{\times}}}
\providecommand{\xmark}{\cross}\providecommand{\crossmark}{\cross}\providecommand{\cmark}{\ensuremath{\checkmark}}
% Ligne de réponse / justification à compléter (inventée par les modèles)
\providecommand{\justification}[1]{\par\noindent\textit{Justification~:} \dotfill\par}
% \textipa (package tipa non chargé) : la police ne couvre pas l'API -> simple passage
\providecommand{\textipa}[1]{#1}
% Soulignement (ulem non chargé) : \uline, \ul
\providecommand{\uline}[1]{\underline{#1}}\providecommand{\ul}[1]{\underline{#1}}
% Mot ou phrase en espagnol à l'intérieur d'un paragraphe français
\newcommand{\esp}[1]{\ifmmode\text{\textit{#1}}\else\begingroup\selectlanguage{spanish}\itshape #1\endgroup\fi}
\newtcolorbox{exemplebox}[1][]{popbase,colframe=poporange,before upper={\popboxhead{Exemple}{poporange}{#1}}}
\newtcolorbox{definition}[1][]{popbase,colframe=popblue,before upper={\popboxhead{Définition}{popblue}{#1}}}
\newtcolorbox{propriete}[1][]{popbase,colframe=poppurple,before upper={\popboxhead{Propriété}{poppurple}{#1}}}
\newtcolorbox{methode}[1][]{popbase,colframe=popgold,before upper={\popboxhead{Méthode}{popgold}{#1}}}
\newtcolorbox{attention}[1][]{popbase,colframe=poppink,before upper={\popboxhead{Attention}{poppink}{#1}}}
\newtcolorbox{experience}[1][]{popbase,colframe=popblue,colback=popblueL,before upper={\popboxhead{Expérience}{popblue}{#1}}}
% « Le savais-tu ? » : carte violette pleine (contexte, culture, Cameroun)
\newtcolorbox{savaistu}[1][]{popbase,colback=poppurple,colframe=popink,
  fontupper=\popsemi\fontsize{10.4}{14.2}\selectfont\color{white},
  before upper={\poppill{Le savais-tu\,?}{white}{poppurple}\ifx\relax#1\relax\else\hspace{6pt}{\popxbold\color{white}#1}\fi\par\vspace{7pt}}}
% Exercice résolu : énoncé en haut, \tcblower puis la solution sur fond crème
\newcounter{popexo}\newcounter{popexr}
\newtcolorbox{exoresolu}[1][]{popbase,colframe=poporange,colbacklower=popcream,
  segmentation style={popink,line width=1.2pt,dashed},
  before upper={\stepcounter{popexr}\popboxhead{Exercice résolu \thepopexr}{poporange}{#1}},
  before lower={\poppill{Solution}{popink}{popcream}\par\vspace{6pt}}}
% Exercice (non résolu en place) — la correction est dans la partie Corrigés
\newtcolorbox{exercice}[1][]{popbase,colframe=popink,
  before upper={\stepcounter{popexo}\popboxhead{Exercice \thepopexo}{popink}{#1}}}
% Corrigé
\newtcolorbox{corrige}[1][]{popbase,colframe=popgreen,colback=popgreenL,
  before upper={\popboxhead{Corrigé}{popgreen}{#1}}}
% Objectifs / prérequis / bilan
\newtcolorbox{objectifs}{popbase,colframe=popink,colback=poporangeL,
  before upper={\popboxhead{Objectifs de la leçon}{popink}{}}}
\newtcolorbox{prerequis}{popbase,colframe=popink,colback=popblueL,
  before upper={\popboxhead{Prérequis}{popblue}{}}}
\newtcolorbox{bilan}{popbase,colframe=popink,colback=white,boxrule=2.8pt,
  before upper={\popboxhead{Fiche bilan}{poporange}{L'essentiel à connaître pour l'examen}}}
% Figure encadrée avec légende
\newtcolorbox{popfigure}[1][]{enhanced,breakable=false,colback=white,colframe=popline,boxrule=1.4pt,arc=8pt,
  left=10pt,right=10pt,top=10pt,bottom=8pt,before skip=10pt,after skip=12pt,halign=center,
  lower separated=false,halign lower=center,
  fontlower=\popsemi\fontsize{9}{11.5}\selectfont\color{popmuted},
  #1}
\newcommand{\legende}[1]{\par\vspace{6pt}{\popsemi\fontsize{9}{11.5}\selectfont\color{popmuted}#1}}

% ============================================================
% Signature OnBuch+
% ============================================================
\newcommand{\poplogoinline}[1]{{\poplogo\fontsize{#1}{#1}\selectfont\color{popink}OnBuch\color{poporange}+}}
\newcommand{\signatureonbuch}{%
  \begin{tcolorbox}[enhanced,colback=popink,colframe=popink,boxrule=2.2pt,arc=10pt,
      drop shadow=poporange,left=14pt,right=14pt,top=10pt,bottom=10pt,before skip=0pt,after skip=0pt]
    \tikz[baseline=-0.7ex]\node[circle,fill=popgreen,draw=popcream,line width=1.3pt,minimum size=22pt,inner sep=0]{\popcheck{white}};%
    \hspace{9pt}\begin{minipage}[c]{0.62\linewidth}
      {\popxbold\fontsize{12}{13}\selectfont\color{popcream}Signé\ \textbullet\ {\poplogo OnBuch\color{poporange}+}}\par\vspace{2pt}
      {\raggedright\popsemi\fontsize{8}{10.5}\selectfont\color{popline}\MakeUppercase{Équipe pédagogique OnBuch+}\\\MakeUppercase{Conforme au programme officiel MINESEC}\par}
    \end{minipage}\hfill
    {\poplogo\fontsize{22}{22}\selectfont\color{popcream}OnBuch\color{poporange}+}
  \end{tcolorbox}%
}

% ============================================================
% Page de garde
% #1 titre · #2 matière · #3 niveau · #4 module · #5 n° leçon · #6 durée ·
% #7 description / objectifs (texte ou itemize)
% ============================================================
\newcommand{\pagegarde}[7]{%
  \thispagestyle{empty}
  \newgeometry{a4paper,margin=1.7cm,top=1.6cm,bottom=1.4cm}%
  \begin{tikzpicture}[remember picture,overlay]
    \fill[poporange!18] ([xshift=-3.2cm,yshift=-3.6cm]current page.north east) circle (4.6cm);
    \fill[poppurple!14] ([xshift=2.4cm,yshift=7.6cm]current page.south west) circle (3.8cm);
  \end{tikzpicture}%
  {\poplogo\fontsize{30}{30}\selectfont\color{popink}OnBuch\color{poporange}+}\hfill
  \raisebox{0.6ex}{\poppill{Cours signé \textbullet\ Premium}{popink}{popcream}}\par
  \vspace{1pt}\popsquiggle{poppurple}\par
  \vspace{8pt}
  \begin{tcolorbox}[enhanced,colback=poporange,colframe=popink,boxrule=2.8pt,arc=13pt,
      drop shadow=popink,left=18pt,right=18pt,top=12pt,bottom=13pt,after skip=10pt]
    \poppill{#2 \textbullet\ #3}{popink}{popcream}\hspace{5pt}\poppill{#4}{white}{popink}\par\vspace{8pt}
    {\popdisplay\fontsize{14}{14}\selectfont\color{popink}LEÇON #5}\par\vspace{4pt}
    {\raggedright\hyphenpenalty=10000\exhyphenpenalty=10000\popdisplay\fontsize{25}{27}\selectfont\color{popcream}\MakeUppercase{#1}\par}
    \vspace{8pt}
    {\popxbold\fontsize{9.5}{9.5}\selectfont\color{popink}\MakeUppercase{Durée conseillée : #6}}
  \end{tcolorbox}
  \begin{tcolorbox}[enhanced,colback=white,colframe=popink,boxrule=2.2pt,arc=10pt,
      drop shadow=popink,left=16pt,right=16pt,top=10pt,bottom=10pt,after skip=8pt]
    \poppill{Dans cette leçon}{poppurple}{white}\par\vspace{5pt}
    \popsemi\fontsize{9.3}{12.6}\selectfont\color{popink}#7
  \end{tcolorbox}
  \noindent\begin{minipage}[t]{0.487\linewidth}
    \begin{tcolorbox}[enhanced,colback=popgreen,colframe=popink,boxrule=2.2pt,arc=10pt,
        drop shadow=popink,left=11pt,right=11pt,top=10pt,bottom=10pt,height=3.1cm,valign=center]
      \tcbox[colback=white,colframe=popink,boxrule=1.3pt,arc=5pt,boxsep=0pt,nobeforeafter,
        left=3pt,right=3pt,top=3pt,bottom=3pt,valign=center]{\qrcode[height=1.9cm]{https://play.google.com/store/apps/details?id=com.luvvix.onbuch&hl=fr}}%
      \hspace{8pt}\begin{minipage}[c]{0.5\linewidth}
        {\popdisplay\fontsize{11}{12.5}\selectfont\color{white}\MakeUppercase{Télécharge OnBuch}}\par\vspace{4pt}
        {\raggedright\popsemi\fontsize{8.4}{11}\selectfont\color{white}Gratuit \textbullet\ Google Play\\Tuteur IA, annales, concours, résultats\par}
      \end{minipage}
    \end{tcolorbox}
  \end{minipage}\hfill
  \begin{minipage}[t]{0.487\linewidth}
    \begin{tcolorbox}[enhanced,colback=poppurple,colframe=popink,boxrule=2.2pt,arc=10pt,
        drop shadow=popink,left=11pt,right=11pt,top=10pt,bottom=10pt,height=3.1cm,valign=center]
      \tikz[baseline=-0.7ex]\node[circle,fill=white,draw=popink,line width=1.3pt,minimum size=1.6cm,inner sep=0]{\popdisplay\fontsize{24}{24}\selectfont\color{poppurple}+};%
      \hspace{8pt}\begin{minipage}[c]{0.6\linewidth}
        {\popdisplay\fontsize{11}{12.5}\selectfont\color{white}\MakeUppercase{Passe à OnBuch+}}\par\vspace{4pt}
        {\raggedright\popsemi\fontsize{8.4}{11}\selectfont\color{white}Tous les cours signés, Tuteur IA illimité, fiches PDF — onglet OnBuch+ de l'app\par}
      \end{minipage}
    \end{tcolorbox}
  \end{minipage}
  \vspace{8pt}
  \popcontenu
  \vfill
  \signatureonbuch
  \restoregeometry
  \newpage
}

% Rangée « Ce cours contient » (page de garde)
\newcommand{\poptile}[3]{% #1 couleur #2 gros texte #3 légende
  \begin{tcolorbox}[enhanced,colback=#1,colframe=popink,boxrule=2pt,arc=9pt,shadow={2.5pt}{-2.5pt}{0pt}{popink},
      left=6pt,right=6pt,top=9pt,bottom=9pt,halign=center,valign=center,height=1.75cm,width=\linewidth,nobeforeafter]
    {\popdisplay\fontsize{12.5}{13}\selectfont\color{white}\MakeUppercase{#2}}\par\vspace{3pt}
    {\centering\popsemi\fontsize{7.8}{9.5}\selectfont\color{white}#3\par}
  \end{tcolorbox}}
\newcommand{\popcontenu}{%
  \par\noindent{\popxboldls\fontsize{7.5}{7.5}\selectfont\color{popmuted}CE COURS SIGNÉ CONTIENT}\par\vspace{7pt}
  \noindent\begin{minipage}[t]{0.235\linewidth}\poptile{popblue}{Cours}{complet, illustré, pas à pas}\end{minipage}\hfill
  \begin{minipage}[t]{0.235\linewidth}\poptile{poporange}{Méthodes}{et exercices résolus}\end{minipage}\hfill
  \begin{minipage}[t]{0.235\linewidth}\poptile{poppink}{Exercices}{type examen + corrigés}\end{minipage}\hfill
  \begin{minipage}[t]{0.235\linewidth}\poptile{popgreen}{Fiche bilan}{l'essentiel à retenir}\end{minipage}\par}

% Sommaire
\newtcolorbox{sommaire}{enhanced,colback=white,colframe=popink,boxrule=2.2pt,arc=10pt,
  drop shadow=popink,left=18pt,right=18pt,top=13pt,bottom=13pt,before skip=6pt,after skip=20pt,
  before upper={\poppill{Sommaire}{poppurple}{white}\par\vspace{9pt}\popsemi\fontsize{10.6}{14.5}\selectfont\color{popink}}}

% Signature de fin de cours — poussée tout en bas de la dernière page
\newcommand{\findecours}{%
  \par\vfill\nopagebreak
  \begin{center}\popsquiggle{poporange}\end{center}\vspace{4pt}
  \signatureonbuch
}
\AtBeginDocument{\def\llbracket{\mathopen{[\mkern-3.5mu[}}\def\rrbracket{\mathclose{]\mkern-3.5mu]}}} % intervalles d'entiers (glyphe ⟦ absent de la police)
% Glyphes absents de Fira Math : redéfinis à partir de glyphes disponibles
\AtBeginDocument{%
  \def\setminus{\mathbin{\backslash}}\let\smallsetminus\setminus
  \def\vdots{\mathord{\vbox{\baselineskip3.2pt\lineskiplimit0pt\kern4pt\hbox{.}\hbox{.}\hbox{.}}}}%
  \def\ddots{\mathinner{\mkern1mu\raise6.5pt\hbox{.}\mkern2mu\raise3.6pt\hbox{.}\mkern2mu\raise0.7pt\hbox{.}\mkern1mu}}%
  \def\triangle{\mathord{\scalebox{0.9}{$\symup{\Delta}$}}}%
  \def\nabla{\mathord{\raisebox{\depth}{\scalebox{1}[-1]{$\symup{\Delta}$}}}}%
  \def\checkmark{\popcheck{popdark}}%
  \def\star{\mathbin{\ast}}\def\bigstar{\mathord{\ast}}%
  \def\diamond{\mathbin{\scalebox{0.6}{$\Diamond$}}}%
  \def\bigcup{\mathop{\mathchoice{\raisebox{-0.3em}{\scalebox{1.7}{$\cup$}}}{\scalebox{1.25}{$\cup$}}{\cup}{\cup}}\displaylimits}%
  \def\bigcap{\mathop{\mathchoice{\raisebox{-0.3em}{\scalebox{1.7}{$\cap$}}}{\scalebox{1.25}{$\cap$}}{\cap}{\cap}}\displaylimits}%
  \def\lesseqqgtr{\mathrel{\lessgtr}}\def\bigoplus{\mathop{\oplus}\displaylimits}%
}
\providecommand{\ssi}{si et seulement si} % abréviation fréquente
\AtBeginDocument{\providecommand{\wideparen}{\overparen}} % arc de cercle

%% FIGFIX-BEGIN v5 — correctifs globaux des figures (docs/onbuch-generation/tools/figfix)
\usepackage{adjustbox}\usepackage{etoolbox}\tcbuselibrary{raster}
% toute figure TikZ/pgfplots plus large que la ligne est réduite à la largeur disponible
\BeforeBeginEnvironment{tikzpicture}{\begin{adjustbox}{max width=\linewidth}}
\AfterEndEnvironment{tikzpicture}{\end{adjustbox}}
% légendes pgfplots lisibles par-dessus les courbes
\pgfplotsset{every axis/.append style={legend style={fill=white,fill opacity=0.92,text opacity=1,draw=black!15,inner sep=3pt}}}
% étiquettes d'arêtes (syntaxe edge["..."]) avec fond blanc
\tikzset{every edge quotes/.append style={fill=white,inner sep=1.2pt}}
%% FIGFIX-END
\begin{document}

\pagegarde{L'emphase ; redacción : carta de motivación (120 palabras)}{Espagnol}{Tle A}{Módulo 4 : Activités économiques et monde du travail}{E18}{2 h}{Cette leçon mixte (2h) permet de maîtriser les mécanismes de l'emphase en espagnol (tournures clivées, ordre des mots, augmentatifs) pour valoriser son profil professionnel, et d'acquérir la méthodologie rigoureuse de rédaction d'une lettre de motivation formelle de 120 mots, exigée au Baccalauréat.
\vspace{4pt}
\begin{multicols}{2}\small
\begin{itemize}
\item Identifier et expliquer les procédés d'emphase (clivage, reprise pronominale, ser de relieve, ordre inversé, suffixes affectifs) dans un texte authentique.
\item Manipuler la structure syntaxique des phrases emphatiques en les transformant depuis une phrase neutre (thème) et vers le français (version).
\item Mobiliser le lexique spécialisé de l'emploi (contrat, compétences, entretien, postuler) dans un contexte formel.
\item Analyser la structure canonique d'une carta de motivación : en-tête, formule d'appel, corps (présentation, motivation, apport), formule de politesse finale.
\item Rédiger une lettre de motivation cohérente de 120 mots (±10\%) en respectant les conventions épistolaires hispaniques et en intégrant au moins trois procédés d'emphase.
\item S'auto-corriger sur les pièges classiques : confusion ser/estar dans l'emphase, accord du participe passé dans les clivées, registres de langue inappropriés.
\end{itemize}\end{multicols}}

\begin{sommaire}
\begin{multicols}{2}
\begin{enumerate}
\item 1. Texte support \& Compréhension : « La carta que cambió mi destino »
\item 2. Lexique thématique : Le monde de l'emploi (Campos semánticos)
\item 3. Grammaire : Les mécanismes de l'emphase (El relieve expresivo)
\item 4. Atelier d'analyse : Commentaire guidé de l'emphase dans le texte
\item 5. Expression écrite : Rédaction de la Carta de Motivación (120 mots)
\item 6. Synthèse, Entraînement Bac \& Auto-évaluation
\item Activité d'intégration
\item Méthodes et exercices résolus
\item Exercices
\item Corrigés des exercices
\item Fiche bilan
\end{enumerate}
\end{multicols}
\end{sommaire}

\begin{objectifs}
\begin{itemize}
\item Identifier et expliquer les procédés d'emphase (clivage, reprise pronominale, ser de relieve, ordre inversé, suffixes affectifs) dans un texte authentique.
\item Manipuler la structure syntaxique des phrases emphatiques en les transformant depuis une phrase neutre (thème) et vers le français (version).
\item Mobiliser le lexique spécialisé de l'emploi (contrat, compétences, entretien, postuler) dans un contexte formel.
\item Analyser la structure canonique d'une carta de motivación : en-tête, formule d'appel, corps (présentation, motivation, apport), formule de politesse finale.
\item Rédiger une lettre de motivation cohérente de 120 mots (±10\%) en respectant les conventions épistolaires hispaniques et en intégrant au moins trois procédés d'emphase.
\item S'auto-corriger sur les pièges classiques : confusion ser/estar dans l'emphase, accord du participe passé dans les clivées, registres de langue inappropriés.
\end{itemize}
\end{objectifs}

\begin{prerequis}
\begin{itemize}
\item Maîtrise des pronoms relatifs (que, quien, el cual) et des antéposés (lo que, lo cual).
\item Conjugaison et valeurs du subjonctif présent (volonté, doute, impersonnels) pour les formules de politesse.
\item Structure de base de la lettre formelle apprise en Première (en-tête, objet, formules de politesse simples).
\item Accord des adjectifs et participe passé (règle de l'accord avec 'tener' vs 'ser' dans les clivées).
\item Vocabulaire de base du monde du travail (trabajo, empresa, jefe, entrevista, currículum).
\end{itemize}
\end{prerequis}

\newpage

\coursec{Texte support \& Compréhension : « La carta que cambió mi destino »}

Imaginez la scène : vous êtes à Douala, fraîchement diplômé(e) de Terminale A, et vous rêvez de décrocher un stage dans une entreprise espagnole implantée au Cameroun ou dans une ONG hispanophone. Vous avez les compétences, la motivation, les idées. Mais voilà : sur le bureau du recruteur, votre dossier n'est qu'un document parmi des dizaines d'autres. Comment faire pour que le vôtre sorte du lot ? La réponse se joue dans deux détails décisifs : une lettre de motivation impeccablement structurée et l'art de souligner vos atouts grâce à l'\cle{emphase}, cette « force de frappe » de la langue espagnole qui transforme une phrase banale en une affirmation percutante. Comparez : \esp{Tengo experiencia en proyectos sociales} (« J'ai de l'expérience dans les projets sociaux ») et \esp{Lo que tengo es una amplia experiencia en proyectos sociales} (« Ce que j'ai, c'est une grande expérience des projets sociaux »). Même information, mais quel impact différent !

\textbf{Problématique :} comment une lettre de motivation bien construite, appuyée par des procédés d'emphase, peut-elle convaincre un recruteur en 120 mots ?

\textbf{Objectifs de la section :} à l'issue de cette première partie, vous saurez (1) identifier la structure en sept blocs d'une lettre de motivation, (2) reconnaître les marques du registre formel, (3) repérer les premières tournures emphatiques qui seront étudiées en grammaire à la section 3.

\courssub{Lecture silencieuse du texte support}

Lisez silencieusement la lettre suivante, écrite par un jeune Camerounais qui postule pour un poste de gestionnaire de projets dans une ONG à Madrid. Comptez environ 120 mots dans le corps de la lettre : c'est exactement la longueur exigée au Bac pour ce type de production.

\begin{textebox}[La carta que cambió mi destino]
Yaoundé, 15 de marzo de 2024\par
\medskip
Estimados Sres. del Departamento de Recursos Humanos:\par
\medskip
Les escribo para presentar mi candidatura al puesto de gestor de proyectos que su ONG ofrece en Madrid.\par
Soy Daniel Mbarga, licenciado en Cooperación Internacional. Lo que me define es mi compromiso con el desarrollo social: durante tres años desempeñé funciones de coordinador en proyectos educativos en Camerún.\par
Su organización me interesa precisamente porque fue usted quien lanzó el primer programa de becas para jóvenes africanos en mi país. Comparto sus valores y deseo aportar mi valor añadido: mi conocimiento del terreno y mi dominio del francés, del español y del inglés.\par
Me gustaría formar parte de su equipo. Quedo a la espera de sus noticias y me pongo a su disposición para una entrevista.\par
Atentamente,\par
Daniel Mbarga Nkoulou
\end{textebox}

\textbf{Glossaire :}
\begin{itemize}
\item \esp{desempeñar funciones} : exercer des fonctions, occuper un rôle. Synonyme courant : \esp{ejercer}.
\item \esp{compromiso (con)} : engagement (envers) ; l'adjectif \esp{comprometido} signifie « engagé, dévoué ».
\item \esp{valor añadido} : valeur ajoutée, ce que le candidat apporte en plus.
\item \esp{quedo a la espera de sus noticias} : je reste dans l'attente de vos nouvelles (formule de clôture formelle).
\item \esp{ponerse a su disposición} : se tenir à votre disposition.
\item \esp{el terreno} : le terrain (au sens figuré : la réalité locale, le travail de terrain).
\item \esp{el puesto} : le poste (de travail) ; attention, \esp{puesto} est aussi le participe passé de \esp{poner} et un nom signifiant « étal, stand » (\esp{un puesto de mercado}).
\item \esp{la beca} : la bourse (d'études) ; \esp{un programa de becas} : un programme de bourses.
\item \esp{licenciado en...} : licencié en..., titulaire d'une licence de...
\end{itemize}

\begin{definition}[La carta de motivación]
La \cle{carta de motivación} (lettre de motivation) est un document formel qui accompagne le CV (\esp{currículum vitae}) et dont le but est d'argumenter l'\cle{adéquation entre le profil du candidat et le poste visé}. Elle se distingue de la \esp{carta de presentación}, plus générale, qui sert simplement à se présenter sans répondre à une offre précise. La lettre de motivation, elle, répond à une annonce (\esp{oferta de empleo}) et doit citer le poste exact.
\end{definition}

\courssub{Compréhension globale : les informations clés}

Avant toute analyse fine, un bon lecteur repère d'abord les cinq informations essentielles d'une lettre professionnelle.

\begin{methode}[Lire une lettre professionnelle en 4 étapes]
\begin{enumerate}
\item \textbf{Identifier l'expéditeur et le destinataire} : qui écrit ? à qui ? (chercher la signature, l'en-tête, la formule d'appel).
\item \textbf{Trouver le poste visé} : il figure dans le premier paragraphe, souvent introduit par \esp{Les escribo para...} ; dans un courrier administratif, on trouve parfois une ligne \esp{Asunto:} ou \esp{Ref.:}.
\item \textbf{Isoler les trois arguments majeurs} du candidat (formation, expérience, qualités personnelles ou langues).
\item \textbf{Noter le niveau de langue} : formel (\esp{usted}, conditionnel de politesse) ou informel (\esp{tú}). Au Bac, une lettre de motivation est toujours formelle.
\end{enumerate}
\end{methode}

Appliquons cette méthode au texte :

\begin{center}
\begin{tabularx}{\linewidth}{|X|X|}\hline
\rowcolor{popblueL} \textbf{Information clé} & \textbf{Réponse dans le texte} \\ \hline
Expéditeur (\esp{remitente}) & Daniel Mbarga Nkoulou, licencié en Coopération Internationale \\ \hline
Destinataire (\esp{destinatario}) & Le Département des Ressources Humaines d'une ONG à Madrid \\ \hline
Poste visé (\esp{puesto}) & Gestionnaire de projets (\esp{gestor de proyectos}) \\ \hline
Arguments du candidat & 3 ans d'expérience de coordination ; connaissance du terrain ; trilinguisme \\ \hline
Disponibilité & Disponible pour un entretien (\esp{me pongo a su disposición}) \\ \hline
\end{tabularx}
\end{center}

\courssub{Repérage visuel : les indices de formalité}

Surlignez dans le texte (ou notez sur votre cahier) tous les indices qui prouvent que le registre est \cle{formel}. Vous devriez trouver :

\begin{itemize}
\item \textbf{Le vouvoiement} \esp{usted} : \esp{Les escribo} (le pronom \esp{les} renvoie à \esp{ustedes}, le Département), \esp{su ONG}, \esp{sus valores}, \esp{su disposición}.
\item \textbf{Le leísmo de courtoisie} : dans \esp{fue usted quien lanzó...}, le recours à \esp{usted} avec le verbe à la troisième personne marque la déférence ; de même, dans la correspondance espagnole, on emploie souvent \esp{le} (au lieu de \esp{lo}) pour une personne masculine : c'est la norme de politesse en Espagne.
\item \textbf{Le conditionnel de politesse} : \esp{Me gustaría formar parte de su equipo} — le conditionnel \esp{gustaría} adoucit la demande, là où l'indicatif \esp{quiero} sonnerait trop direct, voire arrogant.
\item \textbf{Les formules figées} : \esp{Estimados Sres.}, \esp{Quedo a la espera de sus noticias}, \esp{Atentamente}.
\end{itemize}

\begin{exemplebox}[Formules formelles traduites]
\esp{Les escribo para presentar mi candidatura.} « Je vous écris pour présenter ma candidature. »\par
\esp{Me gustaría formar parte de su equipo.} « J'aimerais faire partie de votre équipe. »\par
\esp{Quedo a la espera de sus noticias.} « Je reste dans l'attente de vos nouvelles. »\par
\esp{Me pongo a su disposición para una entrevista.} « Je me tiens à votre disposition pour un entretien. »
\end{exemplebox}

\begin{attention}[Estimado Señor ou Estimado Sr. García ?]
Erreur fréquente au Bac : confondre \esp{Estimado Señor} (sans nom, quand on ne connaît pas le destinataire) et \esp{Estimado Sr. García} (avec le nom, abrégé \esp{Sr.}). Si le nom du destinataire est inconnu, n'inventez rien : utilisez \esp{Estimados Sres.} ou \esp{A la atención del Departamento de Recursos Humanos}. Autre piège : ne jamais écrire \esp{Estimado amigo} dans une lettre de motivation — c'est un registre amical, catastrophique ici ! Enfin, n'oubliez pas les \textbf{deux points} après la formule d'appel en espagnol : \esp{Estimados Sres.:} (et non une virgule comme en français).
\end{attention}

\courssub{Compréhension détaillée : QCM et questions ouvertes}

\begin{exoresolu}[QCM de compréhension]
\textbf{Énoncé.} Choisissez la bonne réponse.\par
1. Daniel postule pour un poste de : a) traducteur ; b) gestor de proyectos ; c) profesor.\par
2. Il a travaillé pendant : a) un an ; b) deux ans ; c) tres años.\par
3. Ce qui l'attire dans l'ONG, c'est : a) le salaire ; b) son programme de bourses ; c) sa taille.\par
4. Daniel parle : a) deux langues ; b) tres idiomas ; c) seulement l'espagnol.\par
5. À la fin, il : a) exige une réponse rapide ; b) se propose pour une entrevista ; c) envoie son CV plus tard.
\tcblower
\textbf{Solution détaillée.}\par
1. \textbf{b)} — \esp{mi candidatura al puesto de gestor de proyectos} (paragraphe 1).\par
2. \textbf{c)} — \esp{durante tres años desempeñé funciones de coordinador} (paragraphe 2).\par
3. \textbf{b)} — \esp{fue usted quien lanzó el primer programa de becas} (paragraphe 3).\par
4. \textbf{b)} — \esp{mi dominio del francés, del español y del inglés} : trois langues (paragraphe 3).\par
5. \textbf{b)} — \esp{me pongo a su disposición para una entrevista} (paragraphe 4). La réponse a) est fausse : le ton reste courtois, jamais exigeant.
\end{exoresolu}

\begin{exercice}[Questions ouvertes]
Répondez en espagnol, par des phrases complètes.\par
1. ¿Qué estudios tiene Daniel Mbarga?\par
2. ¿Qué experiencia profesional aporta al puesto?\par
3. ¿Por qué le interesa precisamente esa ONG?\par
4. ¿Qué lenguas domina y por qué es una ventaja para una ONG internacional?\par
5. ¿Qué formas de cortesía utiliza el candidato? Cite dos ejemplos del texto.
\end{exercice}

\begin{corrige}[Exercice : questions ouvertes]
1. \esp{Es licenciado en Cooperación Internacional.}\par
2. \esp{Aporta tres años de experiencia como coordinador de proyectos educativos en Camerún.}\par
3. \esp{Le interesa porque esa ONG lanzó el primer programa de becas para jóvenes africanos en su país y porque comparte sus valores.}\par
4. \esp{Domina el francés, el español y el inglés; es una ventaja porque una ONG internacional trabaja con países y socios de distintas lenguas.}\par
5. \esp{Utiliza el tratamiento de «usted» («Les escribo», «su equipo») y el condicional de cortesía («Me gustaría»), además de la fórmula final «Quedo a la espera de sus noticias».}
\end{corrige}

\courssub{Analyse du ton : humble mais assertif}

Une lettre de motivation réussie marche sur une ligne de crête : ni arrogante (\esp{Soy el mejor candidato} serait désastreux), ni timide au point de s'effacer. Observez comment Daniel dose son discours :

\begin{itemize}
\item \textbf{Assertivité par l'emphase} : \esp{Lo que me define es mi compromiso} (tournure clivée : « ce qui me définit, c'est mon engagement ») et \esp{fue usted quien lanzó...} (emphase sur le destinataire : « c'est vous qui avez lancé... »). Ces structures seront étudiées en détail dans la section 3.
\item \textbf{Humilité par le conditionnel} : \esp{Me gustaría}, \esp{deseo aportar} — le candidat propose, il n'impose pas.
\item \textbf{Orientation vers l'apport} : le candidat ne dit pas « que pouvez-vous m'offrir ? » mais \esp{deseo aportar mi valor añadido} : il met en avant ce qu'il peut \emph{donner} à l'organisation. C'est la règle d'or du genre : parler des besoins de l'employeur, pas seulement des siens.
\end{itemize}

\begin{aretenir}[L'équilibre du ton]
Lettre réussie = \cle{emphase} pour affirmer ses atouts + \cle{conditionnel de politesse} pour rester humble + vocabulaire de l'\cle{apport} (\esp{aportar, contribuir, colaborar}) plutôt que de la demande (\esp{quiero, necesito, exijo}).
\end{aretenir}

\courssub{Schéma de la structure du texte}

La lettre modèle suit une architecture en sept blocs, à reproduire dans vos propres rédactions :

\begin{popfigure}
\begin{center}
\begin{tikzpicture}[
  node distance=0.32cm,
  bloc/.style={rectangle, rounded corners=3pt, draw=popblue, fill=popblueL, text width=6.6cm, align=center, font=\small, inner sep=4pt}
]
\node[bloc, fill=popgoldL, draw=popgold] (a) {\textbf{En-tête} : lugar y fecha (Yaoundé, 15 de marzo de 2024)};
\node[bloc, below=of a, fill=poporangeL, draw=poporange] (b) {\textbf{Appel} : Estimados Sres. del Dpto. de RR.HH.:};
\node[bloc, below=of b] (c) {\textbf{§1 Qui suis-je ?} : presentación + puesto solicitado};
\node[bloc, below=of c] (d) {\textbf{§2 Pourquoi vous ?} : interés por la ONG y sus valores};
\node[bloc, below=of d] (e) {\textbf{§3 Pourquoi moi ?} : experiencia, idiomas, valor añadido};
\node[bloc, below=of e, fill=poppurpleL, draw=poppurple] (f) {\textbf{Disponibilité} : a su disposición para una entrevista};
\node[bloc, below=of f, fill=popgreenL, draw=popgreen] (g) {\textbf{Politesse + Signature} : Atentamente, Daniel Mbarga};
\foreach \x/\y in {a/b, b/c, c/d, d/e, e/f, f/g}{
  \draw[-{Stealth}, thick, popblue] (\x) -- (\y);
}
\end{tikzpicture}
\end{center}
\legende{Organigramme de la structure de la lettre de motivation : sept blocs obligatoires, de l'en-tête à la signature.}
\end{popfigure}

\begin{savaistu}[Le leísmo de courtoisie]
En Espagne, on dit couramment \esp{Le conozco bien} (avec \esp{le}) pour parler d'un homme, alors que la norme latino-américaine préfère \esp{Lo conozco bien}. Ce phénomène, appelé \cle{leísmo}, est accepté par la Real Academia lorsque le pronom désigne une personne masculine : c'est une marque de politesse très appréciée dans la correspondance formelle espagnole. En Guinée équatoriale, seul pays hispanophone d'Afrique subsaharienne, la lettre formelle suit les mêmes conventions qu'en Espagne — un atout culturel pour les candidats camerounais hispanisants qui postulent à Malabo ou à Bata !
\end{savaistu}

\begin{exercice}[Repérage dans le texte]
1. Relevez dans la lettre les deux tournures emphatiques et réécrivez-les sans emphase (phrase neutre).\par
2. Relevez trois formules de politesse formelle et traduisez-les en français.\par
3. Pourquoi le candidat cite-t-il le programme de bourses de l'ONG ? Quel effet cela produit-il sur le lecteur ?
\end{exercice}

\begin{corrige}[Exercice : repérage]
1. \esp{Lo que me define es mi compromiso} $\rightarrow$ forme neutre : \esp{Me define mi compromiso}. \esp{Fue usted quien lanzó el programa} $\rightarrow$ forme neutre : \esp{Usted lanzó el programa}.\par
2. \esp{Estimados Sres.} « Messieurs » ; \esp{Quedo a la espera de sus noticias} « Je reste dans l'attente de vos nouvelles » ; \esp{Atentamente} « Cordialement / Veuillez agréer... ».\par
3. En citant le programme de bourses, le candidat prouve qu'il connaît réellement l'organisation et qu'il a un lien personnel avec elle : cela flatte le destinataire (emphase sur \esp{usted}) et montre une motivation authentique, pas une candidature envoyée au hasard.
\end{corrige}

Vous savez désormais lire une lettre de motivation comme un recruteur : structure en sept blocs, registre formel, ton humble mais assertif. La section suivante vous donnera les clés lexicales du monde de l'emploi pour enrichir vos propres productions.

\coursec{Lexique thématique : Le monde de l'emploi (Campos semánticos)}

Dans la section précédente, nous avons lu et commenté \esp{« La carta que cambió mi destino »}, la lettre de motivation grâce à laquelle une jeune diplômée de Yaoundé a décroché un poste dans une entreprise hispanophone. Vous avez remarqué que cette lettre regorge de mots précis : \esp{puesto}, \esp{contrato}, \esp{experiencia}, \esp{habilidades}, \esp{entrevista}. Pour rédiger vous-même une lettre convaincante de 120 mots (section 5), il faut d'abord \cle{maîtriser le lexique de l'emploi}. C'est l'objet de cette section : observer, classer, mémoriser et manipuler le vocabulaire du monde du travail, en évitant les pièges qui guettent les francophones.

\courssub{Observer avant de classer : le lexique en contexte}

\begin{textebox}[Anuncio de empleo — Oficina de empleo de Malabo]
\textbf{SE BUSCA: AUXILIAR ADMINISTRATIVO/A}

Empresa de comercio internacional con sede en Malabo busca auxiliar administrativo con experiencia previa de al menos dos años. Se requiere dominio del español y del francés, capacidad de organización y disponibilidad inmediata.

Ofrecemos: contrato indefinido tras un período de prueba de tres meses, jornada completa, sueldo competitivo y posibilidades de promoción dentro de la empresa.

Los candidatos interesados deben enviar su currículum y una carta de motivación antes del día 30 del mes en curso. Los preseleccionados serán convocados a una entrevista personal con el jefe de recursos humanos.

No se atenderán solicitudes incompletas ni enviadas fuera de plazo.
\end{textebox}

\begin{definition}[Glossaire de l'annonce]
\begin{itemize}
\item \esp{se busca} : on recherche (formule impersonnelle typique des annonces)
\item \esp{con sede en} : dont le siège est à
\item \esp{se requiere} : on exige, il est requis
\item \esp{período de prueba} : période d'essai
\item \esp{jornada completa} : temps plein
\item \esp{los preseleccionados} : les candidats présélectionnés
\item \esp{fuera de plazo} : hors délai
\item \esp{recursos humanos (RR. HH.)} : les ressources humaines
\end{itemize}
\end{definition}

\textbf{Observation.} Soulignez dans l'annonce tous les mots qui désignent : (a) la recherche du candidat, (b) le type de contrat, (c) les qualités exigées, (d) les étapes de la sélection. Vous obtenez naturellement quatre groupes : c'est exactement la classification en \cle{champs lexicaux} que nous allons adopter.

\courssub{Les quatre champs lexicaux de l'emploi}

\begin{definition}[Campo semántico]
Un \cle{champ sémantique} (\esp{campo semántico}) est un ensemble de mots unis par une même idée générale. Classer le vocabulaire en champs facilite la mémorisation et garantit la richesse lexicale exigée au Bac (critère « riqueza y corrección del léxico »).
\end{definition}

\begin{center}
\begin{tabularx}{\linewidth}{|X|X|X|}
\hline
\rowcolor{popblueL} \textbf{Champ} & \textbf{Vocabulaire ES} & \textbf{Français} \\
\hline
1. \esp{El reclutamiento} & \esp{oferta de empleo, anuncio, candidato, solicitud, currículum, carta de motivación, presentarse a un puesto} & offre d'emploi, annonce, candidat, candidature, CV, lettre de motivation, postuler \\
\hline
2. \esp{El contrato y el estatuto} & \esp{contrato indefinido, contrato temporal, período de prueba, jornada completa/parcial, sueldo, salario, retribución, despido, ascenso} & CDI, CDD, période d'essai, temps plein/partiel, salaire, rémunération, licenciement, promotion \\
\hline
3. \esp{Las competencias y cualidades} & \esp{experiencia previa, habilidades, capacidad de organización, dominio de idiomas, puntualidad, responsabilidad, trabajo en equipo} & expérience préalable, compétences, sens de l'organisation, maîtrise des langues, ponctualité, responsabilité, travail en équipe \\
\hline
4. \esp{La entrevista y el seguimiento} & \esp{entrevista personal, preseleccionado, convocar, carta de agradecimiento, respuesta, plazo} & entretien individuel, présélectionné, convoquer, lettre de remerciement, réponse, délai \\
\hline
\end{tabularx}
\end{center}

\begin{popfigure}
\begin{tikzpicture}[mindmap, grow cyclic, every node/.style={concept, text=white, align=center, font=\small},
root concept/.append style={concept color=popblue, font=\bfseries},
level 1/.append style={level distance=3.4cm, sibling angle=90},
level 2/.append style={level distance=2.4cm, sibling angle=45, font=\scriptsize}]
\node[root concept]{EMPLEO}
  child[concept color=poporange]{ node {Reclutamiento}
      child { node[concept color=poporange!70, align=center]{oferta de empleo\\ offre d'emploi} }
      child { node[concept color=poporange!70, align=center]{solicitud\\ candidature} } }
  child[concept color=popgreen]{ node {Contrato}
      child { node[concept color=popgreen!70, align=center]{indefinido\\ CDI} }
      child { node[concept color=popgreen!70, align=center]{temporal\\ CDD} } }
  child[concept color=poppurple]{ node {Competencias}
      child { node[concept color=poppurple!70, align=center]{experiencia previa\\ expérience} }
      child { node[concept color=poppurple!70, align=center]{habilidades\\ compétences} } }
  child[concept color=poppink]{ node {Entrevista}
      child { node[concept color=poppink!70, align=center]{convocar\\ convoquer} }
      child { node[concept color=poppink!70, align=center]{plazo\\ délai} } };
\end{tikzpicture}
\legende{Carte mentale du lexique de l'emploi : quatre champs sémantiques à mémoriser pour le Bac.}
\end{popfigure}

\courssub{Tableau de référence : 20 mots indispensables pour le Bac}

\begin{center}
\begin{tabularx}{\linewidth}{|X|X|X|}
\hline
\rowcolor{popblueL} \textbf{Mot ES} & \textbf{Mot FR} & \textbf{Exemple en contexte} \\
\hline
\esp{el puesto (de trabajo)} & le poste & \esp{Solicito el puesto de secretaria bilingüe.} \\
\hline
\esp{la oferta de empleo} & l'offre d'emploi & \esp{He leído su oferta de empleo en el periódico.} \\
\hline
\esp{el/la candidato/a} & le/la candidat(e) & \esp{Los candidatos deben hablar francés.} \\
\hline
\esp{la solicitud} & la candidature (dossier) & \esp{Envío mi solicitud antes del plazo.} \\
\hline
\esp{el currículum (vitae)} & le CV & \esp{Adjunto mi currículum a esta carta.} \\
\hline
\esp{la carta de motivación} & la lettre de motivation & \esp{Redacté una carta de motivación clara.} \\
\hline
\esp{el contrato indefinido} & le CDI & \esp{Busco un contrato indefinido.} \\
\hline
\esp{el contrato temporal} & le CDD & \esp{Acepté un contrato temporal de seis meses.} \\
\hline
\esp{el período de prueba} & la période d'essai & \esp{El período de prueba dura tres meses.} \\
\hline
\esp{la jornada completa / parcial} & le temps plein / partiel & \esp{Prefiero una jornada parcial.} \\
\hline
\esp{el sueldo} & le salaire (fixe mensuel) & \esp{El sueldo se paga a fin de mes.} \\
\hline
\esp{la retribución} & la rémunération (globale) & \esp{La retribución incluye las primas.} \\
\hline
\esp{la experiencia previa} & l'expérience préalable & \esp{Tengo experiencia previa en el sector.} \\
\hline
\esp{las habilidades} & les compétences & \esp{Desarrollé nuevas habilidades informáticas.} \\
\hline
\esp{el trabajo en equipo} & le travail en équipe & \esp{Valoro mucho el trabajo en equipo.} \\
\hline
\esp{la entrevista (personal)} & l'entretien & \esp{Me convocaron a una entrevista personal.} \\
\hline
\esp{el jefe de recursos humanos} & le chef du personnel / DRH & \esp{El jefe de recursos humanos me recibió.} \\
\hline
\esp{el plazo} & le délai & \esp{Cumplí todos los plazos del proyecto.} \\
\hline
\esp{el ascenso} & la promotion & \esp{Obtuvo un ascenso tras dos años.} \\
\hline
\esp{el despido} & le licenciement & \esp{El despido fue injustificado.} \\
\hline
\end{tabularx}
\end{center}

\begin{savaistu}[« Currículum » : un mot à part]
En Espagne, on dit le plus souvent \esp{el currículum} (sans \esp{vitae}) ou simplement \esp{el CV}. Attention au pluriel : les formes admises sont \esp{los currículums} ou \esp{los currículos} ; à l'écrit formel, on préfère souvent \esp{los currículos}. Au Cameroun et en Guinée équatoriale, la photo sur le CV est fréquente ; en Espagne, elle n'est pas obligatoire mais reste courante. Autre détail : on \esp{adjunta} (joint) son CV à une lettre, on ne l'« envoie » pas seul dans une candidature formelle.
\end{savaistu}

\courssub{Les collocations indispensables (verbe + nom)}

En espagnol comme en français, certains verbes « s'attachent » à certains noms : ce sont les \cle{collocations}. Les traduire mot à mot depuis le français produit des phrases incorrectes ou artificielles.

\begin{propriete}[Règle des collocations de l'emploi]
En espagnol :
\begin{itemize}
\item on \esp{adquiere} ou \esp{tiene} experiencia — jamais *\esp{gana experiencia} (calque de « gagner de l'expérience ») ;
\item on \esp{desarrolla} habilidades ; pour un projet, on dit \esp{desarrollar un proyecto} ou \esp{poner en marcha un proyecto} (lancer un projet) ;
\item on \esp{cumple} objetivos, plazos, requisitos (atteindre des objectifs, respecter des délais, remplir des conditions) ;
\item on \esp{firma} un contrato, on \esp{presenta} una solicitud, on \esp{ocupa} un puesto, on \esp{supera} una entrevista (réussir un entretien).
\end{itemize}
\end{propriete}

\begin{exemplebox}[Collocations en phrases]
\esp{Firmé el contrato tras el período de prueba.} « J'ai signé le contrat après la période d'essai. »

\esp{Adquirí experiencia trabajando en un banco de Douala.} « J'ai acquis de l'expérience en travaillant dans une banque de Douala. »

\esp{Desarrollé mis habilidades lingüísticas durante mis estudios.} « J'ai développé mes compétences linguistiques pendant mes études. »

\esp{Siempre cumplo los plazos que me fijan.} « Je respecte toujours les délais qu'on me fixe. »

\esp{Superó la entrevista y ocupó el puesto.} « Elle a réussi l'entretien et obtenu le poste. »
\end{exemplebox}

\begin{attention}[Sueldo, salario, retribución : ne pas confondre]
\esp{El sueldo} = le salaire fixe versé chaque mois. \esp{El salario} = la rémunération au sens large, souvent calculée à l'heure ou à la journée. Pour « rémunération » au sens juridique et formel (celui du Bac et des lettres), le terme le plus sûr est \esp{la retribución}, qui inclut le fixe et les primes (\esp{las primas}). À retenir aussi : \esp{el salario mínimo} = le salaire minimum (en Espagne, on parle du \esp{SMI}, \esp{salario mínimo interprofesional}).
\end{attention}

\courssub{Faux-amis et pièges lexicaux}

\begin{attention}[Les faux-amis du monde du travail]
\begin{itemize}
\item \esp{éxito} = \textbf{succès} (la « sortie » se dit \esp{salida}) : \esp{tener éxito} = réussir, avoir du succès.
\item \esp{empresa} = \textbf{entreprise} (un « empire » se dit \esp{imperio}).
\item \esp{pretender} = \textbf{briguer, postuler, avoir l'intention de} (et non « prétendre » au sens d'affirmer : \esp{afirmar}, \esp{sostener}). \esp{Pretendo este puesto} = je brigue ce poste.
\item \esp{demandar} = exiger, réclamer, poursuivre en justice (la « demande » d'emploi = \esp{la solicitud}).
\item \esp{formación} = formation (et non « information » : \esp{información}).
\item \esp{actualmente} = actuellement, en ce moment (et non « en fait » : \esp{en realidad}).
\item \esp{asistir a} = assister à, être présent à (et non « aider » : \esp{ayudar}).
\item \esp{una entrevista} = un entretien d'embauche ou une interview journalistique ; « interviewer » = \esp{entrevistar}.
\end{itemize}
\end{attention}

\begin{exemplebox}[Exemples traduits — phrases clés de la lettre]
\esp{Tengo experiencia previa en el sector.} « J'ai une expérience préalable dans le secteur. »

\esp{Busco un contrato indefinido.} « Je cherche un CDI (contrat à durée indéterminée). »

\esp{Pretendo el puesto de auxiliar administrativo anunciado por su empresa.} « Je brigue le poste d'assistant administratif annoncé par votre entreprise. »

\esp{Mi formación y mi dominio del francés me permiten cumplir los objetivos del puesto.} « Ma formation et ma maîtrise du français me permettent d'atteindre les objectifs du poste. »
\end{exemplebox}

\courssub{Extension Bac : télétravail et contrats précaires}

Les sujets récents du Bac aiment les thèmes de société : le télétravail, la précarité, les stages. Voici le lexique d'extension à connaître.

\begin{center}
\begin{tabularx}{\linewidth}{|X|X|X|}
\hline
\rowcolor{popblueL} \textbf{Mot ES} & \textbf{Français} & \textbf{Exemple} \\
\hline
\esp{el teletrabajo / trabajar a distancia} & le télétravail & \esp{El teletrabajo aumentó después de la pandemia.} \\
\hline
\esp{el/la becario/a} & le/la stagiaire & \esp{Empecé como becaria en una ONG.} \\
\hline
\esp{las prácticas (en empresa)} & le stage (en entreprise) & \esp{Hice unas prácticas de tres meses.} \\
\hline
\esp{el contrato temporal} & le contrat précaire / CDD & \esp{Encadena contratos temporales desde 2020.} \\
\hline
\esp{la precariedad laboral} & la précarité professionnelle & \esp{Los jóvenes sufren la precariedad laboral.} \\
\hline
\esp{trabajar por cuenta propia} & travailler à son compte & \esp{Muchos jóvenes trabajan por cuenta propia.} \\
\hline
\esp{el desempleo / el paro} & le chômage & \esp{La tasa de desempleo juvenil es alta.} \\
\hline
\esp{la conciliación laboral y familiar} & la conciliation vie pro/vie privée & \esp{El teletrabajo mejora la conciliación.} \\
\hline
\end{tabularx}
\end{center}

\begin{savaistu}[Le stage en espagnol : « prácticas » et « becario »]
En Espagne, l'étudiant en stage est \esp{un becario} et son stage \esp{unas prácticas} (presque toujours au pluriel). Le mot \esp{becario} vient de \esp{beca} (bourse) : à l'origine, le stagiaire recevait une petite bourse. Attention : \esp{hacer unas prácticas} et non *\esp{hacer un stage} — le mot « stage » n'existe pas en espagnol !
\end{savaistu}

\courssub{Exercices de mise en contexte}

\begin{exoresolu}[Phrases à trous : les collocations]
Complétez avec : \esp{adquirí — firmé — cumplo — desarrollé — presenté}.

1. \esp{\ldots\ldots mi solicitud antes del plazo.}
2. \esp{\ldots\ldots experiencia como vendedor en un mercado de Yaoundé.}
3. \esp{\ldots\ldots el contrato después de la entrevista.}
4. \esp{\ldots\ldots mis habilidades informáticas en la universidad.}
5. \esp{Siempre \ldots\ldots los plazos que me fijan.}
\tcblower
\textbf{Corrigé détaillé.}
1. \esp{Presenté mi solicitud antes del plazo.} — on « présente » une candidature (\esp{presentar una solicitud}).
2. \esp{Adquirí experiencia como vendedor\ldots} — l'expérience « s'acquiert » (\esp{adquirir experiencia}), elle ne se « gagne » pas.
3. \esp{Firmé el contrato\ldots} — collocation : \esp{firmar un contrato}.
4. \esp{Desarrollé mis habilidades\ldots} — collocation : \esp{desarrollar habilidades}.
5. \esp{Siempre cumplo los plazos\ldots} — collocation : \esp{cumplir plazos}.
\end{exoresolu}

\begin{exercice}[Association définition / mot]
Associez chaque définition au mot correspondant : \esp{becario — despido — período de prueba — retribución — ascenso}.

1. Ensemble de ce que reçoit le salarié, primes comprises.
2. Période initiale pendant laquelle l'employeur évalue le nouvel employé.
3. Étudiant qui fait un stage en entreprise.
4. Passage à un poste supérieur dans la hiérarchie.
5. Rupture du contrat décidée par l'employeur.
\end{exercice}

\begin{corrige}[Exercice : association définition / mot]
1. \esp{retribución} ; 2. \esp{período de prueba} ; 3. \esp{becario} ; 4. \esp{ascenso} ; 5. \esp{despido}.
\end{corrige}

\begin{exercice}[Du CV à la lettre : reformulation]
Voici trois lignes télégraphiques d'un CV. Transformez-les en phrases complètes et élégantes pour une lettre de motivation, en utilisant les collocations étudiées.

1. \esp{Experiencia: 2 años, tienda de ropa, Douala.}
2. \esp{Idiomas: español, francés, inglés.}
3. \esp{Informática: Word, Excel.}
\end{exercice}

\begin{corrige}[Exercice : reformulation CV $\rightarrow$ lettre]
1. \esp{Adquirí una sólida experiencia de dos años como dependiente en una tienda de ropa de Douala, donde aprendí a atender a los clientes.} « J'ai acquis une solide expérience de deux ans comme vendeur dans un magasin de vêtements de Douala, où j'ai appris à servir les clients. »

2. \esp{Domino el español, el francés y el inglés, lo que me permite comunicarme con clientes de distintos países.} « Je maîtrise l'espagnol, le français et l'anglais, ce qui me permet de communiquer avec des clients de différents pays. »

3. \esp{Desarrollé mis habilidades informáticas, especialmente en Word y Excel, durante mi formación.} « J'ai développé mes compétences informatiques, notamment sur Word et Excel, pendant ma formation. »

Remarquez la méthode : nom du CV $\rightarrow$ verbe de collocation + nom + détail valorisant. C'est exactement ce que l'on attend dans une lettre de motivation au Bac.
\end{corrige}

\begin{aretenir}[L'essentiel du lexique de l'emploi]
\begin{itemize}
\item Quatre champs sémantiques : \esp{reclutamiento}, \esp{contrato}, \esp{competencias}, \esp{entrevista}.
\item Collocations obligatoires : \esp{adquirir experiencia}, \esp{desarrollar habilidades}, \esp{cumplir plazos/objetivos}, \esp{firmar un contrato}, \esp{presentar una solicitud}, \esp{superar una entrevista}.
\item Faux-amis à bannir : \esp{éxito} (succès), \esp{empresa} (entreprise), \esp{pretender} (briguer).
\item Terme noble pour « rémunération » : \esp{la retribución}.
\item Extension Bac : \esp{teletrabajo}, \esp{becario}, \esp{prácticas}, \esp{contrato temporal/indefinido}, \esp{jornada parcial/completa}.
\end{itemize}
\end{aretenir}

Ce lexique désormais classé et manipulé, nous pouvons passer à la section suivante : la grammaire de l'\cle{emphase}, qui vous permettra de mettre en valeur ce vocabulaire dans vos phrases — car une bonne lettre ne dit pas seulement \esp{tengo experiencia}, elle dit \esp{lo que tengo es experiencia}.

\coursec{Grammaire : Les mécanismes de l'emphase (El relieve expresivo)}

Dans la section précédente, nous avons constitué le lexique du monde de l'emploi : \esp{el puesto de trabajo}, \esp{la experiencia}, \esp{el currículum}, \esp{la entrevista}... Mais un vocabulaire riche ne suffit pas pour convaincre. Dans la lettre que nous avons lue, \emph{La carta que cambió mi destino}, la narratrice ne dit pas simplement les choses : elle les \textbf{souligne}, elle les \textbf{met en relief}. C'est précisément ce que la grammaire espagnole appelle \esp{el relieve expresivo} ou \esp{el énfasis}. Cette section est le cœur grammatical de la leçon : nous allons démonter cinq mécanismes qui permettent de passer d'une phrase neutre à une phrase emphatique, puis nous verrons lesquels sont adaptés au registre formel de la lettre de motivation.

\courssub{Observation guidée : du neutre à l'emphatique}

Reprenons cinq phrases extraites de notre texte support et comparons-les à leur version neutre.

\begin{exemplebox}[Cinq phrases du texte sous la loupe]
\begin{enumerate}
\item \esp{Fue mi abuela quien me enseñó a escribir cartas.} (neutre : \esp{Mi abuela me enseñó a escribir cartas.}) « C'est ma grand-mère qui m'a appris à écrire des lettres. »
\item \esp{Lo que yo buscaba era un puesto estable.} (neutre : \esp{Yo buscaba un puesto estable.}) « Ce que je cherchais, c'était un poste stable. »
\item \esp{A mí no me asustaban las entrevistas.} (neutre : \esp{No me asustaban las entrevistas.}) « Moi, les entretiens ne m'effrayaient pas. »
\item \esp{Muchas puertas llamé antes de encontrar empleo.} (neutre : \esp{Llamé muchas puertas antes de encontrar empleo.}) « Bien des portes, j'en ai frappé avant de trouver un emploi. »
\item \esp{Recibí una cartita breve pero decisiva.} (neutre : \esp{Recibí una carta breve pero decisiva.}) « J'ai reçu une petite lettre, brève mais décisive. »
\end{enumerate}
\end{exemplebox}

\begin{methode}[Observer avant de conclure]
\begin{enumerate}
\item Souligne dans chaque phrase emphatique l'élément mis en relief (personne, action, objet, circonstance).
\item Compare avec la version neutre : qu'est-ce qui a changé ? (un mot ajouté ? un déplacement ? une terminaison ?)
\item Demande-toi quel effet produit ce changement : insistance, contraste, émotion, élégance ?
\item Déduis-en le nom du procédé : c'est ce que nous allons formaliser maintenant.
\end{enumerate}
\end{methode}

\courssub{Règle 1 : Les tournures clivées (oraciones hendidas)}

La tournure clivée « coupe » la phrase en deux pour isoler un élément : \esp{Es / Fue + élément focalisé + que / quien / donde / cuando + reste de la phrase}. C'est l'équivalent presque exact du français « c'est... qui / que ».

\begin{propriete}[Structure de la clivée espagnole]
\textbf{Forme} : \esp{SER + SN (élément focalisé) + quien/que/donde/cuando + proposition}.
\begin{itemize}
\item Sujet personne focalisé : \esp{quien} (sg) / \esp{quienes} (pl) : \esp{Fue ella quien me contrató.} « C'est elle qui m'a embauchée. »
\item Objet personne : \esp{a quien} : \esp{Es a ti a quien buscan.} « C'est toi qu'ils cherchent. »
\item Objet chose ou circonstance : \esp{que} : \esp{Es la experiencia que más valoran.} « C'est l'expérience qu'ils valorisent le plus. »
\item Lieu et temps : \esp{donde}, \esp{cuando} : \esp{Fue en Douala donde conseguí mi primer empleo.} « C'est à Douala que j'ai décroché mon premier emploi. » ; \esp{Fue ayer cuando llegó la respuesta.} « C'est hier que la réponse est arrivée. »
\item Le temps de \esp{ser} suit la logique du récit : présent (\esp{es}) pour le présent ou le général, passé (\esp{fue}) pour le passé.
\end{itemize}
\end{propriete}

\begin{propriete}[Accord du verbe \esp{ser} dans la clivée]
Le verbe \esp{ser} se met à la 3\textsuperscript{e} personne (sg ou pl selon le SN focalisé) : \esp{Son los jóvenes quienes más sufren el desempleo.} « Ce sont les jeunes qui souffrent le plus du chômage. »

\textbf{Exception capitale} : si l'élément focalisé est un \cle{pronom personnel sujet} (\esp{yo, tú, nosotros, vosotros}), \esp{ser} s'accorde avec ce pronom :
\begin{itemize}
\item \esp{Soy yo quien lo he hecho.} « C'est moi qui l'ai fait. »
\item \esp{Eres tú quien debe decidir.} « C'est toi qui dois décider. »
\item \esp{Sois vosotros quienes ganasteis el concurso.} « C'est vous qui avez gagné le concours. »
\end{itemize}
\end{propriete}

\begin{attention}[Le piège n°1 du francophone]
Ne traduis jamais mot à mot « C'est moi qui l'ai fait » par *\esp{Es yo quien lo ha hecho} : double faute ! L'espagnol exige l'accord avec \esp{yo} : \esp{Soy yo quien lo he hecho} ou, au passé, \esp{Fui yo quien lo hizo}. De même : « C'est nous qui... » = \esp{Somos nosotros quienes...}, jamais *\esp{Es nosotros...}.
\end{attention}

\courssub{Règle 2 : \esp{Ser} de relieve (oración de relieve)}

\begin{definition}[\esp{Ser} de relieve]
Structure : \esp{Lo que + verbe + ser + SN focalisé}. Elle permet de mettre en relief l'\textbf{action}, le \textbf{complément d'objet} ou le \textbf{complément circonstanciel}, ce que la clivée simple ne fait pas toujours élégamment.
\begin{itemize}
\item \esp{Lo que busco es estabilidad.} « Ce que je cherche, c'est la stabilité. »
\item \esp{Lo que más me atrae de este puesto es el trabajo en equipo.} « Ce qui m'attire le plus dans ce poste, c'est le travail en équipe. »
\item \esp{Lo que necesitan ustedes es un técnico con experiencia.} « Ce qu'il vous faut, c'est un technicien expérimenté. »
\end{itemize}
\textbf{Variante inversée} : le bloc focalisé peut passer en tête : \esp{La experiencia es lo que más valoran.} « L'expérience, c'est ce qu'ils valorisent le plus. »
\end{definition}

\begin{attention}[\esp{Lo que} de relieve ou \esp{lo} de nominalisation ?]
Ne confonds pas :
\begin{itemize}
\item \esp{Lo que me interesa es el salario.} (relieve : « ce qui m'intéresse, c'est le salaire » — phrase clivée en deux blocs autour de \esp{es}, avec un verbe subordonné).
\item \esp{Lo importante es la experiencia.} (nominalisation : \esp{lo + adjectif} = « l'important », sans \esp{que} ni verbe subordonné).
\end{itemize}
Dans les deux cas, le français utilise « ce qui / ce que... c'est », mais seule la première structure focalise un verbe conjugué.
\end{attention}

\courssub{Règle 3 : La reprise pronominale (redundancia pronominal)}

L'espagnol adore « doubler » un complément par un pronom, surtout quand ce complément est \textbf{antéposé} (placé avant le verbe) ou désigne une \textbf{personne}.

\begin{propriete}[La réplication pronominale]
\begin{itemize}
\item \textbf{COI personne} : obligatoire avec \esp{a mí, a ti, a Juan...} : \esp{A mí me gusta este trabajo.} « Moi, ce travail me plaît. » ; \esp{A mi jefa le pareció bien mi carta.} « Ma cheffe a trouvé ma lettre bonne. »
\item \textbf{COD antéposé} : \esp{Ese puesto, yo lo quiero.} « Ce poste, moi je le veux. »
\item \textbf{Sujet redoublé} (insistance, contraste) : \esp{Juan, él lo hizo.} « Jean, c'est lui qui l'a fait. » ; \esp{Yo sí tengo experiencia.} « Moi, j'ai vraiment de l'expérience. »
\end{itemize}
\textbf{Valeurs} : insistance (\esp{A mí me da igual} « moi, ça m'est égal »), contraste (\esp{A ti te gusta, a mí no} « toi tu aimes, moi non »), clarification (\esp{Se lo di a ella} « je le lui ai donné, à elle », pour distinguer de « lui »).
\end{propriete}

\begin{exemplebox}[Traduire le « moi, je » français]
\esp{A esta empresa le interesa mi perfil.} « Cette entreprise, mon profil l'intéresse. »\par
\esp{A mí no me asustan los retos.} « Moi, les défis ne m'effraient pas. »\par
\esp{El inglés, lo domino perfectamente.} « L'anglais, je le maîtrise parfaitement. »
\end{exemplebox}

\courssub{Règle 4 : L'ordre des mots (orden inversado / anteposición)}

L'espagnol dispose d'un ordre des mots beaucoup plus souple que le français : déplacer un élément, c'est déjà le souligner.

\begin{propriete}[Deux mouvements emphatiques]
\begin{enumerate}
\item \textbf{Inversion sujet--verbe} (verbe + sujet), fréquente avec les verbes de mouvement, d'existence, d'apparition : \esp{Llegó el tren.} « Le train arriva » (annonce, nouvelle) vs \esp{El tren llegó.} (constat neutre). \esp{Sonó el teléfono.} « Le téléphone a sonné. »
\item \textbf{Antéposition du complément} (mise en tête), souvent avec reprise pronominale : \esp{Muchas puertas llamé.} « Bien des portes, j'en ai frappé. » ; \esp{¡Bien lo sabes!} « Tu le sais bien ! » vs \esp{Lo sabes bien} (neutre) ; en lettre soignée : \esp{El inglés lo domino perfectamente.} « L'anglais, je le maîtrise parfaitement. »
\end{enumerate}
\end{propriete}

\courssub{Règle 5 : Diminutifs et augmentatifs expressifs}

\begin{propriete}[Suffixes appréciatifs]
\begin{itemize}
\item \textbf{Diminutifs} \esp{-ito/-ita, -cito, -ecito, -illo} : petitesse, tendresse (\esp{cariño}) ou atténuation : \esp{una cartita} « une petite lettre », \esp{un momentito} « un petit moment », \esp{ahorita} « tout de suite » (Amérique).
\item \textbf{Augmentatifs} \esp{-ón/-ona, -azo, -ote} : grandeur ou intensité, souvent familiers, admiratifs ou péjoratifs : \esp{un hombretón} « un grand gaillard », \esp{una casona} « une grande et belle maison », \esp{un jefazo} « un sacré chef » (familier).
\item \textbf{Superlatif absolu} \esp{-ísimo} : \esp{grandísimo} « immense », \esp{muchísimas gracias} « merci infiniment ».
\end{itemize}
\end{propriete}

\begin{attention}[Registre : attention à la lettre de motivation]
Dans une \esp{carta de motivación}, bannis les suffixes familiers (\esp{jefazo}, \esp{jefecito}) : ils détruisent le sérieux du message. En revanche, certains superlatifs sont admis, voire élégants, dans les formules de politesse : \esp{muchísimas gracias por su atención} « merci infiniment de votre attention », \esp{Se lo agradecería infinitamente} « je vous en serais infiniment reconnaissant ». Règle d'or : l'emphase, en lettre formelle, passe par la \textbf{syntaxe} (clivée, relieve, ordre), non par les suffixes affectifs.
\end{attention}

\courssub{Comparaison systématique français / espagnol}

\begin{center}
\begin{tabularx}{\linewidth}{|X|X|X|}
\hline
\rowcolor{popblueL} Français & Español & Remarque \\
\hline
C'est ma grand-mère qui m'a appris. & \esp{Fue mi abuela quien me enseñó.} & \esp{quien} pour une personne sujet \\
\hline
C'est moi qui l'ai fait. & \esp{Soy yo quien lo he hecho.} & accord de \esp{ser} avec \esp{yo} \\
\hline
Ce que je cherche, c'est la stabilité. & \esp{Lo que busco es estabilidad.} & relieve : focalise l'action \\
\hline
Moi, ce travail me plaît. & \esp{A mí me gusta este trabajo.} & reprise pronominale obligatoire \\
\hline
L'anglais, je le maîtrise. & \esp{El inglés lo domino.} & antéposition + pronom \esp{lo} \\
\hline
C'est hier qu'il est arrivé. & \esp{Fue ayer cuando llegó.} & \esp{cuando} pour le temps \\
\hline
\end{tabularx}
\end{center}

\begin{popfigure}
\begin{tikzpicture}[font=\small]
\node[draw=popdark, fill=popblueL, rounded corners, align=center, minimum width=3.2cm] (titre) at (0,0) {\textbf{El relieve}\\\textbf{expresivo}};
\node[draw=popblue, fill=popblue!15, rounded corners, align=center, minimum width=4.4cm] (clivee) at (-5.2,2.6) {\textbf{Clivée}\\\esp{Es/Fue + SN + quien/que}\\\emph{identité d'une personne}};
\node[draw=poppurple, fill=poppurple!15, rounded corners, align=center, minimum width=4.4cm] (relieve) at (5.2,2.6) {\textbf{Ser de relieve}\\\esp{Lo que + V + es + SN}\\\emph{action, objet}};
\node[draw=popgreen, fill=popgreen!15, rounded corners, align=center, minimum width=4.4cm] (reprise) at (-5.2,-2.6) {\textbf{Reprise pronominale}\\\esp{A mí me..., lo...}\\\emph{insistance, contraste}};
\node[draw=poporange, fill=poporange!15, rounded corners, align=center, minimum width=4.4cm] (orden) at (5.2,-2.6) {\textbf{Ordre des mots}\\\esp{Llegó el tren / ¡Bien lo sabes!}\\\emph{nouvelle, style}};
\node[draw=popgold, fill=popgold!25, rounded corners, align=center, minimum width=4.4cm] (sufijos) at (0,-4.6) {\textbf{Diminutifs / augmentatifs}\\\esp{-ito, -ón, -ísimo}\\\emph{affection, intensité (registre !)}};
\draw[-{Stealth}, thick, popblue] (titre) -- (clivee);
\draw[-{Stealth}, thick, poppurple] (titre) -- (relieve);
\draw[-{Stealth}, thick, popgreen] (titre) -- (reprise);
\draw[-{Stealth}, thick, poporange] (titre) -- (orden);
\draw[-{Stealth}, thick, popgold] (titre) -- (sufijos);
\end{tikzpicture}
\legende{Carte mentale des cinq procédés de l'emphase espagnole.}
\end{popfigure}

\begin{aretenir}[Tableau synoptique des cinq procédés]
\begin{center}
\begin{tabularx}{\linewidth}{|l|X|X|X|}
\hline
\rowcolor{popblueL} Procédé & Structure type & Neutre $\rightarrow$ Emphatique & Valeur / registre \\
\hline
Clivée & \esp{Es/Fue + SN + quien/que} & \esp{Ella me contrató} $\rightarrow$ \esp{Fue ella quien me contrató} & identité ; tous registres \\
\hline
\esp{Ser} de relieve & \esp{Lo que + V + es + SN} & \esp{Busco estabilidad} $\rightarrow$ \esp{Lo que busco es estabilidad} & action/objet ; soutenu \\
\hline
Reprise & \esp{A mí me / lo, le...} & \esp{Me gusta} $\rightarrow$ \esp{A mí me gusta} & contraste ; courant \\
\hline
Ordre inversé & V + S / COD antéposé & \esp{El tren llegó} $\rightarrow$ \esp{Llegó el tren} & nouvelle ; littéraire \\
\hline
Suffixes & \esp{-ito, -ón, -ísimo} & \esp{carta} $\rightarrow$ \esp{cartita} & affectif ; familier (sauf \esp{-ísimo} soutenu) \\
\hline
\end{tabularx}
\end{center}
\end{aretenir}

\begin{methode}[Transformer une phrase neutre en phrase emphatique]
\begin{enumerate}
\item \textbf{Identifier} l'élément à souligner : une personne ? une action ? un contraste entre deux personnes ?
\item \textbf{Choisir} le procédé adapté : identité d'une personne $\rightarrow$ clivée ; action ou objet $\rightarrow$ \esp{ser} de relieve ; contraste (\esp{moi... toi...}) $\rightarrow$ reprise pronominale ; élégance narrative $\rightarrow$ ordre inversé ; nuance affective $\rightarrow$ suffixe.
\item \textbf{Construire} la phrase en appliquant la structure exacte.
\item \textbf{Vérifier} : accord de \esp{ser} (surtout avec \esp{yo}, \esp{nosotros}), pronom de reprise obligatoire, registre compatible avec la lettre formelle.
\end{enumerate}
\end{methode}

\begin{exoresolu}[Du neutre à l'emphatique]
\textbf{Énoncé.} Transforme les phrases suivantes en phrases emphatiques selon le procédé indiqué, puis traduis-les en français.\par
1. \esp{Mi formación me permitió conseguir este puesto.} (clivée)\par
2. \esp{Valoro la posibilidad de aprender.} (\esp{ser} de relieve)\par
3. \esp{Los idiomas no me asustan.} (reprise pronominale, contraste avec les autres candidats)
\tcblower
\textbf{Solution détaillée.}\par
1. Élément à focaliser : \esp{mi formación} (chose, sujet). Clivée avec \esp{lo que} : \esp{Fue mi formación lo que me permitió conseguir este puesto.} « C'est ma formation qui m'a permis d'obtenir ce poste. » (On emploie \esp{lo que} car l'élément focalisé est une chose ; \esp{ser} au passé \esp{fue} car le récit est au passé.)\par
2. Action à focaliser : \esp{valoro}. \esp{Ser} de relieve : \esp{Lo que valoro es la posibilidad de aprender.} « Ce que j'apprécie, c'est la possibilité d'apprendre. »\par
3. Contraste : on antépose \esp{a mí} et on garde le pronom \esp{me} : \esp{A mí no me asustan los idiomas.} « Moi, les langues ne m'effraient pas. » (Accord : \esp{asustan}, 3\textsuperscript{e} pers. pl., car le sujet est \esp{los idiomas}.)
\end{exoresolu}

\begin{exercice}[À toi de jouer]
1. Mets en relief l'élément souligné par une clivée : \esp{\textbf{Tú} debes firmar el contrato.}\par
2. Transforme avec \esp{ser} de relieve : \esp{Necesito un empleo estable.}\par
3. Traduis en espagnol : « C'est nous qui avons rédigé la lettre. » ; « Moi, l'expérience ne me manque pas. » ; « Ce qui compte, c'est la motivation. »
\end{exercice}

\begin{corrige}[Exercice : emphase]
1. \esp{Eres tú quien debe firmar el contrato.} (accord de \esp{ser} avec \esp{tú} ; \esp{quien} car personne sujet).\par
2. \esp{Lo que necesito es un empleo estable.}\par
3. \esp{Somos nosotros quienes hemos redactado la carta.} ; \esp{A mí no me falta experiencia.} ; \esp{Lo que cuenta es la motivación.}
\end{corrige}

\begin{attention}[Récapitulatif des pièges francophones]
\begin{itemize}
\item *\esp{Es yo quien...} $\rightarrow$ \esp{Soy yo quien...} (accord avec le pronom).
\item Après \esp{es} dans le relieve, l'espagnol omet souvent l'article là où le français le met : \esp{Lo que busco es estabilidad} (et non systématiquement \esp{la estabilidad}).
\item Ne double pas le sujet lexical par un pronom comme en français populaire : « Mon frère, il travaille » = \esp{Mi hermano trabaja} ou, avec emphase, \esp{Mi hermano, él sí que trabaja}.
\item Réserve \esp{-ito/-ón} à l'oral et aux contextes affectifs ; en lettre de motivation, préfère la syntaxe emphatique.
\end{itemize}
\end{attention}

\begin{savaistu}[L'emphase, une passion hispanique]
L'espagnol est l'une des langues romanes qui exploitent le plus l'ordre des mots à visée expressive : là où le français doit dire « c'est... qui », l'espagnol peut simplement déplacer. En Guinée équatoriale, seul pays hispanophone d'Afrique subsaharienne, la presse de Malabo use volontiers de la clivée dans ses titres (\esp{Fue el deporte lo que unió al país} « C'est le sport qui a uni le pays ») pour dramatiser l'information. Maîtriser ces tournures, c'est donc aussi mieux lire la presse hispanophone... et écrire une lettre de motivation qui sort du lot.
\end{savaistu}

\coursec{Atelier d'analyse : Commentaire guidé de l'emphase dans le texte}

Dans la section précédente, nous avons appris à \emph{reconnaître} et à \emph{former} les mécanismes de l'emphase : la phrase clivée (\esp{Es mi formación la que...}), le \esp{ser} de relieve, la répétition du pronom (\esp{A mí me apasiona...}), l'ordre inversé des mots et les diminutifs expressifs. Mais au baccalauréat, il ne suffit pas de savoir produire ces tournures : il faut aussi savoir les \emph{commenter}, c'est-à-dire expliquer pourquoi l'auteur les a choisies et quel effet elles produisent sur le lecteur. C'est précisément l'objet de cet atelier : passer du statut d'apprenant à celui d'analyste.

\courssub{Qu'est-ce qu'un commentaire de texte au Bac ?}

\begin{definition}[Le commentaire de texte (Bac)]
Le \cle{commentaire de texte} est un exercice écrit noté au baccalauréat. Il exige trois choses : (1) \cle{citer} exactement le passage étudié, (2) employer une \cle{terminologie grammaticale précise} (phrase clivée, \esp{ser} de relieve, inversion...), (3) proposer une \cle{analyse pragmatique} : quelle est l'intention de l'énonciateur, quel effet le procédé produit-il sur le destinataire ?
\end{definition}

La différence entre un bon et un mauvais commentaire tient en une phrase : le mauvais élève \emph{décrit}, le bon élève \emph{interprète}. Dire « il y a une phrase clivée » ne rapporte aucun point ; dire « la phrase clivée focalise l'attention du recruteur sur la formation du candidat afin de le présenter comme techniquement compétent » rapporte tous les points.

\courssub{Le texte support : rappel et relecture active}

Relisons le passage central de notre texte support, \emph{La carta que cambió mi destino}, dans lequel la narratrice, Aïcha, jeune diplômée de Yaoundé, raconte la rédaction de sa lettre de motivation pour un poste dans une entreprise d'exportation de cacao.

\begin{textebox}[La carta que cambió mi destino (fragmento)]
Cuando vi el anuncio en el periódico, supe que era mi oportunidad. Pero ¿cómo convencer a una empresa de Douala de que una joven de veintitrés años merecía el puesto? Pasé tres noches redactando mi carta.

No escribí «Tengo experiencia en comercio», como todo el mundo. Escribí: «Es mi experiencia en el mercado de Mvog-Mbi lo que me enseñó a negociar con clientes exigentes.» Tampoco escribí «Hablo tres idiomas». Preferí: «Lo que más valoro de mi formación es la capacidad de comunicarme en español, francés e inglés.»

A mi jefa actual le sorprendió mi carta. «Es la primera vez que leo una solicitud tan viva», me dijo. Y es que yo había escrito también: «A mí me apasiona el cacao camerunés, y es esa pasión la que quiero poner al servicio de su empresa.»

Dos semanas después, me llamaron. Fue esa carta la que cambió mi destino. Hoy, cuando los estudiantes me piden consejo, siempre repito lo mismo: no enumeréis vuestras cualidades como en una lista de la compra; hacedlas brillar. Porque una carta sin relieve es como un cacao sin aroma: existe, pero nadie lo recuerda.
\end{textebox}

\textbf{Glossaire :} \esp{el anuncio} (l'annonce), \esp{el puesto} (le poste), \esp{redactar} (rédiger), \esp{negociar} (négocier), \esp{exigente} (exigeant), \esp{la solicitud} (la candidature), \esp{poner al servicio de} (mettre au service de), \esp{el relieve} (le relief, l'emphase), \esp{el aroma} (l'arôme).

\textbf{Remarques linguistiques sur le texte.} Notez l'imparfait de subjonctif après \esp{¿cómo convencer... de que...?} : \esp{merecía} (subjonctif exigé par le doute implicite dans la question indirecte). Notez aussi \esp{Es la primera vez que leo} : après \esp{es la primera vez que}, l'espagnol utilise l'indicatif quand le fait est réel et accompli. Enfin, \esp{hacedlas brillar} est l'impératif affirmatif de \esp{vosotros} (\esp{haced} + pronom \esp{las} soudé, avec accent maintenu).

\courssub{La méthode de commentaire en trois étapes}

\begin{methode}[Commenter l'emphase au Bac : 3 étapes obligatoires]
\begin{enumerate}
\item \textbf{CITER} la phrase exacte, entre guillemets, sans la déformer : \esp{« Es mi experiencia en el mercado de Mvog-Mbi lo que me enseñó a negociar... »}
\item \textbf{NOMMER} le procédé avec la terminologie exacte : phrase clivée sujet (\esp{Es... lo que}), \esp{ser} de relieve, répétition du pronom (\esp{A mí me...}), ordre inversé, etc.
\item \textbf{EXPLIQUER} l'intention communicative : « Le candidat met l'accent sur X pour montrer Y » — persuader, rassurer, se différencier, créer de la proximité.
\end{enumerate}
\end{methode}

\begin{popfigure}
\begin{tikzpicture}[
  node distance=0.55cm,
  etape/.style={rectangle, rounded corners=4pt, draw=popblue, fill=popblueL, text width=3.1cm, align=center, minimum height=1.1cm, font=\small},
  fleche/.style={-{Stealth[length=3mm]}, thick, poporange}
]
\node[etape, fill=popcream, draw=popgold, align=center] (t){\textbf{Texte}\\ lire et relire};
\node[etape, right=of t, align=center] (r){\textbf{1. Repérage}\\ surligner les tournures};
\node[etape, right=of r, align=center] (c){\textbf{2. Classification}\\ tableau synoptique};
\node[etape, right=of c, align=center] (i){\textbf{3. Interprétation}\\ Pourquoi ici ?};
\node[etape, right=of i, fill=poppurpleL, draw=poppurple, align=center] (red){\textbf{4. Rédaction}\\ commentaire synthétique};
\draw[fleche] (t) -- (r);
\draw[fleche] (r) -- (c);
\draw[fleche] (c) -- (i);
\draw[fleche] (i) -- (red);
\end{tikzpicture}
\legende{Le flux de la méthode de commentaire : du texte à la rédaction.}
\end{popfigure}

\courssub{Exercice ciblé : repérer et classer les 5 occurrences d'emphase}

\begin{exoresolu}[Repérage et classification dans le texte support]
\textbf{Énoncé.} Retrouvez dans le fragment ci-dessus les 5 occurrences d'emphase et classez-les dans le tableau synoptique de la section 3 (procédé, citation, élément mis en relief).
\tcblower
\textbf{Solution détaillée.}
\begin{enumerate}
\item \esp{« Es mi experiencia en el mercado de Mvog-Mbi lo que me enseñó a negociar... »} — \textbf{phrase clivée} (\esp{Es... lo que}) : relief sur \esp{mi experiencia}.
\item \esp{« Lo que más valoro de mi formación es la capacidad de comunicarme... »} — \textbf{pseudo-clivée} (\esp{Lo que... es...}) : relief sur \esp{la capacidad de comunicarme}.
\item \esp{« A mí me apasiona el cacao camerunés... »} — \textbf{répétition du pronom} (\esp{A mí me...}) : relief sur la personne de l'énonciatrice.
\item \esp{« ...y es esa pasión la que quiero poner al servicio de su empresa. »} — \textbf{phrase clivée} (\esp{es... la que}) : relief sur \esp{esa pasión}.
\item \esp{« Fue esa carta la que cambió mi destino. »} — \textbf{phrase clivée au passé} (\esp{Fue... la que}) : relief sur \esp{esa carta}, qui clôt le récit en écho au titre.
\end{enumerate}
\end{exoresolu}

\begin{center}
\begin{tabularx}{\linewidth}{|l|X|X|X|}
\hline
\rowcolor{popblueL} \textbf{N°} & \textbf{Citation} & \textbf{Procédé} & \textbf{Effet de sens} \\
\hline
1 & \esp{Es mi experiencia... lo que me enseñó...} & Clivée & Prouve la compétence concrète \\
\hline
2 & \esp{Lo que más valoro... es la capacidad...} & Pseudo-clivée & Hiérarchise les qualités \\
\hline
3 & \esp{A mí me apasiona el cacao...} & Pronom redoublé & Affirme la motivation personnelle \\
\hline
4 & \esp{Es esa pasión la que quiero poner...} & Clivée & Transforme l'émotion en offre de service \\
\hline
5 & \esp{Fue esa carta la que cambió mi destino} & Clivée (passé) & Conclut et valorise l'écriture soignée \\
\hline
\end{tabularx}
\end{center}

\courssub{Du repérage à l'interprétation : exemples de commentaires réussis}

\begin{exemplebox}[Modèles de commentaires complets]
\textbf{Texte original :} \esp{Es mi formación la que me permite gestionar los pedidos internacionales.} (« C'est ma formation qui me permet de gérer les commandes internationales. »)\\
\textbf{Commentaire :} « Phrase clivée sujet (\esp{Es... la que}). Le candidat focalise l'attention du recruteur sur \esp{formación} pour prouver sa compétence technique face aux exigences du poste : il ne dit pas "je suis formé", il fait de sa formation l'héroïne de la phrase. »

\medskip
\textbf{Texte original :} \esp{A mí me apasiona el cacao camerunés.} (« Le cacao camerounais, c'est cela qui me passionne, moi. »)\\
\textbf{Commentaire :} « Redoublement pronominal (\esp{A mí me...}). La candidate se met personnellement en avant pour créer une proximité émotionnelle avec le lecteur et se différencier des candidatures anonymes : sa motivation n'est pas un cliché, c'est un engagement. »
\end{exemplebox}

\begin{attention}[L'erreur d'analyse qui coûte des points]
Ne dites jamais « c'est pour insister » sans préciser \textbf{SUR QUOI} porte l'insistance et \textbf{POURQUOI}. La formulation « l'auteur insiste » est vide. L'\cle{intention communicative} — persuader, rassurer, différencier, émouvoir — est \textbf{obligatoire} pour obtenir les points d'analyse. Formule magique : « Le procédé X met en relief Y \emph{afin de} Z. »
\end{attention}

\courssub{Rédaction d'un micro-commentaire structuré (50 mots)}

Consigne de production : rédigez un paragraphe d'environ 50 mots commençant par « Dans ce texte, l'emphase sert à... » en choisissant une fonction dominante : valoriser les compétences, marquer la motivation ou créer la proximité.

\begin{exemplebox}[Corrigé type de micro-commentaire (environ 50 mots)]
« Dans ce texte, l'emphase sert à \textbf{marquer la motivation} de la candidate. La phrase clivée \esp{« es esa pasión la que quiero poner al servicio de su empresa »} transforme un simple sentiment en promesse professionnelle, tandis que le redoublement \esp{« A mí me apasiona »} personnalise l'engagement. L'emphase fait ainsi passer la lettre du registre administratif au registre de la conviction. »
\end{exemplebox}

\courssub{Comparaison : le texte « à plat » et la perte d'impact}

Réécrivons le deuxième paragraphe du texte \emph{sans aucune emphase}, dans un style neutre et plat :

\begin{center}
\begin{tabularx}{\linewidth}{|X|X|}
\hline
\rowcolor{popblueL} \textbf{Version emphatique (originale)} & \textbf{Version plate (réécrite)} \\
\hline
\esp{Es mi experiencia en el mercado de Mvog-Mbi lo que me enseñó a negociar con clientes exigentes.} & \esp{Aprendí a negociar con clientes exigentes gracias a mi experiencia en el mercado de Mvog-Mbi.} \\
\hline
\esp{Lo que más valoro de mi formación es la capacidad de comunicarme en tres idiomas.} & \esp{Mi formación me permite comunicarme en tres idiomas.} \\
\hline
\esp{A mí me apasiona el cacao camerunés.} & \esp{Me gusta el cacao camerunés.} \\
\hline
\end{tabularx}
\end{center}

\textbf{Analyse de la perte.} La version plate est grammaticalement correcte mais rhétoriquement faible : (1) l'information essentielle (\esp{mi experiencia}) n'occupe plus la place d'honneur en tête de phrase ; (2) \esp{me gusta} est banal là où \esp{me apasiona} avec redoublement affirme une vocation ; (3) le lecteur-recruteur ne retient rien, car toutes les phrases se ressemblent. L'emphase est donc un \cle{outil de persuasion} : elle hiérarchise l'information et guide l'œil du destinataire.

\begin{savaistu}[L'ordre inversé dans les lettres administratives modernes]
Dans les lettres administratives et professionnelles espagnoles contemporaines, l'emphase par \cle{ordre inversé} (verbe avant sujet) est très prisée avec les verbes d'action à la première personne : \esp{Aporto experiencia} (« j'apporte de l'expérience »), \esp{Garantizo resultados} (« je garantis des résultats »), \esp{Domino idiomas} (« je maîtrise les langues »). Placer le verbe en tête donne une image dynamique et décidée du candidat — très appréciée des recruteurs, y compris en Guinée équatoriale, seul pays hispanophone d'Afrique.
\end{savaistu}

\courssub{Entraînement Bac}

\begin{exercice}[Question type Bac]
Dans le texte \emph{La carta que cambió mi destino}, expliquez l'effet produit par la structure \esp{« Es mi experiencia lo que... »} (ligne 5). Votre réponse (4 à 6 lignes) devra citer la phrase, nommer le procédé et analyser l'intention de la candidate.
\end{exercice}

\begin{corrige}[Exercice d'entraînement Bac]
La structure \esp{« Es mi experiencia en el mercado de Mvog-Mbi lo que me enseñó a negociar »} est une \textbf{phrase clivée} (\esp{Es... lo que}). Elle place l'expérience concrète de la candidate en position focale, avant même le verbe. Intention : persuader le recruteur que la candidate possède une compétence \emph{réelle et vécue} (le marché de Mvog-Mbi, non pas un simple diplôme), et ainsi se différencier des autres candidats. L'emphase transforme une information banale en argument décisif.
\end{corrige}

\begin{exercice}[Pour aller plus loin]
Reprenez la phrase plate \esp{« Tengo cinco años de experiencia y hablo español »} et réécrivez-la avec trois procédés d'emphase différents (clivée, pseudo-clivée, pronom redoublé), puis commentez en deux lignes l'effet de chacune.
\end{exercice}

\begin{corrige}[Exercice « Pour aller plus loin » — éléments de réponse]
\begin{enumerate}
\item \textbf{Clivée :} \esp{Son cinco años de experiencia los que avalan mi candidatura.} (« Ce sont cinq années d'expérience qui soutiennent ma candidature. ») — l'ancienneté devient l'argument central.
\item \textbf{Pseudo-clivée :} \esp{Lo que distingue mi perfil es que hablo español con fluidez.} (« Ce qui distingue mon profil, c'est que je parle espagnol couramment. ») — la langue est présentée comme un atout unique.
\item \textbf{Pronom redoublé :} \esp{A mí el español me apasiona desde el instituto.} (« Moi, l'espagnol me passionne depuis le lycée. ») — la motivation personnelle est mise en avant.
\end{enumerate}
\end{corrige}

\coursec{Expression écrite : Rédaction de la Carta de Motivación (120 mots)}

Dans la section précédente, nous avons disséqué les procédés d'emphase dans la lettre de Kwame : la tournure clivée \esp{Fue mi carta la que abrió la puerta}, le \esp{ser} de relieve, l'ordre inversé. Il est temps maintenant de passer de l'\cle{analyse} à la \cle{production} : à votre tour de rédiger une lettre de motivation qui, en 120 mots exactement, frappe l'esprit du recruteur. C'est l'exercice d'expression écrite le plus fréquent au baccalauréat dans le module « monde du travail » : autant le maîtriser à la perfection.

\courssub{La structure canonique : le squelette de la lettre}

Avant d'écrire une seule phrase, il faut connaître la \cle{géographie} de la lettre formelle espagnole. Elle diffère légèrement de la lettre française : en espagnol, les coordonnées de l'\cle{expéditeur} se placent en haut à \emph{gauche}, le lieu et la date en haut à \emph{droite} (et non l'inverse), et le destinataire revient à gauche, sous la date.

\begin{popfigure}
\begin{tikzpicture}[
  bloque/.style={draw=popblue, fill=popblueL, rounded corners=2pt, align=left, font=\small, inner sep=4pt},
  etiqueta/.style={font=\scriptsize\itshape, text=popmuted}
]
\node[bloque, anchor=north west, text width=4.6cm] (rem) at (0,0)
  {\textbf{Kwame Ngo Bassa}\\ Calle del Mercado, 12\\ Yaoundé (Camerún)\\ kwame@correo.cm};
\node[bloque, anchor=north east, text width=4.2cm] (fecha) at (10.5,0)
  {Yaoundé, a 15 de marzo de 2025};
\node[bloque, anchor=north west, text width=5.2cm] (dest) at (0,-2.2)
  {\textbf{Sr. Director de Recursos Humanos}\\ Empresa CacaoCam S.A.\\ Douala (Camerún)};
\node[bloque, anchor=north west, text width=10.5cm] (asunto) at (0,-4.1)
  {\textbf{Asunto:} Candidatura al puesto de técnico comercial};
\node[bloque, anchor=north west, text width=10.5cm] (saludo) at (0,-5.3)
  {Estimado Sr. Director:};
\node[bloque, anchor=north west, text width=10.5cm, minimum height=2.6cm, align=left] (cuerpo) at (0,-6.3)
  {\textbf{P1} Presentación + motivo de la carta\\[2pt]
   \textbf{P2} Formación y experiencia (emphase)\\[2pt]
   \textbf{P3} Motivación y disponibilidad (emphase)};
\node[bloque, anchor=north west, text width=10.5cm] (desp) at (0,-9.6)
  {Quedo a la espera de sus noticias, le saluda atentamente,};
\node[bloque, anchor=north west, text width=4cm] (firma) at (0,-10.7)
  {\emph{(firma manuscrita)}\\ \textbf{Kwame Ngo Bassa}};
\node[etiqueta, anchor=west] at (5.2,-11.6) {El nombre impreso va siempre bajo la firma};
\draw[popline, dashed] (-0.3,0.6) rectangle (10.8,-12.1);
\end{tikzpicture}
\legende{Gabarit visuel de la carta de motivación : chaque bloc a sa place exacte.}
\end{popfigure}

\begin{aretenir}[Les cinq blocs obligatoires]
\begin{enumerate}
\item \textbf{Encabezado} : coordonnées de l'expéditeur (gauche), lieu et date (droite), coordonnées du destinataire (gauche).
\item \textbf{Asunto} : l'objet, en une ligne, sans verbe conjugué : \esp{Asunto: Candidatura al puesto de...}
\item \textbf{Saludo} (appel) : suivi de \textbf{deux-points}, jamais de virgule.
\item \textbf{Cuerpo} : 3 ou 4 paragraphes courts (présentation, formation/expérience, motivation).
\item \textbf{Despedida + firma} : formule de politesse, virgule, signature, nom imprimé.
\end{enumerate}
\end{aretenir}

\begin{attention}[La date espagnole : un piège classique]
En espagnol, la date s'écrit \esp{Yaoundé, a 15 de marzo de 2025} : préposition \esp{a} + article facultatif, mois en \textbf{minuscule} (\esp{marzo}, pas « Marzo »), pas d'ordinal (\esp{15}, pas « 15.º » sauf pour le premier du mois : \esp{a 1.º de mayo}). Ne calquez jamais le français « le 15 mars ».
\end{attention}

\courssub{Les formules de politesse codifiées}

La lettre formelle espagnole est un genre \cle{ritualisé} : les formules d'ouverture et de clôture sont figées. Les connaître par cœur, c'est gagner des points de « compétence pragmatique » sans effort.

\begin{propriete}[Formules « passe-partout » pour le Bac]
\textbf{Ouverture (appel) :}
\begin{itemize}
\item Destinataire inconnu : \esp{Muy señores míos:} (toujours suivi de deux-points).
\item Destinataire connu : \esp{Estimado Sr. Director:} / \esp{Estimada Sra. García:}
\end{itemize}
\textbf{Fermeture (politesse finale) :}
\begin{itemize}
\item \esp{Quedo a la espera de sus noticias, le saluda atentamente,} + Nom Prénom.
\item Variantes admises : \esp{Atentamente,} / \esp{Cordialmente,} / \esp{Le saluda atentamente,}
\end{itemize}
\end{propriete}

\begin{center}
\begin{tabularx}{\linewidth}{|X|X|X|}
\hline
\rowcolor{popblueL} Français & Español & Remarque \\
\hline
Monsieur, Madame & Muy señores míos: & destinataire inconnu \\
\hline
Monsieur le Directeur & Estimado Sr. Director: & deux-points obligatoires \\
\hline
Je vous prie d'agréer... & Le saluda atentamente & formule courte, standard \\
\hline
Dans l'attente de votre réponse & Quedo a la espera de su respuesta & \esp{quedar} à la 1\textsuperscript{re} pers. \\
\hline
Veuillez trouver ci-joint mon CV & Le adjunto mi currículum vítae & \esp{adjuntar} + \esp{le} (courtoisie) \\
\hline
Je reste à votre disposition & Quedo a su entera disposición & ton très formel \\
\hline
\end{tabularx}
\end{center}

\begin{savaistu}[Le « leísmo de cortesía »]
En Espagne péninsulaire, la norme formelle emploie \esp{le} (et non \esp{lo}) pour le complément d'objet direct masculin désignant une personne : \esp{Le escribo} « je vous écris », \esp{Le adjunto mi CV} « je vous joins mon CV ». C'est le fameux \emph{leísmo de cortesía}, accepté par la Real Academia. Au Bac, pour un destinataire masculin, utilisez \esp{le} : c'est la marque du registre soutenu. En revanche, pour une femme, on dit \esp{la escribo}.
\end{savaistu}

\courssub{Atelier d'injection d'emphase : de la phrase plate à la phrase impactante}

Une lettre de 120 mots ne peut pas se permettre d'être banale. Chaque paragraphe doit contenir au moins un procédé d'emphase étudié dans la section 3. Voici cinq transformations modèles, à reproduire dans vos propres rédactions.

\begin{exemplebox}[Cinq transformations à imiter]
\begin{enumerate}
\item \textbf{Tournure clivée} (pour le diplôme) :\par
Phrase plate : \esp{Tengo un máster en marketing.} « J'ai un master en marketing. »\par
Emphase : \esp{Es mi máster en marketing lo que avala mi perfil.} « C'est mon master en marketing qui garantit mon profil. »
\item \textbf{Ser de relieve} (pour la motivation) :\par
Phrase plate : \esp{Quiero trabajar con ustedes.} « Je veux travailler avec vous. »\par
Emphase : \esp{Lo que deseo es integrar su equipo.} « Ce que je désire, c'est intégrer votre équipe. »
\item \textbf{Ordre inversé} (pour les langues) :\par
Phrase plate : \esp{Hablo francés, inglés y español.}\par
Emphase : \esp{Tres lenguas domino: el francés, el inglés y el español.} « Trois langues, je les maîtrise... »
\item \textbf{Répétition du pronom} (pour l'expérience) :\par
Phrase plate : \esp{Conozco bien el mercado camerunés.}\par
Emphase : \esp{El mercado camerunés lo conozco desde hace diez años.} « Le marché camerounais, je le connais depuis dix ans. »
\item \textbf{Augmentatif expressif} (pour l'engagement, avec modération) :\par
Phrase plate : \esp{Tengo muchas ganas de empezar.}\par
Emphase : \esp{Traigo unas ganas enormes de empezar.} « J'apporte une envie énorme de commencer. »
\end{enumerate}
\end{exemplebox}

\begin{attention}[Dosage de l'emphase et morphologie]
Trois emphases dans une lettre de 120 mots, c'est l'idéal ; cinq, c'est une caricature. L'emphase fonctionne comme le piment dans un ndolé : une pincée sublime, une cuillerée gâche tout. Placez-en une par paragraphe du corps, jamais dans l'appel ni dans la formule finale, qui doivent rester neutres et codifiés. \textbf{Rappel morphologique} : les augmentatifs (\esp{-ón, -azo, -ote}) et diminutifs (\esp{-ito, -illo}) marquent l'intensité ou l'affectivité (\esp{ganazas, trabajadorazo, ratito}), mais s'utilisent avec une extrême prudence dans le formel ; privilégiez l'intensité lexicale (\esp{enormes, tremendas, total}).
\end{attention}

\courssub{Rédaction guidée pas à pas (Process Writing)}

On ne rédige pas une lettre de Bac d'un seul jet : on la \cle{construit} en cinq étapes.

\begin{methode}[Les cinq étapes de la rédaction]
\begin{enumerate}
\item \textbf{Lluvia de ideas} (brainstorming, 3 min) : remplir une mini-grille SWOT personnelle — \esp{Fortalezas} (points forts : diplômes, langues), \esp{Debilidades} (à taire ou à retourner), \esp{Oportunidades} (ce que l'entreprise offre), \esp{Amenazas} (concurrents : ce qui vous distingue).
\item \textbf{Plan détaillé} (2 min) : noter 3 mots-clés par paragraphe. P1 : \esp{anuncio, puesto, interés}. P2 : \esp{máster, experiencia, idiomas}. P3 : \esp{motivación, disponibilidad, entrevista}.
\item \textbf{Primer borrador} (premier jet, 10 min) : écrire librement, cible 130 mots, sans compter.
\item \textbf{Tallado} (taillage, 5 min) : supprimer les redondances, injecter les 3 emphases, redescendre à 115-125 mots.
\item \textbf{Revisión} (relecture, 5 min) : passer la check-list des 10 points (voir ci-dessous).
\end{enumerate}
\end{methode}

\begin{attention}[Le comptage des mots : convention Cameroun]
Le baccalauréat est strict sur le volume. Selon les académies, les formules d'appel et de politesse finale ne comptent pas toujours, mais le \textbf{corps} compte toujours. \textbf{Convention retenue au Cameroun} : comptez tout le texte compris entre \esp{Estimado...} et \esp{Atentamente} exclus. Visez un corps de \textbf{115 à 125 mots} : en dessous de 110, vous êtes hors sujet quantitatif ; au-dessus de 130, vous êtes pénalisé. Entraînez-vous à compter vite : les mots composés avec trait d'union comptent pour un, les contractions \esp{al}, \esp{del} pour un.
\end{attention}

\begin{textebox}[Check-list de relectura final : los 10 puntos]
1. ¿He contado las palabras del cuerpo? (115-125)\\
2. ¿La fórmula de saludo lleva dos puntos?\\
3. ¿La despedida es formal y completa?\\
4. ¿Hay al menos tres énfasis (clivada, relieve, orden inverso)?\\
5. ¿Los subjuntivos de cortesía son correctos? (\esp{le agradecería que...})\\
6. ¿Las clivadas concuerdan? (\esp{Es mi máster lo que...} / \esp{Son mis estudios los que...})\\
7. ¿He usado léxico del empleo? (\esp{puesto, candidatura, currículum})\\
8. ¿Hay tres párrafos separados y visibles?\\
9. ¿La fecha sigue el modelo \esp{a 15 de marzo de 2025}?\\
10. ¿La firma va seguida del nombre impreso?
\end{textebox}

\emph{Glossaire} : \esp{despedida} = formule de clôture ; \esp{clivada} = tournure clivée ; \esp{concordar} = s'accorder ; \esp{puesto} = poste ; \esp{firma} = signature.

\courssub{Exercice résolu : du brouillon à la lettre parfaite}

\begin{exoresolu}[Transformar y redactar]
\textbf{Énoncé.} Voici le premier jet (trop long et plat) de Kwame : \esp{Me dirijo a ustedes porque vi su anuncio. Tengo un máster en marketing y hablo tres idiomas. Trabajé dos años en una empresa de cacao. Quiero mucho este puesto y puedo empezar ya. Espero su respuesta.} (1) Comptez les mots. (2) Injectez trois emphases (clivée, \esp{ser} de relieve, ordre inversé). (3) Encadrez avec les formules codifiées.
\tcblower
\textbf{Solution détaillée.}
(1) Le brouillon compte 36 mots : il faut développer avec du lexique de l'emploi (\esp{candidatura, trayectoria, disponibilidad inmediata}) pour atteindre 115-125 mots dans la version complète.
(2) Transformations :
\begin{itemize}
\item Clivée : \esp{Fue su anuncio lo que me impulsó a escribirles.} « C'est votre annonce qui m'a poussé à vous écrire. »
\item \esp{Ser} de relieve : \esp{Lo que más me atrae es la dimensión internacional del puesto.} « Ce qui m'attire le plus, c'est la dimension internationale du poste. »
\item Ordre inversé : \esp{Dos años pasé en una empresa cacaotera de Douala.} « Deux ans, je les ai passés dans une entreprise cacaoyère de Douala. »
\end{itemize}
(3) Encadrement : ouverture \esp{Estimado Sr. Director:} ; clôture \esp{Quedo a la espera de sus noticias, le saluda atentamente,} + \esp{Kwame Ngo Bassa}. Le corps ainsi enrichi atteint environ 120 mots, avec trois emphases réparties une par paragraphe : la lettre est conforme aux attentes du Bac.
\end{exoresolu}

\courssub{Critères d'évaluation au baccalauréat}

Le correcteur lit votre lettre avec une grille. Autant la connaître avant lui.

\begin{center}
\begin{tabularx}{\linewidth}{|l|X|X|}
\hline
\rowcolor{popblueL} Critère & Ce que le correcteur attend & Erreur qui coûte cher \\
\hline
Respect des consignes & 115-125 mots, 5 blocs présents & lettre de 80 ou 200 mots \\
\hline
Correction linguistique & accords, subjonctif, temps B2 & \esp{*yo quiero que usted me da} \\
\hline
Richesse lexicale & lexique de l'emploi + 3 emphases & vocabulaire d'école primaire \\
\hline
Cohérence & paragraphes liés, progression logique & idées en vrac \\
\hline
Pragmatique & ton formel, \esp{usted}, \esp{leísmo} & tutoyer le directeur \\
\hline
\end{tabularx}
\end{center}

\begin{exercice}[À votre tour]
Rédigez une lettre de motivation de 120 mots (corps) en réponse à cette annonce : \esp{Empresa AgroYaoundé busca técnico comercial bilingüe. Enviar carta de motivación y CV al Sr. Director de Recursos Humanos.} Appliquez les cinq étapes du Process Writing, injectez trois emphases de procédés différents, puis auto-évaluez-vous avec la check-list des 10 points et la grille ci-dessus.
\end{exercice}

\begin{corrige}[Exercice]
\textbf{Éléments de corrigé (modèle de corps, 114 mots).}\par
\esp{Estimado Sr. Director:}\par
\esp{Fue su anuncio en el periódico El Eco de Yaoundé lo que despertó mi vivo interés por el puesto de técnico comercial bilingüe. Permítame presentarle mi candidatura.}\par
\esp{Es mi formación bilingüe y mi experiencia práctica lo que mejor avala mi perfil: el francés y el español los domino a la perfección desde el bachillerato, y dos años completos pasé en el departamento de ventas de una importante empresa cacaotera de Douala, gestionando carteras de clientes mayoristas.}\par
\esp{Lo que más deseo es integrar un equipo dinámico y profesional como el suyo. Traigo, además, una disponibilidad inmediata, carné de conducir y unas ganas tremendas de contribuir activamente al desarrollo comercial de AgroYaoundé en la zona CEMAC.}\par
\esp{Quedo a la espera de sus noticias, le saluda atentamente,}\par
\esp{Mballa Etoundi, Sandra}\par
Points à vérifier dans votre copie : trois emphases repérables (clivée \esp{Fue su anuncio lo que...}, clivée \esp{Es mi formación... lo que...}, ordre inversé \esp{dos años pasé}, répétition pronominale \esp{los domino}, \esp{ser} de relieve \esp{Lo que más deseo...}), registre \esp{usted} constant, \esp{leísmo} de courtoisie (\esp{le}, \esp{presentarle}), date et blocs conformes au gabarit.
\end{corrige}

\coursec{Synthèse, Entraînement Bac \& Auto-évaluation}

Vous venez de rédiger votre première \esp{carta de motivación} complète : félicitations. Mais au baccalauréat, on ne vous demandera pas seulement de produire : on vous demandera de \cle{reconnaître}, de \cle{traduire} et de \cle{réutiliser} les procédés d'emphase dans un temps limité. Cette dernière section est votre sas de révision : une fiche mémo visuelle, des exercices de thème et de version calibrés, un sujet type Bac complet chronométré, et une grille d'auto-évaluation pour savoir exactement où vous en êtes.

\courssub{Fiche mémo : tout leçon E18 en un coup d'œil}

\begin{aretenir}[Les 5 procédés d'emphase à maîtriser]
\begin{enumerate}
\item \textbf{La clivée} : \esp{Fue + élément + quien/que/donde/cuando + reste}. \esp{Fue mi tío quien me ayudó.} « C'est mon oncle qui m'a aidé. »
\item \textbf{Le ser de relieve (pseudo-clivée)} : \esp{Lo que + verbe + es/son + élément}. \esp{Lo que quiero es trabajar.} « Ce que je veux, c'est travailler. »
\item \textbf{La reprise pronominale (redoublement clitique)} : \esp{A Juan, le ofrecieron el puesto.} Le pronom d'objet (\esp{le, la, lo, les}) reprend le nom mis en tête (antéposé).
\item \textbf{L'ordre marqué (topicalisation)} : \esp{El inglés, lo domino.} « L'anglais, je le maîtrise. » Le complément passe en tête et est repris par le pronom.
\item \textbf{Les suffixes évaluatifs} : diminutif \esp{-ito/-cito} (\esp{un momentito}), augmentatif \esp{-ón/-azo} (\esp{un trabajón}), superlatif \esp{-ísimo} (\esp{importantísimo}).
\end{enumerate}
\end{aretenir}

\begin{center}
\begin{tabularx}{\linewidth}{|l|X|X|}
\hline
\rowcolor{popblueL} \textbf{Procédé} & \textbf{Espagnol} & \textbf{Français} \\
\hline
Clivée & \esp{Fue ayer cuando llegó.} & C'est hier qu'il est arrivé. \\
\hline
Ser de relieve & \esp{Lo que necesito es práctica.} & Ce dont j'ai besoin, c'est de la pratique. \\
\hline
Reprise pronominale & \esp{A mí me gusta el reto.} & Moi, j'aime le défi. \\
\hline
Ordre inversé & \esp{El inglés, lo domino.} & L'anglais, je le maîtrise. \\
\hline
Augmentatif (suffixe) & \esp{una oportunidad grandísima} & une opportunité énorme \\
\hline
\end{tabularx}
\end{center}

\begin{aretenir}[Structure de la carta de motivación (120 palabras)]
\begin{enumerate}
\item \textbf{En-tête} : lieu et date (\esp{Yaoundé, a 15 de mayo de 2025}), destinataire.
\item \textbf{Appel} : \esp{Estimado señor director:} / \esp{Muy señor mío:} (deux-points obligatoires).
\item \textbf{Corps} : présentation (\esp{Me dirijo a usted para solicitar...}), motivations (\esp{Lo que más me atrae es...}), compétences (\esp{Domino el francés y el español.}).
\item \textbf{Clôture} : \esp{Quedo a su disposición... Le saluda atentamente,} + signature.
\end{enumerate}
\end{aretenir}

\begin{aretenir}[Top 10 du lexique de l'emploi]
\begin{center}
\begin{tabularx}{\linewidth}{|X|X|}
\hline
\rowcolor{popgreenL} \textbf{Español} & \textbf{Français} \\
\hline
\esp{el puesto de trabajo} & le poste de travail \\
\hline
\esp{la solicitud / solicitar} & la candidature / postuler \\
\hline
\esp{el currículum vítae} & le CV \\
\hline
\esp{la entrevista de trabajo} & l'entretien d'embauche \\
\hline
\esp{el contrato (fijo / temporal)} & le contrat (CDI / CDD) \\
\hline
\esp{el salario / el sueldo} & le salaire \\
\hline
\esp{la experiencia laboral} & l'expérience professionnelle \\
\hline
\esp{las prácticas / el pasante} & le stage / le stagiaire \\
\hline
\esp{contratar / despedir} & embaucher / licencier \\
\hline
\esp{la carta de motivación} & la lettre de motivation \\
\hline
\end{tabularx}
\end{center}
\end{aretenir}

\begin{attention}[Top 5 des pièges à éviter]
\begin{enumerate}
\item Ne traduisez jamais « c'est... qui » par \esp{es... que} : la clivée espagnole exige \esp{Fue... quien/que}. \esp{Fue mi madre quien me enseñó} (et non \esp{Es mi madre que...}).
\item Après \esp{Fue} + personne, employez \esp{quien} (singulier) ou \esp{quienes} (pluriel), jamais \esp{que} pour un sujet humain.
\item La reprise pronominale exige le pronom : \esp{A Juan le ofrecieron el puesto} (sans \esp{le}, la phrase est fautive).
\item Ne mélangez pas \esp{tú} et \esp{usted} dans la même lettre : la lettre de motivation est toujours en \esp{usted}.
\item Attention aux faux amis : \esp{actualmente} = « actuellement » (pas « en fait ») ; \esp{asistir a} = « être présent à » ; \esp{una oportunidad} (un seul \emph{p}, pas d'accent) et non \esp{una oportunité}.
\end{enumerate}
\end{attention}

\courssub{Entraînement de thème (français $\rightarrow$ espagnol)}

Consigne générale : transformez chaque phrase neutre en phrase emphatique en utilisant le procédé imposé entre parenthèses.

\begin{exoresolu}[Thème guidé : phrases 1 et 2]
\textbf{Énoncé 1.} « C'est mon stage au Mexique qui m'a formé. » (Clivée)

\tcblower

\textbf{Solution.} L'élément mis en relief est \emph{mon stage au Mexique}, une chose : on utilise donc \esp{Fue + élément + que} (et non \esp{quien}, réservé aux personnes). Le verbe \esp{formar} se met au prétérit : \esp{formó}. Accord du relatif avec l'antécédent : \esp{la que} (féminin \esp{etapa}).

\esp{Fue mi etapa en México la que me formó.}

Variante acceptée (neutre) : \esp{Fue mi etapa en México lo que me formó.}
\emph{Attention :} \esp{stage} en espagnol se dit \esp{las prácticas} ou \esp{la estancia/etapa} ; \esp{el estadio} est un faux ami (stade de football).

\bigskip
\textbf{Énoncé 2.} « Ce que je veux, c'est apprendre. » (Ser de relieve)

\tcblower

\textbf{Solution.} Structure \esp{Lo que + verbe + es + infinitif} :

\esp{Lo que quiero es aprender.}

Piège : ne dites pas \esp{El que quiero...} ; le pronom neutre \esp{lo que} est obligatoire devant une idée générale.
\end{exoresolu}

\begin{exercice}[Thème : phrases 3, 4 et 5]
Traduisez en espagnol avec le procédé imposé :
\begin{enumerate}
\item « Moi, je parle couramment espagnol. » (Emphase sur le sujet / pronom tonique)
\item « L'anglais, je le maîtrise. » (Ordre inversé / topicalisation avec reprise)
\item « C'est une opportunité énorme. » (Augmentatif ou adjectif de relief)
\end{enumerate}
\end{exercice}

\begin{corrige}[Exercice de thème 3, 4, 5]
\begin{enumerate}
\item \textbf{Yo hablo español con fluidez.} L'emphase sur le sujet « Moi » s'exprime en espagnol par le pronom tonique \esp{yo} en tête (le sujet nul est la norme, sa présence est marquante). La « reprise pronominale » (redoublement clitique \esp{le/la/lo}) ne s'applique qu'aux compléments d'objet (\esp{A Juan \textbf{lo} vi}), pas au sujet.
\item \esp{El inglés, lo domino.} Le complément \esp{el inglés} passe en tête (topicalisation) et est repris par le pronom d'objet direct \esp{lo}. Jamais \esp{El inglés, domino} sans pronom.
\item \esp{Es una oportunidad grandísima.} (Augmentatif suffixal \esp{-ísima} sur \esp{grande} $\rightarrow$ \esp{grandísima} par apocope). Variantes : \esp{Es una oportunidad inmensa} (adjectif fort) ou \esp{Es una gran oportunidad} (apocope de \esp{grande} devant nom singulier). \emph{Orthographe :} \esp{oportunidad} (un seul \emph{p}, pas d'accent).
\end{enumerate}
\end{corrige}

\courssub{Entraînement de version (espagnol $\rightarrow$ français)}

Traduire une phrase emphatique, ce n'est pas calquer : c'est \cle{retrouver l'effet de sens} en français naturel. Observez d'abord, puis traduisez.

\begin{exoresolu}[Version guidée : phrases 1 et 2]
\textbf{Énoncé 1.} \esp{Fue ayer cuando firmé el contrato.}

\tcblower

\textbf{Analyse.} Clivée temporelle : \esp{Fue + ayer + cuando}. Le français rend la clivée par « c'est... que ».

\textbf{Traduction.} « C'est hier que j'ai signé le contrat. »

Piège : ne pas traduire \esp{cuando} par « quand » ici ; la clivée française exige « que ».

\bigskip
\textbf{Énoncé 2.} \esp{A Juan, a él le ofrecieron el puesto.}

\tcblower

\textbf{Analyse.} Double reprise : \esp{A Juan} + \esp{a él} + pronom \esp{le} (leísmo standard pour personne masculine). Le français ne peut pas tout reprendre ; on choisit une tournure naturelle.

\textbf{Traduction.} « C'est à Juan qu'on a offert le poste. » ou « Juan, c'est à lui qu'on a offert le poste. »

Piège : \esp{ofrecieron} est un pluriel impersonnel (« on a offert »), pas « ils ont offert » sauf si le contexte désigne des personnes précises.
\end{exoresolu}

\begin{exercice}[Version : phrases 3, 4 et 5]
Traduisez en français naturel :
\begin{enumerate}
\item \esp{Lo fundamental es la constancia.}
\item \esp{Ganas muchas experiencias tú.}
\item \esp{Es un trabajador excepcional.}
\end{enumerate}
\end{exercice}

\begin{corrige}[Exercice de version 3, 4, 5]
\begin{enumerate}
\item « Ce qui est fondamental, c'est la constance. » (Ser de relieve : \esp{lo fundamental} = « ce qui compte vraiment ».)
\item « C'est toi qui gagnes beaucoup d'expérience. » ou « Toi, tu accumules vraiment les expériences. » (Postposition du sujet \esp{tú} = insistance sur la personne.)
\item « C'est un travailleur exceptionnel. » (Adjectif postposé + article : simple emphase qualitative ; on peut renforcer : « un employé vraiment exceptionnel ».)
\end{enumerate}
\end{corrige}

\courssub{Sujet type Bac : épreuve unique (2 h)}

Voici un sujet complet, calibré sur l'épreuve officielle. Entraînez-vous en conditions réelles : 2 heures, sans dictionnaire au début, puis auto-correction.

\begin{textebox}[Un empleo para quedarse]
En Douala, la capital económica de Camerún, encontrar un empleo estable no es fácil. Sin embargo, fue en un pequeño taller de reparación de motos donde Aminatou, de veintidós años, encontró su oportunidad. Lo que más le sorprendió fue la confianza del jefe: «A ti te dejo las llaves porque trabajas con seriedad». A ella, que había hecho prácticas en un garaje de Yaoundé, le ofrecieron un contrato fijo después de solo tres meses. Hoy, Aminatou dirige el taller y enseña el oficio a dos jóvenes aprendices. «El inglés lo aprendí sola, y el español, lo estudio por las noches», explica con orgullo. Su historia demuestra que lo fundamental no es el diploma, sino la constancia y las ganas de aprender. Para ella, cada día de trabajo es una oportunidad grandísima de crecer.
\end{textebox}

\textbf{Glossaire.} \esp{el taller} : l'atelier ; \esp{las llaves} : les clés ; \esp{el contrato fijo} : le CDI ; \esp{el oficio} : le métier ; \esp{los aprendices} : les apprentis ; \esp{las ganas de aprender} : l'envie d'apprendre ; \esp{crecer} : grandir, progresser.

\begin{exercice}[Épreuve type Bac (2 h)]
\textbf{I. Comprensión lectora (20 min).} Répondez en espagnol, avec vos propres mots :
\begin{enumerate}
\item ¿Dónde encontró Aminatou su oportunidad y por qué es sorprendente?
\item ¿Qué cualidades de Aminatou valoran en el taller? Justifique con dos elementos del texto.
\end{enumerate}

\textbf{II. Gramática : el relieve (20 min).}
\begin{enumerate}
\item Relevez dans le texte \textbf{une phrase clivée} et \textbf{un ser de relieve} ; identifiez l'élément mis en relief dans chaque cas.
\item Relevez \textbf{deux reprises pronominales} (redoublement clitique).
\item Transformez : \esp{Aminatou dirige el taller hoy.} $\rightarrow$ clivée sur \esp{hoy}.
\end{enumerate}

\textbf{III. Léxico (10 min).} Trouvez dans le texte l'équivalent de : « CDI », « stage », « apprentis », « salaire » (attention, piège : ce mot n'y figure pas — dites-le en espagnol).

\textbf{IV. Expresión escrita (60 min).} \esp{Aminatou busca un nuevo pasante. Escriba usted una carta de motivación (unas 120 palabras) para solicitar ese puesto: preséntese, explique sus motivaciones y sus competencias, y termine con una fórmula de cortesía.}
\end{exercice}

\begin{corrige}[Sujet type Bac — éléments de réponse]
\textbf{I.} 1. \esp{Encontró su oportunidad en un pequeño taller de reparación de motos; es sorprendente porque es un lugar modesto y no una gran empresa.} 2. \esp{El jefe valora su seriedad («trabajas con seriedad») y su experiencia previa (las prácticas en un garaje de Yaoundé).}

\textbf{II.} 1. Clivée : \esp{fue en un pequeño taller... donde Aminatou encontró su oportunidad} (relief du lieu). Ser de relieve : \esp{Lo que más le sorprendió fue la confianza del jefe} (relief du sujet \esp{la confianza}). 2. Reprises : \esp{A ella... le ofrecieron un contrato fijo} ; \esp{A ti te dejo las llaves}. 3. \esp{Es hoy cuando Aminatou dirige el taller.} (ou \esp{Es hoy cuando dirige Aminatou el taller}).

\textbf{III.} CDI = \esp{un contrato fijo} ; stage = \esp{las prácticas} ; apprentis = \esp{los aprendices} ; salaire = \esp{el salario / el sueldo} (absent du texte).

\textbf{IV.} Barème indicatif : respect du format lettre (appel, date, clôture) 4 pts ; 120 mots $\pm$ 10 \% 2 pts ; au moins deux procédés d'emphase corrects 4 pts ; lexique de l'emploi 4 pts ; correction grammaticale (\esp{usted} constant, accords) 6 pts.
\end{corrige}

\begin{popfigure}
\begin{tikzpicture}[font=\small]
\draw[-{Stealth}, thick, popblue] (0,0) -- (14.5,0);
\foreach \x/\t in {0/0h, 2.4/20', 6/50', 8.4/1h10, 13/1h50, 14.4/2h}
  \draw[popdark] (\x,0.12) -- (\x,-0.12) node[below=2pt, popdark]{\t};
\draw[fill=popblueL, draw=popblue, rounded corners=2pt] (0.1,0.5) rectangle (2.3,1.4) node[midway, align=center]{Lecture\\du texte};
\draw[fill=poporangeL, draw=poporange, rounded corners=2pt] (2.5,0.5) rectangle (5.9,1.4) node[midway, align=center]{Questions\\(compréhension)};
\draw[fill=popgreenL, draw=popgreen, rounded corners=2pt] (6.1,0.5) rectangle (8.3,1.4) node[midway, align=center]{Emphase\\+ lexique};
\draw[fill=poppurpleL, draw=poppurple, rounded corners=2pt] (8.5,0.5) rectangle (12.9,1.4) node[midway, align=center]{Rédaction de la lettre\\(120 palabras)};
\draw[fill=poppinkL, draw=poppink, rounded corners=2pt] (13.1,0.5) rectangle (14.3,1.4) node[midway, align=center]{Re-\\lecture};
\end{tikzpicture}
\legende{Gestion du temps pour l'épreuve unique de 2 h : cinq étapes, du texte à la relecture.}
\end{popfigure}

\begin{attention}[Dernier conseil Bac]
N'inventez \textbf{jamais} de formules de politesse : utilisez uniquement les modèles canoniques appris (\esp{Muy señor mío:}, \esp{Le saluda atentamente,}, \esp{Quedo a su disposición para cualquier información adicional}). Ne mélangez pas \esp{tú} et \esp{usted} : relisez chaque verbe. Enfin, dans les clivées passives (rares mais possibles), vérifiez l'accord du participe : \esp{Fue el proyecto el que fue realizado por mí} — \esp{realizado} s'accorde avec \esp{el proyecto}.
\end{attention}

\begin{savaistu}[Le poids de l'expression écrite au Bac camerounais]
L'épreuve d'\esp{Expresión escrita} pèse souvent entre 20 et 25 \% de la note finale d'espagnol. Bonne nouvelle : les correcteurs valorisent la \cle{structure}. Une lettre bien organisée (appel, trois paragraphes, clôture) contenant trois emphases correctes (une clivée, un ser de relieve, une reprise) garantit souvent la moyenne, même avec quelques fautes mineures. L'emphase est donc un investissement rentable : cinq procédés à maîtriser, des points assurés.
\end{savaistu}

\courssub{Grille d'auto-évaluation}

\begin{definition}[Auto-évaluation]
L'auto-évaluation est un outil \cle{métacognitif} : l'élève juge lui-même son niveau de maîtrise en cochant, pour chaque objectif, « Je sais », « Je hésite » ou « Je ne sais pas ». Exemple : « Je sais transformer une phrase neutre en clivée » [Oui / Presque / Non]. Elle permet de cibler les révisions avant l'examen au lieu de tout relire au hasard.
\end{definition}

Cochez honnêtement chaque ligne, puis revoyez les sections correspondantes de la leçon pour tout ce qui n'est pas « Je sais ».

\begin{center}
\begin{tabularx}{\linewidth}{|X|c|c|c|}
\hline
\rowcolor{popblueL} \textbf{Compétence} & \textbf{Je sais} & \textbf{Je hésite} & \textbf{Je ne sais pas} \\
\hline
Identifier une clivée (\esp{Fue... quien/que}) dans un texte & & & \\
\hline
Transformer une phrase neutre en clivée & & & \\
\hline
Employer le ser de relieve (\esp{Lo que... es...}) & & & \\
\hline
Construire une reprise pronominale (\esp{A Juan le...}) & & & \\
\hline
Former diminutifs et augmentatifs (\esp{-ito}, \esp{-ón}, \esp{-ísimo}) & & & \\
\hline
Traduire une phrase emphatique en français naturel & & & \\
\hline
Rédiger une lettre de motivation de 120 mots & & & \\
\hline
Utiliser les formules de politesse formelles (\esp{usted}) & & & \\
\hline
Mobiliser le lexique de l'emploi (10 mots clés) & & & \\
\hline
Gérer mon temps sur une épreuve de 2 h & & & \\
\hline
\end{tabularx}
\end{center}

\begin{experience}[Bilan en binôme (10 min)]
Échangez votre grille avec un camarade. Choisissez ensemble \textbf{deux} cases « Je hésite » ou « Je ne sais pas » : l'un explique la règle à l'autre avec un exemple, puis inversez les rôles. Terminez en rédigeant chacun, sur une feuille, une phrase emphatique par procédé (cinq phrases au total) ; votre binôme vérifie la correction à l'aide de la fiche mémo 6.1. La leçon E18 est réussie lorsque toute la ligne « emphase » est cochée « Je sais ».
\end{experience}

\newpage

\coursec{Activité d'intégration}

\courssub{Objectifs de l'activité}

Cette activité d'intégration clôt la leçon sur l'emphase et la lettre de motivation. Elle vise à faire passer les élèves de l'écrit à l'oral, en mobilisant simultanément toutes les compétences travaillées dans le module : compréhension et expression écrites (la lettre), compréhension et expression orales (l'entretien), grammaire (les mécanismes du relief expressif) et lexique (le monde de l'emploi). À l'issue de la simulation, l'élève sera capable de :

\begin{itemize}
\item lire à voix haute sa lettre de motivation avec une prononciation claire et une intonation expressive ;
\item défendre oralement sa candidature en employant les mécanismes de l'emphase : phrases clivées avec \esp{ser}, redoublement du pronom (\esp{A mí me...}), ordre des mots marqué, \esp{ser} de relief ;
\item comprendre et répondre à des questions inattendues en situation formelle (registre \esp{usted}) ;
\item utiliser spontanément le lexique de l'emploi étudié dans la leçon ;
\item évaluer la production d'un pair selon des critères explicites et argumenter un choix collectif.
\end{itemize}

\courssub{Situation}

\textbf{Contexte.} La multinationale agroalimentaire \esp{ChocoCamer S.A.}, installée à Douala, exporte du cacao transformé vers l'Espagne et le Mexique. Elle recherche un(e) jeune assistant(e) bilingue français-espagnol pour son service commercial. L'entreprise organise un \esp{job dating} : chaque candidat dispose d'un entretien éclair de cinq minutes face au responsable des ressources humaines. Votre classe de Terminale A a été invitée à participer à la simulation grandeur nature organisée par le club d'espagnol du lycée.

\begin{textebox}[Anuncio de empleo — ChocoCamer S.A., Douala]
\textbf{¡TRABAJA CON NOSOTROS!}

ChocoCamer S.A., empresa camerunesa de transformación y exportación de cacao, busca un(a) \textbf{asistente comercial bilingüe} (francés-español) para su oficina de Douala.

\textbf{Perfil buscado :} joven dinámico, con bachillerato y buen nivel de español ; capacidad de redacción (cartas, correos comerciales) ; gusto por el contacto con clientes extranjeros ; disponibilidad inmediata.

\textbf{Ofrecemos :} contrato de un año renovable, formación continua, ambiente internacional.

\textbf{Proceso de selección :} envío de una carta de motivación (120 palabras) + entrevista breve en español con el jefe de Recursos Humanos.

\emph{Es tu talento el que nos interesa. ¡Anímate!}

Contacto : rrhh@chococamer.cm
\end{textebox}

\textbf{Glossaire :} \esp{asistente comercial bilingüe} = assistant commercial bilingue ; \esp{bachillerato} = baccalauréat ; \esp{gusto por el contacto} = goût du contact ; \esp{contrato renovable} = contrat renouvelable ; \esp{entrevista} = entretien ; \esp{¡Anímate!} = Lance-toi !

\begin{attention}[L'emphase dès l'annonce]
Remarquez la dernière phrase de l'annonce : \esp{Es tu talento el que nos interesa.} Il s'agit d'une \textbf{phrase clivée} : la forme neutre serait \esp{Tu talento nos interesa} (« Ton talent nous intéresse »). En clivant, l'annonceur met en relief \esp{tu talent} : « C'est ton talent qui nous intéresse. » Notez l'accord : \esp{el que} reprend \esp{tu talent} (masculin singulier déterminé). L'emphase n'est donc pas un ornement réservé à la littérature : elle est partout, y compris dans les offres d'emploi.
\end{attention}

\courssub{Consignes}

Travail en binômes : un élève joue le rôle du \textbf{Candidat}, l'autre celui du \textbf{Recruteur RH}. Durée totale : 30 minutes.

\begin{enumerate}
\item \textbf{Préparation (5 min).} Le Candidat relit sa lettre de motivation rédigée pendant l'atelier d'expression écrite et souligne deux tournures emphatiques qu'il devra prononcer avec force. Le Recruteur prépare trois questions difficiles à partir de la banque ci-dessous, en veillant à utiliser le registre formel (\esp{usted}).
\item \textbf{Présentation de la lettre (1 min).} Le Candidat lit sa lettre à voix haute, debout, en regardant le Recruteur. Interdiction de lire mot à mot sans intonation : l'emphase doit s'entendre.
\item \textbf{Entretien (4 min).} Le Recruteur pose ses trois questions. Le Candidat doit répondre à chacune en employant \textbf{au moins une structure emphatique} : phrase clivée (\esp{Es mi experiencia la que...}), redoublement pronominal (\esp{A mí me apasiona...}), ou mise en relief par l'ordre des mots (\esp{Eso sí lo sé hacer}).
\item \textbf{Échange des rôles (5 min).} On inverse : le Recruteur devient Candidat et inversement, avec une autre lettre et d'autres questions.
\item \textbf{Évaluation croisée (5 min).} Chaque binôme remplit la fiche d'observation : nombre de tournures emphatiques utilisées, correction grammaticale, richesse du lexique de l'emploi, registre formel respecté.
\item \textbf{Débrief collectif et vote (10 min).} La classe désigne les binômes les plus convaincants. Deux ou trois candidats passent devant toute la classe. Discussion guidée par le professeur : \emph{Quelles emphases ont été les plus persuasives ? Pourquoi ?} Vote à main levée du \esp{Mejor Candidato}, qui reçoit symboliquement le \esp{contrato} de ChocoCamer.
\end{enumerate}

\begin{exemplebox}[Banque de questions pièges du Recruteur]
\begin{itemize}
\item \esp{¿Por qué usted y no otro candidato?} — Pourquoi vous et pas un autre candidat ?
\item \esp{¿Qué es lo que más le apasiona de este puesto?} — Qu'est-ce qui vous passionne le plus dans ce poste ?
\item \esp{¿No le parece que le falta experiencia profesional?} — Ne pensez-vous pas qu'il vous manque de l'expérience professionnelle ?
\item \esp{¿Qué sabe usted de nuestra empresa?} — Que savez-vous de notre entreprise ?
\item \esp{¿Cómo reaccionaría usted ante un cliente enfadado?} — Comment réagiriez-vous face à un client en colère ?
\end{itemize}
\end{exemplebox}

\courssub{Ressources à mobiliser}

\begin{aretenir}[Boîte à outils de l'emphase à l'oral]
\begin{itemize}
\item \textbf{Phrase clivée :} \esp{Es} + élément mis en relief + \esp{que / quien / el que / lo que} + verbe. \esp{Es mi nivel de español lo que me distingue.} « C'est mon niveau d'espagnol qui me distingue. »
\item \textbf{Redoublement pronominal :} \esp{A mí me gusta... / A mí me apasiona...} — le pronom tonique \esp{a mí} renforce la personne qui parle.
\item \textbf{Ordre des mots marqué :} placer le complément en tête de phrase : \esp{Eso sí lo sé hacer.} « Ça, je sais le faire. »
\item \textbf{\esp{Ser} de relief :} \esp{Lo que yo quiero es aprender.} « Ce que je veux, c'est apprendre. »
\item \textbf{Augmentatifs et diminutifs expressifs} (avec prudence au registre formel) : \esp{un reto grandísimo} « un énorme défi ».
\end{itemize}
\end{aretenir}

\begin{propriete}[Rappel : l'accord dans la phrase clivée]
Après l'élément mis en relief, le relatif s'accorde avec lui :
\begin{itemize}
\item antécédent \textbf{déterminé} (article, possessif) : \esp{el que / la que / los que / las que} — \esp{Es mi experiencia \textbf{la que} me distingue.} ; \esp{Fueron mis profesores \textbf{los que} me animaron.}
\item antécédent \textbf{indéterminé ou indéfini} : \esp{que} — \esp{Es talento lo que buscamos} reste possible, mais on préfère : \esp{Es un candidato \textbf{que} habla dos lenguas.}
\item personne mise en relief : \esp{quien / quienes} — \esp{Es usted \textbf{quien} decide.}
\item idée globale, verbe à l'infinitif : \esp{lo que} — \esp{Lo que me interesa \textbf{es} aprender.}
\end{itemize}
\end{propriete}

\textbf{Lexique de l'emploi à réactiver :} \esp{el puesto} (le poste), \esp{la experiencia} (l'expérience), \esp{las ganas de aprender} (l'envie d'apprendre), \esp{trabajar en equipo} (travailler en équipe), \esp{estar disponible} (être disponible), \esp{el contrato} (le contrat), \esp{la entrevista} (l'entretien), \esp{los clientes extranjeros} (les clients étrangers).

\textbf{Registre formel :} vouvoiement obligatoire (\esp{usted} + 3\textsuperscript{e} personne du singulier), formules de politesse : \esp{Muchas gracias por su atención} (« Merci beaucoup de votre attention »), \esp{Quedo a su disposición} (« Je reste à votre disposition »), \esp{Es un honor para mí} (« C'est un honneur pour moi »).

\begin{attention}[Pièges à éviter pendant l'entretien]
\begin{itemize}
\item Ne jamais répondre par un simple \esp{sí} ou \esp{no} : développer toujours avec une mise en relief.
\item Éviter le tutoiement (\esp{tú}) avec le recruteur : \esp{¿Qué le parece?} et non \esp{¿Qué te parece?}
\item Ne pas calquer le français : on dit \esp{tengo muchas ganas de trabajar} (« j'ai très envie de travailler ») ; le calque direct du français est incorrect.
\item Attention à l'accord dans la phrase clivée : \esp{Es mi experiencia \textbf{la que} me distingue} (antécédent déterminé : on préfère \esp{la que} à \esp{lo que}).
\item \esp{estar disponible} et non \esp{ser disponible} : la disponibilité est un état, pas une caractéristique permanente.
\end{itemize}
\end{attention}

\begin{popfigure}
\begin{tikzpicture}[node distance=0.6cm, every node/.style={font=\small}]
\node[draw=popblue, fill=popblueL, rounded corners, align=center, minimum width=3.4cm] (carta) {\textbf{La carta}\\ (escrita, 120 palabras)};
\node[draw=poporange, fill=poporangeL, rounded corners, align=center, minimum width=3.4cm, right=1.6cm of carta] (entrevista) {\textbf{La entrevista}\\ (oral, 5 minutos)};
\node[draw=popgreen, fill=popgreenL, rounded corners, align=center, minimum width=3.4cm, right=1.6cm of entrevista] (contrato) {\textbf{El contrato}\\ (¡éxito!)};
\draw[-{Stealth}, very thick, popdark] (carta) -- (entrevista) node[midway, above, font=\footnotesize\itshape] {selección};
\draw[-{Stealth}, very thick, popdark] (entrevista) -- (contrato) node[midway, above, font=\footnotesize\itshape] {decisión};
\node[below=0.5cm of entrevista, align=center, text=popmuted, font=\footnotesize\itshape] {La émfasis : clivar, redoblar, desplazar};
\end{tikzpicture}
\legende{Le parcours professionnel reproduit par la simulation : de la lettre écrite à l'entretien oral.}
\end{popfigure}

\courssub{Proposition de production et corrigé commenté}

\begin{exoresolu}[Entretien modèle : le Candidat face au Recruteur]
Énoncé : rédiger et jouer les réponses du Candidat aux trois questions pièges du Recruteur, en utilisant au moins trois mécanismes d'emphase différents et le lexique de l'emploi.

\tcblower

\textbf{Recruteur :} \esp{¿Por qué usted y no otro candidato?}

\textbf{Candidat :} \esp{Es mi dominio del español lo que me distingue de los demás. Otros candidatos tienen quizás más experiencia, pero a mí me apasiona el contacto con los clientes extranjeros, y eso no se aprende en los libros.}

« C'est ma maîtrise de l'espagnol qui me distingue des autres. D'autres candidats ont peut-être plus d'expérience, mais moi, le contact avec les clients étrangers me passionne, et cela ne s'apprend pas dans les livres. »

\textbf{Recruteur :} \esp{¿No le parece que le falta experiencia profesional?}

\textbf{Candidat :} \esp{Experiencia, es verdad que todavía me falta. Pero lo que yo ofrezco es entusiasmo y ganas de aprender. A mí me formaron en un liceo exigente de Yaoundé, y trabajar en equipo es precisamente lo que mejor sé hacer.}

« L'expérience, il est vrai qu'elle me manque encore. Mais ce que j'offre, moi, c'est de l'enthousiasme et l'envie d'apprendre. C'est dans un lycée exigeant de Yaoundé que j'ai été formé, et travailler en équipe est précisément ce que je sais le mieux faire. »

\textbf{Recruteur :} \esp{¿Qué es lo que más le apasiona de este puesto?}

\textbf{Candidat :} \esp{Lo que más me atrae es el ambiente internacional de ChocoCamer. Es el cacao camerunés el que merece ser conocido en el mundo, y a mí me encantaría contribuir a ello. Quedo a su disposición y le agradezco mucho su atención.}

« Ce qui m'attire le plus, c'est l'environnement international de ChocoCamer. C'est le cacao camerounais qui mérite d'être connu dans le monde, et j'adorerais y contribuer. Je reste à votre disposition et vous remercie de votre attention. »

\textbf{Commentaire des choix de langue.}
\begin{itemize}
\item \textbf{Phrase clivée} : \esp{Es mi dominio del español lo que me distingue} — le candidat met en avant son atout principal dès la première phrase ; \esp{lo que} est ici correct car l'antécédent est une notion abstraite globale (on pourrait aussi dire \esp{Es mi dominio del español \textbf{el que} me distingue}).
\item \textbf{Redoublement pronominal} : \esp{a mí me apasiona}, \esp{a mí me formaron}, \esp{a mí me encantaría} — insistance sur la personne du candidat, effet persuasif.
\item \textbf{Ordre des mots marqué} : \esp{Experiencia, es verdad que todavía me falta} — le complément placé en tête concède le point faible avant de le retourner en argument : stratégie rhétorique classique de la concession.
\item \textbf{\esp{Ser} de relief} : \esp{lo que yo ofrezco es entusiasmo}, \esp{lo que más me atrae es...} — structure très naturelle à l'oral.
\item \textbf{Clivée avec antécédent déterminé} : \esp{Es el cacao camerunés \textbf{el que} merece...} — \esp{el que} (et non \esp{que}) car l'antécédent porte l'article défini ; ancrage camerounais valorisant.
\item \textbf{Registre formel} : vouvoiement constant (\esp{le falta}, \esp{le agradezco}), conditionnel de politesse (\esp{me encantaría}), formule de clôture \esp{Quedo a su disposición} reprise de la lettre de motivation.
\end{itemize}
\end{exoresolu}

\begin{exemplebox}[Grille d'évaluation croisée (fiche d'observation)]
\begin{center}
\begin{tabularx}{\linewidth}{|X|c|}
\hline
\rowcolor{popblueL} \textbf{Critère} & \textbf{Points (sur 2)} \\
\hline
Au moins trois tournures emphatiques différentes & \\
\hline
Correction grammaticale (accords, conjugaisons) & \\
\hline
Lexique de l'emploi (au moins quatre termes) & \\
\hline
Registre formel respecté (\esp{usted}, politesse) & \\
\hline
Prononciation et intonation expressives & \\
\hline
\end{tabularx}
\end{center}
Barème : 2 = réussi, 1 = partiellement réussi, 0 = absent. Total sur 10.
\end{exemplebox}

\courssub{Bilan de l'activité}

Cette simulation a permis de vérifier que les savoir-faire de la leçon sont réellement intégrés : les élèves ne se contentent plus de reconnaître l'emphase dans un texte, ils la produisent spontanément sous la pression d'une situation d'interaction réelle. Le passage de l'écrit (la lettre) à l'oral (l'entretien) reproduit fidèlement le parcours professionnel authentique : on est d'abord sélectionné sur dossier, puis sur entretien.

Le débrief collectif a mis en évidence que les tournures les plus convaincantes sont la phrase clivée (\esp{Es... lo que...}) pour affirmer un atout, et le redoublement \esp{A mí me...} pour exprimer la motivation personnelle. Le vote du \esp{Mejor Candidato} développe enfin l'esprit critique : évaluer, c'est aussi apprendre.

\begin{aretenir}[Ce que je retiens de cette activité]
À l'oral comme à l'écrit, l'emphase espagnole repose sur trois gestes : \textbf{cliver} (\esp{Es X lo que...}), \textbf{redoubler} (\esp{A mí me...}), \textbf{déplacer} (\esp{Eso sí lo sé}). Dans un entretien formel, je vouvoie toujours, je ne réponds jamais par un seul mot, et je transforme mes faiblesses en arguments grâce à la concession suivie d'une mise en relief.
\end{aretenir}

Pour prolonger l'activité à la maison, chaque élève pourra enregistrer sur son téléphone une version audio de deux minutes de sa défense de candidature, en veillant à y placer au moins quatre tournures emphatiques différentes : c'est un excellent entraînement à l'épreuve orale du baccalauréat.

\newpage

\coursec{Méthodes et exercices résolus}

\courssub{Fiche méthode 1 : Construire une phrase clivée (\esp{oración hendida})}

\begin{methode}[Mettre un élément en relief par la clivage]
La phrase clivée espagnole suit le schéma : \esp{Fue + élément mis en relief + que + reste de la phrase}. Pour la maîtriser, procède en quatre étapes :
\begin{enumerate}
\item \cle{Identifier} l'élément à mettre en valeur (personne, objet, lieu, moment).
\item \cle{Placer} \esp{Fue} (ou \esp{Es}, si le fait est présent) en tête de phrase.
\item \cle{Introduire} l'élément mis en relief, puis la conjonction \esp{que}.
\item \cle{Vérifier} la concordance des temps : \esp{Fue... que + passé} ; \esp{Es... que + présent}.
\end{enumerate}
Réflexe de traduction : le français « C'est... qui/que » se rend presque toujours par \esp{Es/Fue... que}. Attention : en espagnol, on n'emploie JAMAIS \esp{quien} dans la clivée, même pour une personne.
\end{methode}

\begin{exoresolu}[Transformation : de la phrase neutre à la phrase clivée]
Énoncé. Transformez les phrases suivantes en phrases clivées en mettant en relief l'élément souligné : 1. \esp{\underline{Mi abuela} me enseñó a escribir cartas.} 2. \esp{Encontré el empleo \underline{en Douala}.} 3. \esp{\underline{La carta de motivación} decidió mi destino.}
\tcblower
Solution détaillée.
\begin{enumerate}
\item L'élément à mettre en relief est une personne (\esp{mi abuela}). On applique le schéma : \esp{Fue mi abuela quien...} ? Non ! La clivée exige \esp{que}, même pour les personnes. Réponse : \esp{Fue mi abuela la que me enseñó a escribir cartas} ou plus simplement \esp{Fue mi abuela quien me enseñó...} est à éviter au Bac ; la forme sûre est : \esp{Fue mi abuela la que me enseñó a escribir cartas.} « C'est ma grand-mère qui m'a appris à écrire des lettres. »
\item L'élément est un complément de lieu : \esp{Fue en Douala donde encontré el empleo.} « C'est à Douala que j'ai trouvé l'emploi. » Remarque : avec un lieu, on peut employer \esp{donde} ; avec un moment, \esp{cuando} : \esp{Fue en 2023 cuando...}.
\item L'élément est le sujet : \esp{Fue la carta de motivación la que decidió mi destino.} « C'est la lettre de motivation qui a décidé de mon destin. »
\end{enumerate}
Conclusion : la clivée espagnole repose sur \esp{Fue/Es + X + que/la que/donde}. Le temps de \esp{ser} s'accorde avec celui de la phrase d'origine (ici, passé simple \esp{enseñó} $\rightarrow$ \esp{Fue}).
\end{exoresolu}

\courssub{Fiche méthode 2 : \esp{Ser de relieve} et répétition du pronom}

\begin{methode}[Souligner le sujet avec \esp{ser de relieve}]
Le \cle{ser de relieve} permet d'insister sur l'agent ou l'auteur d'une action :
\begin{enumerate}
\item Structure : \esp{Es/Fue + sujet + quien/el que + verbe}. Exemple : \esp{Fui yo quien escribió la carta.} « C'est moi qui ai écrit la lettre. »
\item Le verbe conjugué s'\cle{accorde} avec le sujet mis en relief : \esp{Fui yo quien escribió} (première personne), \esp{Fuiste tú quien lo hiciste} ? Non : \esp{Fuiste tú quien lo hizo} est incorrect ; on dit \esp{Fuiste tú quien lo hiciste}.
\item La \cle{répétition du pronom} est un autre procédé d'emphase : \esp{Yo, lo que quiero, es trabajar.} ; \esp{A mí, lo que me importa es el esfuerzo.} Le pronom tonique (\esp{a mí, a ti, a nosotros}) renforce le pronom atone obligatoire (\esp{me, te, nos}).
\item Réflexe de traduction : « Moi, je... » $\rightarrow$ \esp{A mí me...} ; « C'est moi qui... » $\rightarrow$ \esp{Soy yo quien...}.
\end{enumerate}
\end{methode}

\begin{exoresolu}[Thème : emphase sur la personne]
Énoncé. Traduisez en espagnol : 1. « C'est moi qui ai envoyé la candidature. » 2. « Moi, ce qui me passionne, c'est le commerce international. » 3. « C'est toi qui dois décider, pas tes parents. »
\tcblower
Solution détaillée.
\begin{enumerate}
\item Sujet « moi » + passé composé $\rightarrow$ passé simple en espagnol dans ce contexte : \esp{Fui yo quien envié la candidatura.} Accord du verbe à la première personne : \esp{envié}. « Candidature » se dit \esp{candidatura} (ou \esp{solicitud}), jamais \esp{candidación}.
\item « Moi » mis en relief $\rightarrow$ \esp{A mí} + pronom atone obligatoire \esp{me} : \esp{A mí, lo que me apasiona es el comercio internacional.} Sans le \esp{me}, la phrase est fautive. « Passionne » : \esp{apasiona} (présent, car vérité actuelle).
\item Deuxième personne : \esp{Eres tú quien debe decidir, no tus padres.} Le verbe \esp{debe} reste à la troisième personne après \esp{quien} quand le sujet est \esp{tú} ? Non : la forme rigoureuse est \esp{Eres tú quien debes decidir}. Les deux tournures existent dans l'usage, mais au Bac, accorde le verbe avec la personne du sujet : \esp{debes}.
\end{enumerate}
Conclusion : l'emphase sur la personne exige deux réflexes : l'accord du verbe avec le sujet mis en relief, et la redondance du pronom atone après \esp{a mí / a ti / a nosotros}.
\end{exoresolu}

\courssub{Fiche méthode 3 : L'ordre des mots et les diminutifs expressifs}

\begin{methode}[Jouer sur l'ordre des mots et les suffixes]
\begin{enumerate}
\item \cle{Antéposition du complément} : en espagnol, placer le complément en tête le met en valeur : \esp{Mucho esfuerzo puso ella en esa carta.} « Elle a mis beaucoup d'efforts dans cette lettre. » Le sujet suit souvent le verbe (inversion).
\item \cle{Antéposition de l'adjectif} : \esp{una gran mujer} (une femme remarquable) n'a pas le même sens que \esp{una mujer grande} (une femme de grande taille). L'adjectif avant le nom porte une valeur subjective, emphatique.
\item \cle{Diminutifs} (\esp{-ito/-ita, -cito}) : ils expriment l'affection, la délicatesse ou l'atténuation : \esp{una cartita} (une petite lettre touchante), \esp{un momentito} (un petit moment).
\item \cle{Augmentatifs} (\esp{-ón/-ona, -azo}) : ils expriment la grandeur ou l'intensité : \esp{un trabajón} (un travail énorme), \esp{un exitazo} (un immense succès).
\item Réflexe : à l'écrit, une seule inversion bien placée suffit ; n'en abuse pas dans une lettre formelle.
\end{enumerate}
\end{methode}

\begin{exoresolu}[Version : repérer et traduire l'emphase]
Énoncé. Traduisez en français en rendant la valeur emphatique : \esp{1. Poco a poco construyó ella su futuro. 2. Era una gran profesional. 3. Con una cartita sencilla consiguió un exitazo.}
\tcblower
Solution détaillée.
\begin{enumerate}
\item \esp{Poco a poco} en tête + inversion sujet-verbe (\esp{construyó ella}) : l'emphase porte sur la lenteur progressive. Traduction : « Petit à petit, c'est elle qui a construit son avenir » ou « C'est petit à petit qu'elle a bâti son avenir. »
\item Adjectif antéposé \esp{gran} : valeur subjective d'admiration, non de taille. Traduction : « C'était une remarquable professionnelle » (et non « une professionnelle grande »). \esp{Gran} est l'apocope de \esp{grande} devant un nom singulier.
\item Diminutif \esp{cartita} : « une modeste petite lettre » ; augmentatif \esp{exitazo} : « un succès retentissant ». Traduction complète : « Avec une simple petite lettre, elle a décroché un succès retentissant. »
\end{enumerate}
Conclusion : en version, toute inversion, tout adjectif antéposé, tout suffixe diminutif ou augmentatif doit être rendu en français par un effet équivalent (clivage, adverbe, intensificateur).
\end{exoresolu}

\courssub{Fiche méthode 4 : Rédiger une \esp{carta de motivación} (120 mots)}

\begin{methode}[Structurer la lettre de motivation en cinq blocs]
\begin{enumerate}
\item \cle{En-tête et appel} : lieu et date (\esp{Yaoundé, a 15 de mayo de 2025}), puis formule d'appel formelle : \esp{Estimado señor Director:} ou \esp{Muy señor mío:} (deux-points obligatoires, jamais de virgule).
\item \cle{Présentation} : qui tu es et le poste visé : \esp{Me llamo... y me dirijo a usted para solicitar el puesto de... que vi anunciado en...}.
\item \cle{Motivations et qualités} : relie tes compétences au poste avec des connecteurs : \esp{Además, cuento con experiencia en... ; Por otra parte, domino el francés y el español, lo que me permite...}.
\item \cle{Demande} : \esp{Quedo a su disposición para una entrevista.}
\item \cle{Formule de politesse} : \esp{Le saluda atentamente,} ou \esp{Atentamente,} suivie de la signature.
\end{enumerate}
Réflexes : vouvoiement partout (\esp{usted}) ; compter les mots (110--130) ; un connecteur par paragraphe ; une emphase bien placée (\esp{Es precisamente esta experiencia la que...}) valorise la copie.
\end{methode}

\begin{exoresolu}[Rédaction type Bac : lettre de motivation]
Énoncé. \esp{Has visto un anuncio: una empresa de exportación de cacao de Douala busca un secretario bilingüe. Escribe una carta de motivación de unas 120 palabras.}
\tcblower
Solution détaillée. Plan respecté : appel, présentation, motivations, demande, politesse. Modèle rédigé :

\esp{Douala, a 10 de junio de 2025}

\esp{Estimado señor Director:}

\esp{Me llamo Andrés Mbarga y me dirijo a usted para solicitar el puesto de secretario bilingüe anunciado en su empresa.}

\esp{Acabo de obtener el bachillerato y domino el francés y el español, lo que me permite redactar documentos en ambas lenguas. Además, cuento con experiencia como voluntario en una cooperativa de cacao de mi región. Es precisamente esta experiencia la que me enseñó el valor del trabajo riguroso.}

\esp{Quedo a su disposición para una entrevista cuando lo considere oportuno.}

\esp{Le saluda atentamente,}

\esp{Andrés Mbarga}

Justifications : vouvoiement constant (\esp{se dirijo a usted}, \esp{su disposición}) ; deux-points après l'appel ; emphase par clivée (\esp{Es precisamente esta experiencia la que...}) ; connecteurs \esp{Además}, \esp{lo que} ; vocabulaire de l'emploi (\esp{puesto, solicitar, experiencia, entrevista}). Comptage : environ 95 mots hors en-tête et signature ; on peut ajouter une phrase de qualités (\esp{Soy puntual, organizado y me adapto con facilidad.}) pour atteindre 120 mots. Conclusion : une lettre efficace est courte, structurée, formelle et contient au moins une tournure emphatique.
\end{exoresolu}

\courssub{Fiche méthode 5 : Commenter l'emphase dans un texte}

\begin{methode}[Analyser un procédé d'emphase en quatre temps]
\begin{enumerate}
\item \cle{Repérer} le procédé : clivée, ser de relieve, inversion, répétition pronominale, diminutif.
\item \cle{Citer} le passage exact en espagnol, entre guillemets.
\item \cle{Décrire} le mécanisme : nommer la structure et la reformuler de façon neutre (\esp{Fue la carta la que lo cambió todo} $\rightarrow$ forme neutre : \esp{La carta lo cambió todo}).
\item \cle{Interpréter} : dire quel effet produit l'emphase (insistance, émotion, mise en valeur du sujet, ton affectueux).
\end{enumerate}
Phrase-type d'analyse : « L'auteur met en relief \esp{X} par une phrase clivée, ce qui souligne l'importance décisive de \esp{X} dans son parcours. »
\end{methode}

\begin{exoresolu}[Commentaire guidé]
Énoncé. Relevez et analysez deux procédés d'emphase dans cet extrait : \esp{Fue aquella carta la que cambió mi destino. Yo, que nunca había salido de Yaoundé, recibí de pronto una respuesta de Madrid. ¡Qué emocionante resultó aquel momentito!}
\tcblower
Solution détaillée.
\begin{enumerate}
\item \esp{Fue aquella carta la que cambió mi destino} : phrase clivée (\esp{Fue + sujet + la que + verbe}). Forme neutre : \esp{Aquella carta cambió mi destino}. Effet : la lettre devient l'unique héroïne du récit ; le narrateur souligne qu'aucun autre facteur n'a compté.
\item \esp{aquel momentito} : diminutif en \esp{-ito} appliqué à \esp{momento}. Effet : affection et nostalgie ; le narrateur minimise la durée pour mieux en exalter l'intensité émotionnelle. On peut aussi relever l'exclamation \esp{¡Qué emocionante...!}, autre procédé d'emphase expressive.
\end{enumerate}
Conclusion : un bon commentaire nomme le procédé, le cite, le reformule et en interprète l'effet. Quatre étapes, quatre points gagnés.
\end{exoresolu}

\begin{attention}[Les 5 erreurs qui coûtent des points au Bac]
\begin{enumerate}
\item \cle{Faute d'accent} : \esp{carta} (lettre) sans accent, mais \esp{está} (il est) avec accent ; \esp{si} (si) $\neq$ \esp{sí} (oui) ; \esp{tu} (ton) $\neq$ \esp{tú} (toi). Une lettre de motivation avec \esp{tu} au lieu de \esp{usted} est éliminatoire sur le plan de la correction sociale.
\item \cle{Faux amis} : « une candidature » = \esp{una candidatura}, pas \esp{candidación} ; « un emploi » = \esp{un empleo/un puesto}, pas \esp{una implicación} ; « assister à » = \esp{asistir a} (et non « aider », qui se dit \esp{ayudar}) ; \esp{actualmente} signifie « actuellement », jamais « actuel » au sens de « réel ».
\item \cle{Calque du français} : « Je suis intéressé par ce poste » ne se traduit pas par \esp{Estoy interesado} partout : préfère la forme soutenue \esp{Me interesa el puesto}. « Je vous écris pour... » $\rightarrow$ \esp{Me dirijo a usted para...}, pas un calque de « écrire ».
\item \cle{Accord} : après \esp{Fui yo quien...}, le verbe s'accorde avec la personne (\esp{Fui yo quien escribió}, première personne) ; \esp{lo que} reste invariable ; l'adjectif s'accorde : \esp{una secretaria bilingüe}.
\item \cle{Choix du temps} : dans une lettre, l'expérience passée se raconte au passé simple ou au parfait (\esp{trabajé / he trabajado}), les qualités au présent (\esp{soy puntual}), et la demande au présent de politesse (\esp{Quedo a su disposición}). Ne mélange pas futur et conditionnel : \esp{Le agradecería que me concediera una entrevista} (conditionnel + imparfait du subjonctif), jamais \esp{agradecería que me concederá}.
\end{enumerate}
\end{attention}

\newpage

\coursec{Exercices}

\courssub{Je vérifie mes connaissances}

\begin{exercice}[QCM : l'emphase]
Entoure la bonne réponse.
\begin{enumerate}
\item Dans la phrase \esp{Fue Ana quien me ayudó}, le procédé d'emphase utilisé est :
\begin{enumerate}
\item le \esp{ser de relieve} \item la répétition du pronom \item le diminutif
\end{enumerate}
\item \esp{A mí me gusta trabajar en equipo} est une phrase :
\begin{enumerate}
\item passive \item avec redoublement pronominal \item impersonnelle
\end{enumerate}
\item \esp{Lo que busco es un contrato estable} est une phrase :
\begin{enumerate}
\item clivée avec \esp{lo que} \item au subjonctif \item interrogative indirecte
\end{enumerate}
\item Le diminutif de \esp{carta} est :
\begin{enumerate}
\item \esp{cartaza} \item \esp{cartita} \item \esp{cartona}
\end{enumerate}
\end{enumerate}
\end{exercice}

\begin{exercice}[Vrai ou faux ?]
Indique si chaque affirmation est vraie ou fausse, puis justifie ta réponse par un exemple.
\begin{enumerate}
\item En espagnol, on peut placer le complément direct avant le verbe en le reprenant par un pronom : \esp{El currículum lo envié ayer}.
\item La formule \esp{Muy señor mío} convient pour ouvrir une lettre de motivation formelle.
\item Le mot \esp{habilidades} est un faux ami : il signifie « habitudes ».
\item Une lettre de motivation se termine par une formule de politesse comme \esp{Atentamente}.
\end{enumerate}
\end{exercice}

\begin{exercice}[Vocabulaire : le monde de l'emploi]
Complète chaque phrase avec le mot espagnol qui convient : \esp{entrevista, contrato, currículum, cualificación, jornada}.
\begin{enumerate}
\item Antes de la \_\_\_\_\_\_\_\_\_\_\_\_ de trabajo, repaso mis respuestas.
\item He enviado mi \_\_\_\_\_\_\_\_\_\_\_\_ vitae a la empresa de Valencia.
\item Me han ofrecido un \_\_\_\_\_\_\_\_\_\_\_\_ indefinido.
\item Trabajo a \_\_\_\_\_\_\_\_\_\_\_\_ completa, ocho horas al día.
\item Mi \_\_\_\_\_\_\_\_\_\_\_\_ profesional es el bachillerato.
\end{enumerate}
\end{exercice}

\begin{exercice}[Conjugaison à compléter]
Complète les phrases en conjuguant le verbe entre parenthèses au présent de l'indicatif.
\begin{enumerate}
\item \esp{Yo \_\_\_\_\_\_\_\_\_\_\_\_ (querer) este puesto de trabajo.}
\item \esp{Nosotros \_\_\_\_\_\_\_\_\_\_\_\_ (tener) mucha experiencia.}
\item \esp{Usted \_\_\_\_\_\_\_\_\_\_\_\_ (poder) contar conmigo.}
\item \esp{Yo \_\_\_\_\_\_\_\_\_\_\_\_ (ser) una persona seria y puntual.}
\item \esp{Ellos \_\_\_\_\_\_\_\_\_\_\_\_ (buscar) un asistente comercial.}
\end{enumerate}
\end{exercice}

\courssub{Je m'entraîne}

\begin{exercice}[Traduction : thème]
Traduis en espagnol les phrases suivantes en utilisant une tournure emphatique.
\begin{enumerate}
\item C'est mon expérience qui fait la différence.
\item Moi, j'aime le travail en équipe.
\item Ce que je cherche, c'est un poste stable.
\item C'est hier que j'ai envoyé ma lettre.
\end{enumerate}
\end{exercice}

\begin{exercice}[Transformation de phrases]
Transforme les phrases suivantes en utilisant la structure indiquée entre parenthèses.
\begin{enumerate}
\item \esp{Mi experiencia es amplia.} (Phrase clivée avec \esp{lo que})
\item \esp{Quiero un contrato estable.} (\esp{Ser de relieve} : \esp{Es... lo que...})
\item \esp{Yo envié el CV ayer.} (Reprise pronominale + ordre complément-verbe-sujet)
\item \esp{Me interesa este puesto.} (Redoublement : \esp{A mí...})
\end{enumerate}
\end{exercice}

\begin{exercice}[Texte à trous]
Complète le paragraphe avec les mots de la liste : \esp{contrato indefinido, jornada, habilidades blandas, currículum, entrevista, cualificación}.

\esp{Después de enviar mi \_\_\_\_\_\_\_\_\_\_\_\_ vitae, la empresa me llamó para una \_\_\_\_\_\_\_\_\_\_\_\_. Allí hablé de mi \_\_\_\_\_\_\_\_\_\_\_\_ profesional y de mis \_\_\_\_\_\_\_\_\_\_\_\_ \_\_\_\_\_\_\_\_\_\_\_\_, como la comunicación y el trabajo en equipo. Al final, me ofrecieron un \_\_\_\_\_\_\_\_\_\_\_\_ \_\_\_\_\_\_\_\_\_\_\_\_ a \_\_\_\_\_\_\_\_\_\_\_\_ completa. ¡Estoy muy contento!}
\end{exercice}

\begin{exercice}[Appariements : formules de la lettre]
Relie chaque élément de la lettre de motivation (1-5) à la formule espagnole correspondante (a-e).
\begin{enumerate}
\item Formule d'appel
\item Présentation
\item Expression de la motivation
\item Disponibilité pour un entretien
\item Formule de politesse finale
\end{enumerate}
\begin{enumerate}
\item[\textsf{a.}] \esp{Quedo a su disposición para una entrevista.}
\item[\textsf{b.}] \esp{Estimado señor director:}
\item[\textsf{c.}] \esp{Atentamente,}
\item[\textsf{d.}] \esp{Me llamo... y soy titulado en...}
\item[\textsf{e.}] \esp{Me dirijo a ustedes porque me interesa el puesto de...}
\end{enumerate}
\end{exercice}

\courssub{Je me prépare au Bac}

\begin{exercice}[Type Bac : compréhension et langue]
Lis le texte suivant, puis réponds aux questions.

\begin{textebox}[Un primer empleo en Malabo]
Cuando terminé mis estudios en Bata, pensé que encontrar trabajo sería fácil. Pero no fue así: fue la perseverancia la que me abrió las puertas. A mí me habían dicho que sin contactos no había oportunidades, y sin embargo, lo que de verdad importó fue mi esfuerzo.

Envié mi currículum a muchas empresas. Un día, una compañía de Malabo me llamó para una entrevista. Fue aquel día el que cambió mi vida. El director me preguntó por mis habilidades y yo le respondí con sinceridad: lo que yo ofrezco es seriedad y ganas de aprender. Me contrataron con un contrato indefinido y a jornada completa.

Hoy, a los jóvenes que buscan empleo, les digo siempre lo mismo: preparad bien vuestro currículum, cuidad cada cartita que enviéis y no perdáis nunca la esperanza. Es el esfuerzo el que siempre triunfa.
\end{textebox}

\textbf{Glossaire} : \esp{el esfuerzo} = l'effort ; \esp{las ganas de aprender} = l'envie d'apprendre ; \esp{contratado} = embauché ; \esp{la esperanza} = l'espérance ; \esp{enviéis} = que vous envoyiez (subjonctif présent, 2\up{e} pers. du pluriel).

\textbf{Comprensión lectora} (réponds en espagnol, avec des phrases complètes) :
\begin{enumerate}
\item ¿Dónde terminó el narrador sus estudios?
\item ¿Qué le abrió las puertas, según el texto?
\item ¿Qué tipo de contrato le ofrecieron?
\item ¿Qué consejos da el narrador a los jóvenes?
\end{enumerate}

\textbf{Lengua} :
\begin{enumerate}
\item Releve dans le texte trois tournures emphatiques différentes et nomme le procédé utilisé (clivée, \esp{ser de relieve}, redoublement pronominal, diminutif).
\item Transforme : \esp{El esfuerzo triunfa siempre.} $\rightarrow$ phrase avec \esp{ser de relieve}.
\item Traduis en français : \esp{A mí me habían dicho que sin contactos no había oportunidades.}
\end{enumerate}
\end{exercice}

\begin{exercice}[Type Bac : thème et version]
\textbf{Version.} Traduis en français le passage suivant :

\esp{Fue su oferta la que captó mi atención. A mí me gustan los retos. Lo que ofrezco es seriedad. Es por eso por lo que les escribo esta carta.}

\textbf{Thème.} Traduis en espagnol :
\begin{enumerate}
\item C'est votre entreprise qui m'intéresse le plus.
\item Moi, je parle trois langues : le français, l'anglais et l'espagnol.
\item Ce que je veux, c'est travailler avec vous.
\item C'est à Yaoundé que j'ai obtenu mon baccalauréat.
\end{enumerate}
\end{exercice}

\begin{exercice}[Type Bac : rédaction]
\textbf{Sujet.} Vous êtes titulaire du Baccalauréat A. Vous répondez à une annonce pour un poste d'assistant commercial export dans une entreprise agroalimentaire à Valence (Espagne). Rédigez votre lettre de motivation en espagnol.

\textbf{Consignes :}
\begin{itemize}
\item Longueur : environ \textbf{120 mots}.
\item Respectez la structure de la lettre formelle : lieu et date, formule d'appel, présentation, motivations, disponibilité pour un entretien, formule de politesse.
\item Intégrez \textbf{au moins deux tournures emphatiques distinctes} (phrase clivée, \esp{ser de relieve}, redoublement pronominal).
\item Employez le lexique de l'emploi étudié dans la leçon (\esp{puesto, currículum, experiencia, habilidades, entrevista}...).
\end{itemize}
\end{exercice}



\newpage

\coursec{Corrigés des exercices}

\courssub{Je vérifie mes connaissances}

\begin{corrige}[QCM : l'emphase]
\begin{enumerate}
\item Réponse a : \esp{Fue Ana quien me ayudó} est la tournure du \cle{ser de relieve} (\esp{ser} + élément mis en relief + \esp{quien/que}).
\item Réponse b : \esp{A mí me gusta...} présente un \cle{redoublement pronominal} (\esp{a mí} reprend \esp{me}).
\item Réponse a : \esp{Lo que busco es...} est une \cle{phrase clivée} avec \esp{lo que}.
\item Réponse b : le diminutif de \esp{carta} est \esp{cartita} (\esp{cartaza} serait un augmentatif péjoratif, \esp{cartón} signifie « carton »).
\end{enumerate}
\end{corrige}

\begin{corrige}[Vrai ou faux ?]
\begin{enumerate}
\item \textbf{Vrai.} L'antéposition du complément d'objet direct exige la reprise pronominale : \esp{El currículum lo envié ayer} (« Le CV, je l'ai envoyé hier »).
\item \textbf{Vrai.} \esp{Muy señor mío:} est une formule d'appel très formelle, adaptée à une lettre de motivation (avec deux-points et majuscule à la ligne suivante).
\item \textbf{Faux.} \esp{Habilidades} signifie « compétences, aptitudes » ; « habitudes » se dit \esp{costumbres} ou \esp{hábitos}. Exemple : \esp{Tengo buenas habilidades comunicativas} (« J'ai de bonnes compétences en communication »).
\item \textbf{Vrai.} \esp{Atentamente,} (« Cordialement ») est une formule de politesse finale classique, comme \esp{Le saluda atentamente,}.
\end{enumerate}
\end{corrige}

\begin{corrige}[Vocabulaire : le monde de l'emploi]
\begin{enumerate}
\item \esp{entrevista} : \esp{Antes de la entrevista de trabajo, repaso mis respuestas.} (« Avant l'entretien d'embauche, je révise mes réponses. »)
\item \esp{currículum} : \esp{He enviado mi currículum vitae a la empresa de Valencia.}
\item \esp{contrato} : \esp{Me han ofrecido un contrato indefinido.} (« On m'a proposé un CDI. »)
\item \esp{jornada} : \esp{Trabajo a jornada completa, ocho horas al día.} (« Je travaille à temps plein. »)
\item \esp{cualificación} : \esp{Mi cualificación profesional es el bachillerato.}
\end{enumerate}
\end{corrige}

\begin{corrige}[Conjugaison à compléter]
\begin{enumerate}
\item \esp{quiero} (verbe à diphtongaison e $\rightarrow$ ie) : \esp{Yo quiero este puesto de trabajo.}
\item \esp{tenemos} (diphtongaison seulement aux trois personnes du singulier et à la 3\up{e} du pluriel) : \esp{Nosotros tenemos mucha experiencia.}
\item \esp{puede} (o $\rightarrow$ ue) : \esp{Usted puede contar conmigo.}
\item \esp{soy} (irrégulier) : \esp{Yo soy una persona seria y puntual.}
\item \esp{buscan} (régulier) : \esp{Ellos buscan un asistente comercial.}
\end{enumerate}
\end{corrige}

\courssub{Je m'entraîne}

\begin{corrige}[Traduction : thème]
\begin{enumerate}
\item \esp{Es mi experiencia la que marca la diferencia.} (\esp{ser de relieve})
\item \esp{A mí me gusta el trabajo en equipo.} (redoublement pronominal)
\item \esp{Lo que busco es un puesto estable.} (phrase clivée avec \esp{lo que})
\item \esp{Fue ayer cuando envié mi carta.} (clivée avec \esp{cuando} pour un complément de temps)
\end{enumerate}
\end{corrige}

\begin{corrige}[Transformation de phrases]
\begin{enumerate}
\item \esp{Lo que es amplia es mi experiencia.}
\item \esp{Es un contrato estable lo que quiero.}
\item \esp{El CV lo envié yo ayer.}
\item \esp{A mí me interesa este puesto.}
\end{enumerate}
\end{corrige}

\begin{corrige}[Texte à trous]
\esp{Después de enviar mi \textbf{currículum} vitae, la empresa me llamó para una \textbf{entrevista}. Allí hablé de mi \textbf{cualificación} profesional y de mis \textbf{habilidades blandas}, como la comunicación y el trabajo en equipo. Al final, me ofrecieron un \textbf{contrato indefinido} a \textbf{jornada} completa. ¡Estoy muy contento!}

Traduction : « Après avoir envoyé mon CV, l'entreprise m'a appelé pour un entretien. Là, j'ai parlé de ma qualification professionnelle et de mes compétences relationnelles, comme la communication et le travail en équipe. Finalement, on m'a proposé un CDI à temps plein. Je suis très content ! »
\end{corrige}

\begin{corrige}[Appariements : formules de la lettre]
1 -- b ; 2 -- d ; 3 -- e ; 4 -- a ; 5 -- c.

Rappel de l'ordre canonique de la lettre : appel (\esp{Estimado señor director:}) $\rightarrow$ présentation (\esp{Me llamo...}) $\rightarrow$ motivation (\esp{Me dirijo a ustedes...}) $\rightarrow$ disponibilité (\esp{Quedo a su disposición...}) $\rightarrow$ politesse finale (\esp{Atentamente,}).
\end{corrige}

\courssub{Je me prépare au Bac}

\begin{corrige}[Type Bac : compréhension et langue]
\textbf{Comprensión lectora} (réponses modèles) :
\begin{enumerate}
\item \esp{El narrador terminó sus estudios en Bata.}
\item \esp{Según el texto, fue la perseverancia la que le abrió las puertas.}
\item \esp{Le ofrecieron un contrato indefinido a jornada completa.}
\item \esp{Les aconseja preparar bien el currículum, cuidar cada carta que envíen y no perder nunca la esperanza.}
\end{enumerate}

\textbf{Lengua} :
\begin{enumerate}
\item Exemples de réponses : \esp{fue la perseverancia la que me abrió las puertas} (\esp{ser de relieve}) ; \esp{lo que de verdad importó fue mi esfuerzo} (phrase clivée avec \esp{lo que}) ; \esp{A mí me habían dicho} (redoublement pronominal) ; \esp{cada cartita} (diminutif affectif).
\item \esp{Es el esfuerzo lo que siempre triunfa.} (ou \esp{Es el esfuerzo el que siempre triunfa.})
\item « On m'avait dit, à moi, que sans relations il n'y avait pas d'opportunités. »
\end{enumerate}
\end{corrige}

\begin{corrige}[Type Bac : thème et version]
\textbf{Version.} « C'est votre offre qui a retenu mon attention. Moi, j'aime les défis. Ce que j'offre, c'est le sérieux. C'est pour cela que je vous écris cette lettre. »

\textbf{Thème.}
\begin{enumerate}
\item \esp{Es su empresa la que más me interesa.} (\esp{ser de relieve})
\item \esp{A mí me gustan los idiomas: hablo tres, el francés, el inglés y el español.} ou plus simplement \esp{Yo hablo tres idiomas: francés, inglés y español.} avec redoublement possible : \esp{A mí me apasionan los idiomas...}
\item \esp{Lo que quiero es trabajar con ustedes.} (phrase clivée)
\item \esp{Fue en Yaoundé donde obtuve mi bachillerato.} (clivée avec \esp{donde} pour un complément de lieu)
\end{enumerate}
\end{corrige}

\begin{corrige}[Type Bac : rédaction]
\textbf{Modèle de rédaction} (environ 120 mots) :

\esp{Yaoundé, 15 de mayo de 2025}

\esp{Estimado señor director:}

\esp{Me llamo André Mbarga y soy titulado del Bachillerato A. Me dirijo a ustedes porque me interesa el puesto de asistente comercial export anunciado por su empresa.}

\esp{Fue su oferta la que captó mi atención: es el comercio internacional lo que me apasiona. A mí me gustan los idiomas y hablo francés, inglés y español con fluidez. Además, lo que yo ofrezco es seriedad, puntualidad y verdaderas ganas de aprender.}

\esp{Adjunto mi currículum vitae y quedo a su disposición para una entrevista cuando lo consideren oportuno.}

\esp{Atentamente,}

\esp{André Mbarga}

\textbf{Grille d'auto-évaluation} : lieu et date en haut (2 points) ; formule d'appel avec deux-points (2 points) ; présentation et motivation (4 points) ; deux tournures emphatiques repérables (\esp{Fue su oferta la que...}, \esp{es el comercio... lo que...}, \esp{A mí me gustan...}, \esp{lo que yo ofrezco es...}) (4 points) ; lexique de l'emploi (4 points) ; formule finale et signature (2 points) ; correction linguistique (2 points).
\end{corrige}

\newpage

\coursec{Fiche bilan}

\begin{popfigure}
\begin{tikzpicture}[mindmap, grow cyclic, every node/.style={concept, font=\small},
root concept/.append style={concept color=popblue, font=\bfseries},
level 1/.append style={level distance=3.4cm, sibling angle=51.4},
level 2/.append style={level distance=1.9cm}]
\node[root concept, align=center]{E18\\ Emphase \&\\ Carta de\\ motivación}
child[concept color=poporange]{node[align=center]{Frase escindida\\ \esp{Lo que... es...}}}
child[concept color=popgreen]{node[align=center]{Ser de relieve\\ \esp{Fue Ana quien...}}}
child[concept color=popgold]{node[align=center]{Reduplicación\\ del pronombre\\ \esp{A mí me...}}}
child[concept color=poppurple]{node[align=center]{Orden de\\ las palabras}}
child[concept color=poppink]{node[align=center]{Diminutivos y\\ aumentativos}}
child[concept color=popblue!70]{node[align=center]{Carta de\\ motivación\\ (120 palabras)}}
child[concept color=popgreen!70]{node[align=center]{Léxico\\ del empleo}};
\end{tikzpicture}
\legende{Carte mentale de la leçon E18 : les mécanismes de l'emphase et la lettre de motivation.}
\end{popfigure}

\begin{bilan}
\textbf{1. Lexique clé du monde de l'emploi}
\begin{itemize}
\item \esp{el puesto de trabajo} : le poste ; \esp{la oferta de empleo} : l'offre d'emploi ; \esp{la entrevista} : l'entretien.
\item \esp{el currículum vítae (CV)} : le CV ; \esp{la carta de motivación} : la lettre de motivation.
\item \esp{el/la candidato/a} : le/la candidat(e) ; \esp{el/la empleador(a)} : l'employeur ; \esp{la empresa} : l'entreprise.
\item \esp{solicitar un empleo} : postuler à un emploi ; \esp{contratar} : embaucher ; \esp{el sueldo / el salario} : le salaire.
\item \esp{la experiencia laboral} : l'expérience professionnelle ; \esp{la formación} : la formation ; \esp{las cualidades} : les qualités.
\item \esp{quedo a su disposición} : je reste à votre disposition ; \esp{adjunto mi CV} : je joins mon CV.
\end{itemize}

\textbf{2. Grammaire à connaître par cœur : les mécanismes de l'emphase}
\begin{center}
\begin{tabularx}{\linewidth}{|l|X|X|}
\hline
\rowcolor{popblueL} \textbf{Procédé} & \textbf{Forme} & \textbf{Exemple} \\
\hline
Phrase clivée & \esp{Lo que} + verbe + \esp{es} + élément mis en relief & \esp{Lo que busco es un empleo estable.} (Ce que je cherche, c'est un emploi stable.) \\
\hline
\emph{Ser} de relief & \esp{Ser} + élément + \esp{quien/que} + verbe & \esp{Fue mi madre quien me animó.} (C'est ma mère qui m'a encouragé.) \\
\hline
Reduplication du pronom & COI/COD pronominal + complément repris & \esp{A mí me gusta trabajar en equipo.} (Moi, j'aime travailler en équipe.) \\
\hline
Antéposition & Complément placé en tête de phrase & \esp{Ese puesto lo quiero yo.} (Ce poste, c'est moi qui le veux.) \\
\hline
Diminutifs / augmentatifs & suffixes \esp{-ito/-cito}, \esp{-ón/-azo} & \esp{un momentito} (un petit moment) ; \esp{un problemón} (un gros problème) \\
\hline
\end{tabularx}
\end{center}

\textbf{3. Méthode express : la lettre de motivation en 120 mots}
\begin{enumerate}
\item En-tête : lieu et date (\esp{Yaundé, a 15 de mayo de 2024}) + coordonnées du destinataire.
\item Formule d'appel : \esp{Estimado señor Director:} / \esp{Muy señor mío:} (jamais de \esp{querido} dans un courrier formel).
\item Paragraphe 1 : l'objet — \esp{Me dirijo a usted para solicitar el puesto de...}
\item Paragraphe 2 : le profil — formation, expérience, qualités (\esp{poseo una sólida formación en...}).
\item Paragraphe 3 : la motivation — pourquoi cette entreprise, avec une emphase (\esp{Lo que más me atrae de su empresa es...}).
\item Formule de politesse : \esp{Le saluda atentamente,} / \esp{Atentamente,} + signature.
\end{enumerate}
\end{bilan}

\begin{aretenir}[Checklist avant le Bac]
\begin{itemize}
\item[$\square$] Je sais transformer une phrase neutre en phrase clivée avec \esp{Lo que... es...}.
\item[$\square$] Je maîtrise le \emph{ser} de relief : \esp{Fue + nom + quien/que + verbe}.
\item[$\square$] Je sais utiliser la reduplication pronominale : \esp{A mí me...}, \esp{A nosotros nos...}.
\item[$\square$] Je connais les suffixes \esp{-ito/-cito/-ecito} et \esp{-ón/-azo} et leur valeur affective.
\item[$\square$] Je respecte la structure en trois paragraphes de la lettre de motivation.
\item[$\square$] J'utilise le registre formel : \esp{usted}, jamais \esp{tú}, dans une lettre professionnelle.
\item[$\square$] Je connais au moins trois formules d'appel et trois formules de politesse finales.
\item[$\square$] Je sais écrire la date à l'espagnole : \esp{a 15 de mayo de 2024} (mois sans majuscule).
\item[$\square$] Je compte mes mots : environ 120, ni trop peu ni trop.
\item[$\square$] Je place au moins une tournure d'emphase dans ma lettre pour valoriser ma copie.
\item[$\square$] Je vérifie les accents (\esp{motivación}, \esp{atención}) et la ponctuation d'ouverture ¡ ¿.
\item[$\square$] Je me relis : accord des adjectifs, \emph{ser/estar}, pas de calque du français (\esp{realizar} $\neq$ « réaliser » au sens de « se rendre compte »).
\end{itemize}
\end{aretenir}

\findecours

\end{document}
